\documentclass[11pt,a4paper]{article}
\usepackage[margin=1in]{geometry}
\usepackage{amsmath,amssymb}
\usepackage{enumitem}
\usepackage{multicol}
\usepackage{titlesec}
\usepackage{fancyhdr}

\pagestyle{fancy}
\fancyhf{}
\fancyhead[C]{The Hope Science Academy -- Mathematics 11th (Combined)}
\fancyfoot[C]{\thepage}
\setlength{\headheight}{14pt}

\titleformat{\section}{\large\bfseries}{}{0em}{}

\newlist{mcqlist}{enumerate}{1}
\setlist[mcqlist,1]{label=\textbf{\arabic*.}, leftmargin=2.2em, itemsep=8pt, topsep=4pt}
\newlist{shortlist}{enumerate}{1}
\setlist[shortlist,1]{label=\textbf{\arabic*.}, leftmargin=2.5em, itemsep=8pt, topsep=4pt}
\newlist{longlist}{enumerate}{1}
\setlist[longlist,1]{label=\textbf{\arabic*.}, leftmargin=*, itemsep=10pt, topsep=6pt}

\title{\textbf{The Hope Science Academy}\\[2pt]\large Mathematics -- Class 11th\\[2pt]\normalsize Combined Question Bank (Full Book)}
\author{}
\date{}

\begin{document}
\maketitle
\tableofcontents
\newpage

\section{Multiple Choice Questions (1--585)}

\section*{Choose the correct answer.}

\begin{mcqlist}

\item The number $\sqrt{-1}$ is called:
\begin{quote}\begin{multicols}{2}
\noindent (A) Real number \\
(B) Natural number \\
(C) Complex number \\
(D) Rational number
\end{multicols}\end{quote}

\item $(-i)^{19}$ is equal to:
\begin{quote}\begin{multicols}{2}
\noindent (A) $-i$ \\
(B) $i$ \\
(C) $1$ \\
(D) $-1$
\end{multicols}\end{quote}

\item $(-1)^{-21/2}$ equals:
\begin{quote}\begin{multicols}{2}
\noindent (A) $i$ \\
(B) $-i$ \\
(C) $1$ \\
(D) $-1$
\end{multicols}\end{quote}

\item The multiplicative inverse of complex number $(0,-1)$ is equal to:
\begin{quote}\begin{multicols}{2}
\noindent (A) $(1,0)$ \\
(B) $(0,1)$ \\
(C) $(-1,0)$ \\
(D) $(0,0)$
\end{multicols}\end{quote}

\item Multiplicative inverse of $(1,0)$ is:
\begin{quote}\begin{multicols}{2}
\noindent (A) $(-1,0)$ \\
(B) $(0,1)$ \\
(C) $(0,-1)$ \\
(D) $(1,0)$
\end{multicols}\end{quote}

\item Multiplicative inverse of $-i$ is:
\begin{quote}\begin{multicols}{2}
\noindent (A) $i$ \\
(B) $-i$ \\
(C) $1$ \\
(D) $-1$
\end{multicols}\end{quote}

\item Real part of $\dfrac{3}{\sqrt{6}-\sqrt{-12}}$ is:
\begin{quote}\begin{multicols}{2}
\noindent (A) $\dfrac{1}{\sqrt6}+\dfrac{1}{\sqrt3}$ \\
(B) $\dfrac{i}{\sqrt6}+\dfrac{1}{\sqrt3}$ \\
(C) $\dfrac{1}{\sqrt6}+\dfrac{i}{\sqrt3}$ \\
(D) $\dfrac{i}{\sqrt6}+\dfrac{i}{\sqrt3}$
\end{multicols}\end{quote}

\item Any real number ``$a$'' is equal to:
\begin{quote}\begin{multicols}{2}
\noindent (A) $ia$ \\
(B) $a,b$ \\
(C) $a$ \\
(D) $b,a$
\end{multicols}\end{quote}

\item The real part of $\dfrac{1+3i}{2i}$ is:
\begin{quote}\begin{multicols}{2}
\noindent (A) $\dfrac{3}{2}$ \\
(B) $\dfrac{2}{3}$ \\
(C) $1$ \\
(D) $2$
\end{multicols}\end{quote}

\item If $Z = 3-4i$ then $|\overline{Z}|$ is:
\begin{quote}\begin{multicols}{2}
\noindent (A) $4$ \\
(B) $5$ \\
(C) $-5$ \\
(D) $1$
\end{multicols}\end{quote}

\item Modulus of complex number $-5i$ is:
\begin{quote}\begin{multicols}{2}
\noindent (A) $\pm5$ \\
(B) $-5$ \\
(C) $\sqrt5$ \\
(D) $5$
\end{multicols}\end{quote}

\item If $Z$ is a complex number, then $|Z|^2$ is:
\begin{quote}\begin{multicols}{2}
\noindent (A) $Z^2$ \\
(B) $|\overline Z|^2$ \\
(C) $Z\overline Z$ \\
(D) $\dfrac{Z}{\overline Z}$
\end{multicols}\end{quote}

\item If $Z=-2+3i$ then $\overline Z = ?$:
\begin{quote}\begin{multicols}{2}
\noindent (A) $-2-3i$ \\
(B) $2-3i$ \\
(C) $-2+3i$ \\
(D) $2+3i$
\end{multicols}\end{quote}

\item If $Z=-2+3i$ then $\overline Z = ?$:
\begin{quote}\begin{multicols}{2}
\noindent (A) $-2-3i$ \\
(B) $2-3i$ \\
(C) $-2+3i$ \\
(D) $2+3i$
\end{multicols}\end{quote}

\item If $Z_1=1+i$, $Z_2=1-i$ then $\sqrt{z}$ is:
\begin{quote}\begin{multicols}{2}
\noindent (A) $Z_1<Z_2$ \\
(B) $Z_1>Z_2$ \\
(C) $\overline{Z_1}>\overline{Z_2}$ \\
(D) $|Z_1|=|Z_2|$
\end{multicols}\end{quote}

\item If $Z=-7-24i$ then real part of $\sqrt Z$ is:
\begin{quote}\begin{multicols}{2}
\noindent (A) $2$ \\
(B) $3$ \\
(C) $4$ \\
(D) $1$
\end{multicols}\end{quote}

\item The equation $x^2+1=0$ has solution in:
\begin{quote}\begin{multicols}{2}
\noindent (A) $\mathbb R$ \\
(B) $\mathbb C$ \\
(C) $\mathbb Q$ \\
(D) $\mathbb P$
\end{multicols}\end{quote}

\item $a^2+b^2$ has factors:
\begin{quote}\begin{multicols}{2}
\noindent (A) $(a+b)(a-b)$ \\
(B) $(a+b)(a-ib)$ \\
(C) $(a+ib)(a-ib)$ \\
(D) $(a+ib)(a-b)$
\end{multicols}\end{quote}

\item Complex cube roots of $-1$ are:
\begin{quote}\begin{multicols}{2}
\noindent (A) $\omega,\omega^2$ \\
(B) $1,\omega,\omega^2$ \\
(C) $-1,-\omega,-\omega^2$ \\
(D) $-\omega,-\omega^2$
\end{multicols}\end{quote}

\item Three cube roots of $27$ are:
\begin{quote}\begin{multicols}{2}
\noindent (A) $3,\omega,\omega^2$ \\
(B) $-3,-3\omega,3\omega^2$ \\
(C) $3,3\omega,3\omega^2$ \\
(D) $1,\omega,\omega^2$
\end{multicols}\end{quote}

\item Cube roots of unity are:
\begin{quote}\begin{multicols}{2}
\noindent (A) $-1,-2,1$ \\
(B) $1,-1,\omega$ \\
(C) $1,\omega,\omega^2$ \\
(D) $-1,-\omega,\omega^2$
\end{multicols}\end{quote}

\item Sum of all three cube roots of unity is:
\begin{quote}\begin{multicols}{2}
\noindent (A) $1$ \\
(B) $-1$ \\
(C) $3$ \\
(D) $0$
\end{multicols}\end{quote}

\item Sum of complex roots of unity is:
\begin{quote}\begin{multicols}{2}
\noindent (A) $0$ \\
(B) $-1$ \\
(C) $1$ \\
(D) $2$
\end{multicols}\end{quote}

\item Product of all cube roots of unity is:
\begin{quote}\begin{multicols}{2}
\noindent (A) $1$ \\
(B) $0$ \\
(C) $-1$ \\
(D) $2$
\end{multicols}\end{quote}

\item If $\omega$ is the imaginary/complex cube root of unity, then $\omega^2 = $:
\begin{quote}\begin{multicols}{2}
\noindent (A) $\omega$ \\
(B) $\omega^{-1}$ \\
(C) $\dfrac{1}{\omega^2}$ \\
(D) $\omega^3$
\end{multicols}\end{quote}

\item If $\omega$ is a complex cube root of unity then $\omega^{12}=$:
\begin{quote}\begin{multicols}{2}
\noindent (A) $\omega$ \\
(B) $-\omega$ \\
(C) $1$ \\
(D) $1$
\end{multicols}\end{quote}

\item If $\omega$ is a complex cube root of unity then $\omega^{29}+\omega^{28}+1=$:
\begin{quote}\begin{multicols}{2}
\noindent (A) $0$ \\
(B) $1$ \\
(C) $2$ \\
(D) $3$
\end{multicols}\end{quote}

\item If $\omega$ is the cube root of unity then $(1+\omega-\omega^2)^8=$:
\begin{quote}\begin{multicols}{2}
\noindent (A) $256$ \\
(B) $-256$ \\
(C) $-256\omega$ \\
(D) $256\omega$
\end{multicols}\end{quote}

\item If $\omega$ is the cube root of unity then $\omega+\omega^2$ equals:
\begin{quote}\begin{multicols}{2}
\noindent (A) $0$ \\
(B) $-1$ \\
(C) $1$ \\
(D) $\dfrac{1}{\omega}$
\end{multicols}\end{quote}

\item If $\omega$ be the cube root of unity, then $\omega^2=$:
\begin{quote}\begin{multicols}{2}
\noindent (A) $\dfrac{-1-\sqrt3 i}{2}$ \\
(B) $\dfrac{1-\sqrt3 i}{2}$ \\
(C) $1$ \\
(D) $\dfrac{1+\sqrt3 i}{2}$
\end{multicols}\end{quote}

\item If $\omega$ is a complex cube root of unity, $\omega^3=$:
\begin{quote}\begin{multicols}{2}
\noindent (A) $\omega$ \\
(B) $\omega^2$ \\
(C) $-1$ \\
(D) $1$
\end{multicols}\end{quote}

\item If $\omega$ is a complex cube root of unity then $\omega=$:
\begin{quote}\begin{multicols}{2}
\noindent (A) $0$ \\
(B) $1$ \\
(C) $\omega^2$ \\
(D) $\dfrac{1}{\omega^2}$
\end{multicols}\end{quote}

\item If $\omega$ is a complex cube root of unity then conjugate of $\omega$ is:
\begin{quote}\begin{multicols}{2}
\noindent (A) $\omega^2$ \\
(B) $-\omega^2$ \\
(C) $-\omega$ \\
(D) $i$
\end{multicols}\end{quote}

\item $1+\omega^4=$:
\begin{quote}\begin{multicols}{2}
\noindent (A) $\omega$ \\
(B) $-\omega$ \\
(C) $\omega^2$ \\
(D) $-\omega^2$
\end{multicols}\end{quote}

\item $(-1+\sqrt{-3})^5+(1-\sqrt3)^5=$:
\begin{quote}\begin{multicols}{2}
\noindent (A) $16$ \\
(B) $48$ \\
(C) $-32$ \\
(D) $-1$
\end{multicols}\end{quote}

\item The sum of the fourth roots of unity is:
\begin{quote}\begin{multicols}{2}
\noindent (A) $0$ \\
(B) $1$ \\
(C) $-1$ \\
(D) $2$
\end{multicols}\end{quote}

\item The sum of fourth roots of $16$ is:
\begin{quote}\begin{multicols}{2}
\noindent (A) $0$ \\
(B) $2$ \\
(C) $4$ \\
(D) $16$
\end{multicols}\end{quote}

\item The product of all four fourth roots of unity is:
\begin{quote}\begin{multicols}{2}
\noindent (A) $0$ \\
(B) $1$ \\
(C) $-1$ \\
(D) $4$
\end{multicols}\end{quote}

\item Four fourth roots of $81$ are:
\begin{quote}\begin{multicols}{2}
\noindent (A) $\pm4,\pm4i$ \\
(B) $\pm2,\pm2i$ \\
(C) $\pm3,\pm3i$ \\
(D) $\pm9,\pm9i$
\end{multicols}\end{quote}

\item For any $n\in\mathbb Z$, $\omega^n$ is equivalent to:
\begin{quote}\begin{multicols}{2}
\noindent (A) $1$ or $\omega$ or $\omega^2$ \\
(B) $2,2\omega,2\omega^2$ \\
(C) $3,3\omega,3\omega^2$ \\
(D) $\pm i,2$
\end{multicols}\end{quote}

\item If $\omega$ is a complex cube root of unity then value of $(3+\omega)(3+\omega^2)=?$:
\begin{quote}\begin{multicols}{2}
\noindent (A) $6$ \\
(B) $7$ \\
(C) $9$ \\
(D) $13$
\end{multicols}\end{quote}

\item Four fourth roots of $16$ are:
\begin{quote}\begin{multicols}{2}
\noindent (A) $2,-2,2i,-2i$ \\
(B) $1,-1,i,-i$ \\
(C) $3,-3,3i,-3i$ \\
(D) $4,-4,4i,-4i$
\end{multicols}\end{quote}

\item For a complex number, $(0,1)i=$:
\begin{quote}\begin{multicols}{2}
\noindent (A) $(1,0)$ \\
(B) $(-1,0)$ \\
(C) $(0,1)$ \\
(D) $(0,-1)$
\end{multicols}\end{quote}

\item If $z=a+ib$, what is the value of $\cos\theta$:
\begin{quote}\begin{multicols}{2}
\noindent (A) $\dfrac{a}{|z|}$ \\
(B) $\dfrac{b}{|z|}$ \\
(C) $\dfrac{a}{b}$ \\
(D) $\dfrac{b}{a}$
\end{multicols}\end{quote}

\item A complex number $1+i$ can also be expressed as:
\begin{quote}\begin{multicols}{2}
\noindent (A) $2(\cos60^\circ+i\sin30^\circ)$ \\
(B) $\cos60^\circ+i\sin30^\circ$ \\
(C) $\sqrt2(\cos60^\circ+i\sin60^\circ)$ \\
(D) $\cos30^\circ+i\sin30^\circ$
\end{multicols}\end{quote}

\item If $Z=\cos\theta+i\sin\theta$ then $|Z|^2=$:
\begin{quote}\begin{multicols}{2}
\noindent (A) $0$ \\
(B) $1$ \\
(C) $2$ \\
(D) $3$
\end{multicols}\end{quote}

\item Polar form of complex number $r\cos\theta+ir\sin\theta$ then $\theta=$:
\begin{quote}\begin{multicols}{2}
\noindent (A) $\tan^{-1}\dfrac{y}{x}$ \\
(B) $\tan\dfrac{y}{x}$ \\
(C) $\tan^{-1}\dfrac{x}{y}$ \\
(D) $\tan\dfrac{x}{y}$
\end{multicols}\end{quote}

\item $\left(\cos\dfrac{\pi}{6}+i\sin\dfrac{\pi}{6}\right)$ is equal to:
\begin{quote}\begin{multicols}{2}
\noindent (A) $1$ \\
(B) $0$ \\
(C) $-1$ \\
(D) $i$
\end{multicols}\end{quote}

\item Argument $\theta$ of $(\sqrt3+i)$ is:
\begin{quote}\begin{multicols}{2}
\noindent (A) $60^\circ$ \\
(B) $30^\circ$ \\
(C) $45^\circ$ \\
(D) $90^\circ$
\end{multicols}\end{quote}

\item If $f:A\rightarrow B$ is a function such that $\mathrm{range}\, f \subset B$, then $f$ is called a/an ...... function:
\begin{quote}\begin{multicols}{2}
\noindent (A) Onto \\
(B) Injective \\
(C) Bijective \\
(D) Into
\end{multicols}\end{quote}

\item If a function $f:A\rightarrow B$ is such that range $f=B$ then $f$ is called:
\begin{quote}\begin{multicols}{2}
\noindent (A) Injective \\
(B) Surjective \\
(C) Into \\
(D) Periodic
\end{multicols}\end{quote}

\item A bijective function is:
\begin{quote}\begin{multicols}{2}
\noindent (A) Both one-one and onto \\
(B) One-one but not onto \\
(C) Onto, but not one-one \\
(D) Neither one-one nor onto
\end{multicols}\end{quote}

\item If $f(x)=\sqrt{2x-3}$ then $f(1)=?$:
\begin{quote}\begin{multicols}{2}
\noindent (A) $\pm1$ \\
(B) $\pm i$ \\
(C) $\pm2$ \\
(D) $0$
\end{multicols}\end{quote}

\item Point of intersection $f(x)=2x+5$ and $g(x)=-x+5$ is a point which satisfies:
\begin{quote}\begin{multicols}{2}
\noindent (A) $f(x)$ \\
(B) $g(x)$ \\
(C) Both $f(x)$ and $g(x)$ \\
(D) None
\end{multicols}\end{quote}

\item For a table of a graph, we require minimum values of each function:
\begin{quote}\begin{multicols}{2}
\noindent (A) $1$ \\
(B) $2$ \\
(C) $3$ \\
(D) $4$
\end{multicols}\end{quote}

\item Graph of a quadratic equation will be:
\begin{quote}\begin{multicols}{2}
\noindent (A) Straight line \\
(B) Curve \\
(C) Set of lines \\
(D) None
\end{multicols}\end{quote}

\item Vertex is a minimum point when:
\begin{quote}\begin{multicols}{2}
\noindent (A) $a=0$ \\
(B) $a>0$ \\
(C) $a<0$ \\
(D) None
\end{multicols}\end{quote}

\item Inverse of a quadratic function should be:
\begin{quote}\begin{multicols}{2}
\noindent (A) Into \\
(B) Onto \\
(C) One to one \\
(D) Not to one
\end{multicols}\end{quote}

\item To solve absolute value quadratic function we must eliminate roots which are:
\begin{quote}\begin{multicols}{2}
\noindent (A) Real \\
(B) Imaginary \\
(C) Extraneous \\
(D) Double
\end{multicols}\end{quote}

\item A quadratic equation has degree:
\begin{quote}\begin{multicols}{2}
\noindent (A) $0$ \\
(B) $1$ \\
(C) $2$ \\
(D) $3$
\end{multicols}\end{quote}

\item An equation of the form $ax^2+bx+c=0$ is called quadratic if:
\begin{quote}\begin{multicols}{2}
\noindent (A) $a=0$ \\
(B) $a\neq0$ \\
(C) $b=0$ \\
(D) $b\neq0$
\end{multicols}\end{quote}

\item Maximum number of roots of a quadratic equation is:
\begin{quote}\begin{multicols}{2}
\noindent (A) $1$ \\
(B) $2$ \\
(C) $3$ \\
(D) $4$
\end{multicols}\end{quote}

\item One root of the equation $x^2-3x+a=0$ is $2$, then $a$ is:
\begin{quote}\begin{multicols}{2}
\noindent (A) $-2$ \\
(B) $2$ \\
(C) $3$ \\
(D) $-3$
\end{multicols}\end{quote}

\item What are the distinct roots of the equation $x^2+x^{-2}-2=0$:
\begin{quote}\begin{multicols}{2}
\noindent (A) $\pm1,\pm i$ \\
(B) $\pm i$ \\
(C) $\pm2$ \\
(D) $\pm1$
\end{multicols}\end{quote}

\item An equation involving a radical expression of the variable is called:
\begin{quote}\begin{multicols}{2}
\noindent (A) Radical equation \\
(B) Algebraic equation \\
(C) Exponential equation \\
(D) None of these
\end{multicols}\end{quote}

\item Name the type of equation $\sqrt{x+a}+\sqrt{x+b}=\sqrt{x+c}$:
\begin{quote}\begin{multicols}{2}
\noindent (A) Radical \\
(B) Reciprocal \\
(C) Rational \\
(D) None of these
\end{multicols}\end{quote}

\item The solution of $\sqrt{y+3}=\sqrt{3y-5}$ is:
\begin{quote}\begin{multicols}{2}
\noindent (A) $2$ \\
(B) $4$ \\
(C) $1$ \\
(D) $-4$
\end{multicols}\end{quote}

\item If a matrix $A$ is of order $m\times n$ then $A$ has number of elements:
\begin{quote}\begin{multicols}{2}
\noindent (A) $m\times n$ \\
(B) $m+n$ \\
(C) $m-n$ \\
(D) $m$
\end{multicols}\end{quote}

\item $[8]$ is a:
\begin{quote}\begin{multicols}{2}
\noindent (A) Square matrix \\
(B) Unit matrix \\
(C) Scalar matrix \\
(D) Rectangular matrix
\end{multicols}\end{quote}

\item A matrix which has only one row is called:
\begin{quote}\begin{multicols}{2}
\noindent (A) Row matrix \\
(B) Column matrix \\
(C) Square matrix \\
(D) None of these
\end{multicols}\end{quote}

\item If order $(A)=m\times n$, order $(B)=n\times p$, then order $(AB)^t$ is:
\begin{quote}\begin{multicols}{2}
\noindent (A) $m\times m$ \\
(B) $m\times n$ \\
(C) $p\times m$ \\
(D) $p\times n$
\end{multicols}\end{quote}

\item The $\begin{bmatrix}1&0\\0&2\end{bmatrix}$ matrix is:
\begin{quote}\begin{multicols}{2}
\noindent (A) Null \\
(B) Identity \\
(C) Scalar \\
(D) Diagonal
\end{multicols}\end{quote}

\item Transpose of a diagonal matrix is:
\begin{quote}\begin{multicols}{2}
\noindent (A) Scalar matrix \\
(B) Row matrix \\
(C) Null matrix \\
(D) Diagonal matrix
\end{multicols}\end{quote}

\item The matrix $\begin{bmatrix}0&0&0\\0&2&0\\0&0&2\end{bmatrix}$ is a:
\begin{quote}\begin{multicols}{2}
\noindent (A) Rectangular matrix \\
(B) Diagonal matrix \\
(C) Scalar matrix \\
(D) Hermitian matrix
\end{multicols}\end{quote}

\item If $A$ is a matrix of order $m\times n$ and $B$ is a matrix of order $n\times l$ then what is the order of $A\times B$:
\begin{quote}\begin{multicols}{2}
\noindent (A) $m\times n$ \\
(B) $l\times m$ \\
(C) $l\times n$ \\
(D) $m\times l$
\end{multicols}\end{quote}

\item Given that $S=\begin{bmatrix}3&0\\0&3\end{bmatrix}$ then $S^4$ is given by:
\begin{quote}\begin{multicols}{2}
\noindent (A) $\begin{bmatrix}81&0\\0&81\end{bmatrix}$ \\
(B) $\begin{bmatrix}0&81\\81&0\end{bmatrix}$ \\
(C) $\begin{bmatrix}27&0\\0&27\end{bmatrix}$ \\
(D) $\begin{bmatrix}0&27\\27&0\end{bmatrix}$
\end{multicols}\end{quote}

\item If each element of a matrix is zero, then it is called:
\begin{quote}\begin{multicols}{2}
\noindent (A) Null matrix \\
(B) Identity matrix \\
(C) Scalar matrix \\
(D) Diagonal matrix
\end{multicols}\end{quote}

\item The matrix obtained by interchanging its rows and columns is called:
\begin{quote}\begin{multicols}{2}
\noindent (A) Inverse \\
(B) Transpose \\
(C) Symmetric \\
(D) Adjoint
\end{multicols}\end{quote}

\item Transpose of a row matrix is a:
\begin{quote}\begin{multicols}{2}
\noindent (A) Square matrix \\
(B) Column matrix \\
(C) Symmetric matrix \\
(D) Skew symmetric matrix
\end{multicols}\end{quote}

\item If the order of a matrix $A$ is $m\times n$, then the order of $A^t$ is:
\begin{quote}\begin{multicols}{2}
\noindent (A) $m\times n$ \\
(B) $n\times n$ \\
(C) $n\times m$ \\
(D) $m\times m$
\end{multicols}\end{quote}

\item If $A=[a\ b\ c]$ then order of $A^t$ is:
\begin{quote}\begin{multicols}{2}
\noindent (A) $1\times3$ \\
(B) $3\times1$ \\
(C) $3\times3$ \\
(D) $1\times1$
\end{multicols}\end{quote}

\item If $A$ is a matrix of order $3\times2$ then order of $A^tA$ is:
\begin{quote}\begin{multicols}{2}
\noindent (A) $3\times2$ \\
(B) $2\times3$ \\
(C) $2\times2$ \\
(D) $3\times3$
\end{multicols}\end{quote}

\item If $A=\begin{bmatrix}1&2\\4&1\end{bmatrix}$, $B=\begin{bmatrix}-1&3\\5&2\end{bmatrix}$ then $A+B=$:
\begin{quote}\begin{multicols}{2}
\noindent (A) $\begin{bmatrix}0&5\\9&3\end{bmatrix}$ \\
(B) $\begin{bmatrix}1&5\\6&7\end{bmatrix}$ \\
(C) $\begin{bmatrix}1&0\\0&1\end{bmatrix}$ \\
(D) $\begin{bmatrix}2&2\\1&1\end{bmatrix}$
\end{multicols}\end{quote}

\item If $A=\begin{bmatrix}2\\-1\end{bmatrix}$ and $B=[5\ 0]$, then $AB$ is equal to:
\begin{quote}\begin{multicols}{2}
\noindent (A) $\begin{bmatrix}10&0\\-5&0\end{bmatrix}$ \\
(B) $\begin{bmatrix}10\\0\end{bmatrix}$ \\
(C) $[0\ 0]$ \\
(D) $\begin{bmatrix}10\\-5\end{bmatrix}$
\end{multicols}\end{quote}

\item If $\begin{bmatrix}x+3&1\\-3&3y-4\end{bmatrix}=\begin{bmatrix}2&1\\-3&2\end{bmatrix}$ then the value of $y$ will be:
\begin{quote}\begin{multicols}{2}
\noindent (A) $2$ \\
(B) $-2$ \\
(C) $4$ \\
(D) $-4$
\end{multicols}\end{quote}

\item A matrix of order $1\times n$ is called:
\begin{quote}\begin{multicols}{2}
\noindent (A) Unit matrix \\
(B) Diagonal matrix \\
(C) Column matrix \\
(D) Row matrix
\end{multicols}\end{quote}

\item If the matrix $\begin{bmatrix}\lambda&4\\1&2\end{bmatrix}$ is singular then $\lambda$:
\begin{quote}\begin{multicols}{2}
\noindent (A) $4$ \\
(B) $0$ \\
(C) $2$ \\
(D) $1$
\end{multicols}\end{quote}

\item If $\begin{bmatrix}x&1\\1&1\end{bmatrix}$ and $|A|=1$, then $x=$:
\begin{quote}\begin{multicols}{2}
\noindent (A) $0$ \\
(B) $1$ \\
(C) $2$ \\
(D) $3$
\end{multicols}\end{quote}

\item The matrix $\begin{bmatrix}5&1\\15&3\end{bmatrix}$ is:
\begin{quote}\begin{multicols}{2}
\noindent (A) Non-singular \\
(B) Symmetric \\
(C) Skew symmetric \\
(D) Singular
\end{multicols}\end{quote}

\item The value of the determinant of an identity matrix is equal to:
\begin{quote}\begin{multicols}{2}
\noindent (A) $0$ \\
(B) $1$ \\
(C) $-1$ \\
(D) $2$
\end{multicols}\end{quote}

\item If $A$ is a non-singular square matrix, then $AA^{-1}$ equals:
\begin{quote}\begin{multicols}{2}
\noindent (A) $A$ \\
(B) $A^{-1}$ \\
(C) $O$ \\
(D) $1$
\end{multicols}\end{quote}

\item If $A=\begin{bmatrix}2&1\\6&3\end{bmatrix}$ then the cofactor of $6$ is:
\begin{quote}\begin{multicols}{2}
\noindent (A) $1$ \\
(B) $-1$ \\
(C) $-6$ \\
(D) $3$
\end{multicols}\end{quote}

\item If the matrices $A$ and $B$ are conformable for multiplication then:
\begin{quote}\begin{multicols}{2}
\noindent (A) $AB=BA$ \\
(B) $AB=A$ \\
(C) $AB=B$ \\
(D) None
\end{multicols}\end{quote}

\item If $A$ and $B$ are square matrices of the same order then:
\begin{quote}\begin{multicols}{2}
\noindent (A) $(A+B)^2=A^2+2AB+B^2$ \\
(B) $(A+B)^2=A^2+2B\,A+B^2$ \\
(C) Both (a) and (b) \\
(D) None of these
\end{multicols}\end{quote}

\item If $A=\begin{bmatrix}1&-2&3\\-2&3&1\\4&-3&2\end{bmatrix}$ then $A_{33}$ equals:
\begin{quote}\begin{multicols}{2}
\noindent (A) $-1$ \\
(B) $1$ \\
(C) $7$ \\
(D) $-7$
\end{multicols}\end{quote}

\item If $A$ is a square matrix then which is true:
\begin{quote}\begin{multicols}{2}
\noindent (A) $|A|=|-A|$ \\
(B) $|A|=|A^t|$ \\
(C) $|A|=|-A^t|$ \\
(D) $|A^2|=|A|$
\end{multicols}\end{quote}

\item If $|A|=5$ then $|A^t|=$:
\begin{quote}\begin{multicols}{2}
\noindent (A) $-5$ \\
(B) $1/5$ \\
(C) $5$ \\
(D) $0$
\end{multicols}\end{quote}

\item $\begin{vmatrix}2&3&0\\3&9&6\\2&15&1\end{vmatrix}=a\begin{vmatrix}2&1&0\\1&1&2\\2&5&1\end{vmatrix}$ then $a=$:
\begin{quote}\begin{multicols}{2}
\noindent (A) $3$ \\
(B) $6$ \\
(C) $9$ \\
(D) $12$
\end{multicols}\end{quote}

\item Let $A=\begin{bmatrix}1&2&3x\\2&3&6x\\3&5&9x\end{bmatrix}$ then $|A|$ is equal to:
\begin{quote}\begin{multicols}{2}
\noindent (A) $1$ \\
(B) $3$ \\
(C) $2$ \\
(D) $0$
\end{multicols}\end{quote}

\item If $A$ is a square matrix of order $3\times3$ then $|kA|=$:
\begin{quote}\begin{multicols}{2}
\noindent (A) $k|A|$ \\
(B) $k^2|A|$ \\
(C) $k^3|A|$ \\
(D) $k|A|^3$
\end{multicols}\end{quote}

\item The value of $\begin{vmatrix}1&12&25\\0&3&15\\0&0&8\end{vmatrix}$ is:
\begin{quote}\begin{multicols}{2}
\noindent (A) $0$ \\
(B) $1$ \\
(C) $8$ \\
(D) $24$
\end{multicols}\end{quote}

\item $(AB)^t=$:
\begin{quote}\begin{multicols}{2}
\noindent (A) $AB$ \\
(B) $BA$ \\
(C) $A^tB^t$ \\
(D) $B^tA^t$
\end{multicols}\end{quote}

\item $(A^{-1})^t$ is equal to:
\begin{quote}\begin{multicols}{2}
\noindent (A) $A$ \\
(B) $A^{-1}A^t$ \\
(C) $(A^t)^{-1}$ \\
(D) $-A^t$
\end{multicols}\end{quote}

\item The inverse of a square matrix exists if $A$ is:
\begin{quote}\begin{multicols}{2}
\noindent (A) Singular \\
(B) Non-singular \\
(C) Symmetric \\
(D) Rectangular
\end{multicols}\end{quote}

\item Rank of the matrix $\begin{bmatrix}1&0\\0&1\end{bmatrix}$ is:
\begin{quote}\begin{multicols}{2}
\noindent (A) $1$ \\
(B) $2$ \\
(C) $3$ \\
(D) $4$
\end{multicols}\end{quote}

\item The rank of matrix $[1\ 0\ 3]$ is:
\begin{quote}\begin{multicols}{2}
\noindent (A) Zero \\
(B) $1$ \\
(C) $2$ \\
(D) $3$
\end{multicols}\end{quote}

\item Rank of matrix $\begin{bmatrix}1&0&0\\0&1&0\\0&0&0\end{bmatrix}$ is:
\begin{quote}\begin{multicols}{2}
\noindent (A) $1$ \\
(B) $2$ \\
(C) $3$ \\
(D) $4$
\end{multicols}\end{quote}

\item Rank of the matrix $\begin{bmatrix}1\\1\\1\\0\end{bmatrix}$ is:
\begin{quote}\begin{multicols}{2}
\noindent (A) $0$ \\
(B) $1$ \\
(C) $2$ \\
(D) $3$
\end{multicols}\end{quote}

\item If $AX=B$ represents the system of non-homogeneous linear equations then $X=$:
\begin{quote}\begin{multicols}{2}
\noindent (A) $BA^{-1}$ \\
(B) $AB$ \\
(C) $A^{-1}B^{-1}$ \\
(D) $A^{-1}B$
\end{multicols}\end{quote}

\item For non-homogeneous linear equations, system $AX=B$, a solution is possible if:
\begin{quote}\begin{multicols}{2}
\noindent (A) Consistent \\
(B) Inconsistent \\
(C) Unique \\
(D) Homogeneous
\end{multicols}\end{quote}

\item The $(0,0)$ solution of the system $a_1x+b_1y=0$ and $a_2x+b_2y=0$ is called:
\begin{quote}\begin{multicols}{2}
\noindent (A) Consistent \\
(B) Inconsistent \\
(C) Non-trivial \\
(D) Trivial
\end{multicols}\end{quote}

\item The trivial solution of a homogeneous linear equation is:
\begin{quote}\begin{multicols}{2}
\noindent (A) $(1,0,0)$ \\
(B) $(0,1,0)$ \\
(C) $(0,0,1)$ \\
(D) $(0,0,0)$
\end{multicols}\end{quote}

\item For homogeneous linear equations, system $AX=0$, a non-trivial solution is possible if $|A|=$:
\begin{quote}\begin{multicols}{2}
\noindent (A) Non-zero \\
(B) Zero \\
(C) Negative \\
(D) Positive
\end{multicols}\end{quote}

\item A fraction in which degree of numerator is less than degree of denominator is called:
\begin{quote}\begin{multicols}{2}
\noindent (A) Algebraic relation \\
(B) Improper fraction \\
(C) Proper fraction \\
(D) Equation
\end{multicols}\end{quote}

\item A relation in which the equality is true for all finite values of unknowns is called an:
\begin{quote}\begin{multicols}{2}
\noindent (A) Identity \\
(B) Equation \\
(C) Trigonometric equation \\
(D) Algebraic equation
\end{multicols}\end{quote}

\item The improper fraction can be changed into a proper fraction by:
\begin{quote}\begin{multicols}{2}
\noindent (A) Addition \\
(B) Subtraction \\
(C) Multiplication \\
(D) Division
\end{multicols}\end{quote}

\item A fraction in which the degree of numerator is equal to or greater than the degree of denominator is called:
\begin{quote}\begin{multicols}{2}
\noindent (A) Proper fraction \\
(B) An improper fraction \\
(C) Equation \\
(D) Algebraic fraction
\end{multicols}\end{quote}

\item An open sentence formed by using an equality is called:
\begin{quote}\begin{multicols}{2}
\noindent (A) An equation \\
(B) An identity \\
(C) Sentence \\
(D) Partial fraction
\end{multicols}\end{quote}

\item $(x+3)(x+4)=x^2+7x+12$ is:
\begin{quote}\begin{multicols}{2}
\noindent (A) Equation \\
(B) Function \\
(C) Identity \\
(D) Conditional equation
\end{multicols}\end{quote}

\item The conditional equation $5x+4$ is true if $x=$:
\begin{quote}\begin{multicols}{2}
\noindent (A) $4$ \\
(B) $5$ \\
(C) $5/4$ \\
(D) $4/5$
\end{multicols}\end{quote}

\item If $(2x+1)=A(x+1)+B(x+2)$, then $A=$:
\begin{quote}\begin{multicols}{2}
\noindent (A) $3$ \\
(B) $4$ \\
(C) $5$ \\
(D) $1$
\end{multicols}\end{quote}

\item $x^2+x-6=0$ is:
\begin{quote}\begin{multicols}{2}
\noindent (A) Conditional equation \\
(B) Identity \\
(C) Proper fraction \\
(D) Improper fraction
\end{multicols}\end{quote}

\item Conditional equation $3x-1=0$ is truly only:
\begin{quote}\begin{multicols}{2}
\noindent (A) $x=3$ \\
(B) $x=-3$ \\
(C) $x=\dfrac{1}{3}$ \\
(D) $x=-\dfrac{1}{3}$
\end{multicols}\end{quote}

\item $\dfrac{x^2+1}{q(x)}$ will be a proper fraction if degree of $q(x)=$:
\begin{quote}\begin{multicols}{2}
\noindent (A) $0$ \\
(B) $3$ \\
(C) $2$ \\
(D) $1$
\end{multicols}\end{quote}

\item $\dfrac{x^3+2x-6}{Q(x)}$ is improper if:
\begin{quote}\begin{multicols}{2}
\noindent (A) Degree of $Q(x)=2$ \\
(B) Degree of $Q(x)=3$ \\
(C) Degree of $Q(x)=4$ \\
(D) Degree of $Q(x)=5$
\end{multicols}\end{quote}

\item If $\dfrac{7x+25}{(x+3)(x+4)}=\dfrac{A}{x+3}+\dfrac{B}{x+4}$ then the value of $B$ is:
\begin{quote}\begin{multicols}{2}
\noindent (A) $3$ \\
(B) $-3$ \\
(C) $4$ \\
(D) $-4$
\end{multicols}\end{quote}

\item Partial fraction of $\dfrac{1}{x^2-1}$ will be of the form:
\begin{quote}\begin{multicols}{2}
\noindent (A) $\dfrac{Ax+B}{x^2-1}$ \\
(B) $\dfrac{A}{x+1}+\dfrac{B}{x-1}$ \\
(C) $2$ \\
(D) $4$
\end{multicols}\end{quote}

\item $\dfrac{9}{(x+2)^2(x-1)}=\dfrac{A}{x-1}+\dfrac{B}{x+2}+\ldots$:
\begin{quote}\begin{multicols}{2}
\noindent (A) $\dfrac{C}{x+2}$ \\
(B) $\dfrac{C}{(x+2)^2}$ \\
(C) $\dfrac{C}{x-2}$ \\
(D) None of these
\end{multicols}\end{quote}

\item The fraction $\dfrac{x+1}{x^2+2}$ is:
\begin{quote}\begin{multicols}{2}
\noindent (A) Proper fraction \\
(B) Improper fraction \\
(C) Identity \\
(D) Mixed
\end{multicols}\end{quote}

\item $\dfrac{P(x)}{x^2+1}$ will be a proper fraction if degree of $P(x)$ equals:
\begin{quote}\begin{multicols}{2}
\noindent (A) $1$ \\
(B) $2$ \\
(C) $3$ \\
(D) $0$
\end{multicols}\end{quote}

\item $\dfrac{x^4}{1-x^4}$ is an:
\begin{quote}\begin{multicols}{2}
\noindent (A) Proper fraction \\
(B) Improper fraction \\
(C) Decimal \\
(D) Equation
\end{multicols}\end{quote}

\item Partial fraction of $\dfrac{x^2+1}{x^3+1}$ is:
\begin{quote}\begin{multicols}{2}
\noindent (A) $\dfrac{A}{x+1}+\dfrac{B}{x^2+x+1}$ \\
(B) $\dfrac{A}{x+1}+\dfrac{Bx+C}{x^2-x+1}$ \\
(C) $\dfrac{A}{x+1}+\dfrac{Bx+C}{x^2+x+1}$ \\
(D) $\dfrac{A}{x+1}+\dfrac{Bx+C}{x^2-x-1}$
\end{multicols}\end{quote}

\item Partial fraction of $\dfrac{x+4}{(x-1)(x^2+2)}$ will be:
\begin{quote}\begin{multicols}{2}
\noindent (A) $\dfrac{A}{x-1}+\dfrac{Bx+C}{x^2-x+1}$ \\
(B) $\dfrac{A}{x-1}+\dfrac{B}{x^2+2}$ \\
(C) $\dfrac{Ax}{x-1}+\dfrac{Bx+C}{x^2+2}$ \\
(D) $\dfrac{Ax}{x-1}+\dfrac{Bx}{x^2+2}$
\end{multicols}\end{quote}

\item Partial fraction of $\dfrac{1}{x^3+1}$ will be of the form:
\begin{quote}\begin{multicols}{2}
\noindent (A) $\dfrac{A}{x+1}+\dfrac{Bx+C}{x^2-x+1}$ \\
(B) $\dfrac{A}{x-1}+\dfrac{Bx+C}{x^2+x+1}$ \\
(C) $\dfrac{A}{x+1}+\dfrac{B}{x^2-x+1}$ \\
(D) $\dfrac{Ax+B}{x^2+1}+\dfrac{C}{x^2-x+1}$
\end{multicols}\end{quote}

\item The expression $\dfrac{x}{x^2+1}-\dfrac{1}{x+1}$ is resolved into partial fraction of:
\begin{quote}\begin{multicols}{2}
\noindent (A) $\dfrac{x+1}{(x^2+1)(x+1)}$ \\
(B) $\dfrac{x}{(x^2+1)(x+1)}$ \\
(C) $\dfrac{x-1}{(x^2+1)(x+1)}$ \\
(D) $\dfrac{2x}{(x^2+1)(x+1)}$
\end{multicols}\end{quote}

\item Partial fraction of $\dfrac{2x+3}{x^2(x^2-1)}$ will be of the form:
\begin{quote}\begin{multicols}{2}
\noindent (A) $\dfrac{Ax+B}{x^2}+\dfrac{Cx+D}{x^2-1}$ \\
(B) $\dfrac{A}{x}+\dfrac{B}{x^2}+\dfrac{C}{x+1}+\dfrac{D}{x-1}$ \\
(C) $\dfrac{Ax+B}{x^2}+\dfrac{C}{x+1}+\dfrac{D}{x-1}$ \\
(D) $\dfrac{A}{x}+\dfrac{B}{x^2}+\dfrac{Cx+D}{x^2-1}$
\end{multicols}\end{quote}

\item Partial fraction of $\dfrac{4x^2+8x}{x^4+2x^2+9}$ will be of the form:
\begin{quote}\begin{multicols}{2}
\noindent (A) $\dfrac{A}{x^4}+\dfrac{B}{2x^2}+\dfrac{C}{9}$ \\
(B) $\dfrac{A}{9}+\dfrac{Bx+C}{x^4+2x^2}$ \\
(C) $\dfrac{Ax+B}{x^2+2x+3}+\dfrac{Cx-D}{x^2-2x+3}$ \\
(D) $\dfrac{Ax+B}{x^2+2x+3}+\dfrac{Cx+D}{x^2-2x}$
\end{multicols}\end{quote}

\item A sequence is a function whose domain is a subset of the set of:
\begin{quote}\begin{multicols}{2}
\noindent (A) Identity \\
(B) Natural numbers \\
(C) Integers \\
(D) None of these
\end{multicols}\end{quote}

\item A sequence is denoted by:
\begin{quote}\begin{multicols}{2}
\noindent (A) $a_n$ \\
(B) $\{a_n\}$ \\
(C) $\{S_n\}$ \\
(D) None of these
\end{multicols}\end{quote}

\item For a sequence, $a_n$ represents:
\begin{quote}\begin{multicols}{2}
\noindent (A) $n$th term \\
(B) General term \\
(C) Tenth term \\
(D) Last term
\end{multicols}\end{quote}

\item If the $n$th term of an A.P. is $\dfrac{1}{2}(3-n)$, then the first three terms are:
\begin{quote}\begin{multicols}{2}
\noindent (A) $3,2,1$ \\
(B) $1,2,3$ \\
(C) $1,2,1$ \\
(D) $1,1/2,0$
\end{multicols}\end{quote}

\item The 5th term of the sequence $a_n=(-1)^n(2n-3)$ is:
\begin{quote}\begin{multicols}{2}
\noindent (A) $8$ \\
(B) $-7$ \\
(C) $-8$ \\
(D) $7$
\end{multicols}\end{quote}

\item The fifth term of the sequence $a_n=2n+1$ is:
\begin{quote}\begin{multicols}{2}
\noindent (A) $7$ \\
(B) $9$ \\
(C) $11$ \\
(D) $13$
\end{multicols}\end{quote}

\item If $a_n=(-1)^{n+1}$ then the 26th term is:
\begin{quote}\begin{multicols}{2}
\noindent (A) $1$ \\
(B) $-1$ \\
(C) $26$ \\
(D) $24$
\end{multicols}\end{quote}

\item If $a_n-a_{n-1}$, $a_1=2$ then $a_3=$:
\begin{quote}\begin{multicols}{2}
\noindent (A) $6$ \\
(B) $4$ \\
(C) $11$ \\
(D) $17$
\end{multicols}\end{quote}

\item Next term of the sequence $1,3,7,15,31,\ldots$ is:
\begin{quote}\begin{multicols}{2}
\noindent (A) $63$ \\
(B) $62$ \\
(C) $80$ \\
(D) $81$
\end{multicols}\end{quote}

\item The next two terms of the sequence $7,9,12,17,\ldots$ are:
\begin{quote}\begin{multicols}{2}
\noindent (A) $21,27$ \\
(B) $21,26$ \\
(C) $20,27$ \\
(D) $20,26$
\end{multicols}\end{quote}

\item The 7th term of the sequence $2,6,11,17,\ldots$ is:
\begin{quote}\begin{multicols}{2}
\noindent (A) $24$ \\
(B) $36$ \\
(C) $30$ \\
(D) $41$
\end{multicols}\end{quote}

\item The $n$th term of the sequence $\dfrac{1}{3},\dfrac{2}{5},\dfrac{3}{7},\ldots$ is:
\begin{quote}\begin{multicols}{2}
\noindent (A) $\dfrac{n}{2n-1}$ \\
(B) $\dfrac{n}{2n+1}$ \\
(C) $\dfrac{n}{3n-1}$ \\
(D) $\dfrac{n}{3n+1}$
\end{multicols}\end{quote}

\item If $a_n=\dfrac{(-1)^{n+1}}{2^n}$, then $a_6=$:
\begin{quote}\begin{multicols}{2}
\noindent (A) $1/8$ \\
(B) $1/16$ \\
(C) $1/32$ \\
(D) $1/64$
\end{multicols}\end{quote}

\item A sequence $\{a_n\}$ is an arithmetic sequence if $\forall n\in N$ and $n>1$:
\begin{quote}\begin{multicols}{2}
\noindent (A) $\dfrac{a_n}{a_{n-1}}$ is same \\
(B) $\dfrac{a_{n-1}}{a_n}$ is same \\
(C) $a_n-a_{n-1}$ is different \\
(D) None of these
\end{multicols}\end{quote}

\item Common difference of arithmetic sequence $3,5,7,\ldots$ is:
\begin{quote}\begin{multicols}{2}
\noindent (A) $-\dfrac{2}{15}$ \\
(B) $\dfrac{2}{15}$ \\
(C) $2$ \\
(D) $\dfrac{1}{2}$
\end{multicols}\end{quote}

\item If $\dfrac{1}{a},\dfrac{1}{b},\dfrac{1}{c}$ are in A.P., then $b=$:
\begin{quote}\begin{multicols}{2}
\noindent (A) $\dfrac{a+c}{2ac}$ \\
(B) $\dfrac{a-c}{2ac}$ \\
(C) $\dfrac{2ac}{a+c}$ \\
(D) $\dfrac{2ac}{a-c}$
\end{multicols}\end{quote}

\item If $\dfrac{1}{a},\dfrac{1}{b},\dfrac{1}{c}$ are in A.P. then the common difference is:
\begin{quote}\begin{multicols}{2}
\noindent (A) $\dfrac{2ac}{a-c}$ \\
(B) $\dfrac{2ac}{a+c}$ \\
(C) $\dfrac{a-c}{2ac}$ \\
(D) $\dfrac{a+c}{2ac}$
\end{multicols}\end{quote}

\item $2,4,6,8,10,\ldots$ is:
\begin{quote}\begin{multicols}{2}
\noindent (A) G.P. \\
(B) A.P. \\
(C) Geometric series \\
(D) None of these
\end{multicols}\end{quote}

\item The $n$th term of an A.P. with usual notation is:
\begin{quote}\begin{multicols}{2}
\noindent (A) $a_1+(n-1)d$ \\
(B) $a_2+(n-1)d$ \\
(C) $a_1+nd$ \\
(D) $a_2+nd$
\end{multicols}\end{quote}

\item The A.P. whose $n$th term is $2n-1$ is:
\begin{quote}\begin{multicols}{2}
\noindent (A) $1,3,6,\ldots$ \\
(B) $2,3,5,\ldots$ \\
(C) $1,3,5,\ldots$ \\
(D) $5,3,1,\ldots$
\end{multicols}\end{quote}

\item If $a_{n-2}=3n-11$ then the $n$th term is:
\begin{quote}\begin{multicols}{2}
\noindent (A) $3n+5$ \\
(B) $3n-3$ \\
(C) $3n-5$ \\
(D) $3n+2$
\end{multicols}\end{quote}

\item If $a_{n-3}=2n-5$, then the 7th term is:
\begin{quote}\begin{multicols}{2}
\noindent (A) $9$ \\
(B) $15$ \\
(C) $11$ \\
(D) $13$
\end{multicols}\end{quote}

\item The 9th term of the sequence $1,-3,5,-7,\ldots$ is:
\begin{quote}\begin{multicols}{2}
\noindent (A) $15$ \\
(B) $-15$ \\
(C) $17$ \\
(D) $-17$
\end{multicols}\end{quote}

\item If $a_1=5$ and other consecutive terms are $23,26,29,\ldots$ then the 3rd term is:
\begin{quote}\begin{multicols}{2}
\noindent (A) $8$ \\
(B) $11$ \\
(C) $14$ \\
(D) $17$
\end{multicols}\end{quote}

\item Which term of $64,60,56,52,\ldots$ is zero:
\begin{quote}\begin{multicols}{2}
\noindent (A) 16th \\
(B) 17th \\
(C) 14th \\
(D) 15th
\end{multicols}\end{quote}

\item The $n$th term of an A.P. is $\dfrac{1}{2}(3n)$; the first three terms are:
\begin{quote}\begin{multicols}{2}
\noindent (A) $3,2,1$ \\
(B) $1,2,3$ \\
(C) $\dfrac{3}{2},3,\dfrac{9}{2}$ \\
(D) $\dfrac{3}{2},9,\dfrac{9}{2}$
\end{multicols}\end{quote}

\item The A.M. between two numbers $a$ and $b$ is:
\begin{quote}\begin{multicols}{2}
\noindent (A) $\dfrac{a+b}{2ab}$ \\
(B) $\dfrac{2ab}{a+b}$ \\
(C) $\dfrac{a+b}{2}$ \\
(D) $\dfrac{2}{a+b}$
\end{multicols}\end{quote}

\item If $a,A,b$ are in A.P., then $2A=$:
\begin{quote}\begin{multicols}{2}
\noindent (A) $\dfrac{a+b}{2}$ \\
(B) $a+b$ \\
(C) $2$ \\
(D) $a-b$
\end{multicols}\end{quote}

\item A.M. between $3\sqrt5$ and $5\sqrt5$ is:
\begin{quote}\begin{multicols}{2}
\noindent (A) $4\sqrt5$ \\
(B) $5\sqrt5$ \\
(C) $10$ \\
(D) $2\sqrt5$
\end{multicols}\end{quote}

\item Arithmetic mean between two numbers $\dfrac{1}{a}$ and $\dfrac{1}{b}$ is:
\begin{quote}\begin{multicols}{2}
\noindent (A) $\dfrac{a+b}{ab}$ \\
(B) $\dfrac{a+b}{2ab}$ \\
(C) $\dfrac{2ab}{a+b}$ \\
(D) $\dfrac{ab}{a+b}$
\end{multicols}\end{quote}

\item The arithmetic mean between $2+\sqrt{2}$ and $2-\sqrt{2}$ is:
\begin{quote}\begin{multicols}{2}
\noindent (A) 2 \\
(B) $\sqrt{2}$ \\
(C) 0 \\
(D) 4
\end{multicols}\end{quote}

\item A.M between two numbers $-2$ and $-6$ is:
\begin{quote}\begin{multicols}{2}
\noindent (A) 4 \\
(B) $-4$ \\
(C) $\pm\sqrt{12}$ \\
(D) $-3$
\end{multicols}\end{quote}

\item The A.M between $x-3$ and $x+5$ is:
\begin{quote}\begin{multicols}{2}
\noindent (A) $x+1$ \\
(B) $x-1$ \\
(C) $x-3$ \\
(D) $x+5$
\end{multicols}\end{quote}

\item If 5 is A.M between $-5$ and $b$, then $b$ is equal to:
\begin{quote}\begin{multicols}{2}
\noindent (A) 0 \\
(B) 15 \\
(C) $-5$ \\
(D) 10
\end{multicols}\end{quote}

\item For what value of $n$, $\dfrac{a^n+b^n}{a^{n-1}+b^{n-1}}$ is the A.M between $a$ and $b$:
\begin{quote}\begin{multicols}{2}
\noindent (A) $-1$ \\
(B) 1 \\
(C) $\dfrac{1}{2}$ \\
(D) 0
\end{multicols}\end{quote}

\item The sum of $n$ arithmetic means between $a$ and $b$ is:
\begin{quote}\begin{multicols}{2}
\noindent (A) $\dfrac{n(a+b)}{2}$ \\
(B) $\dfrac{a+b}{2}$ \\
(C) $\dfrac{2ab}{a+b}$ \\
(D) $ab$
\end{multicols}\end{quote}

\item A.M between $\dfrac{a}{2}$ and $\dfrac{2}{a}$ is:
\begin{quote}\begin{multicols}{2}
\noindent (A) $\dfrac{a^2+4}{a}$ \\
(B) $\dfrac{a^2+4}{2}$ \\
(C) $\dfrac{a^2+4}{2a}$ \\
(D) $\dfrac{a^2+4}{4a}$
\end{multicols}\end{quote}

\item Sum of terms of an A.P is called:
\begin{quote}\begin{multicols}{2}
\noindent (A) Geometric series \\
(B) Arithmetic series \\
(C) Arithmetic progression \\
(D) Harmonic progression
\end{multicols}\end{quote}

\item Sum of $n$ terms of an A.P is:
\begin{quote}\begin{multicols}{2}
\noindent (A) $\dfrac{a(1-r^n)}{1-r}$ \\
(B) $\dfrac{n}{2}\{2a+(n-1)d\}$ \\
(C) $\dfrac{2ab}{a+b}$ \\
(D) $\dfrac{a+b}{2}$
\end{multicols}\end{quote}

\item 21st term of $2+4+6+\ldots$ is:
\begin{quote}\begin{multicols}{2}
\noindent (A) 40 \\
(B) 42 \\
(C) 44 \\
(D) $-42$
\end{multicols}\end{quote}

\item Sum of the A.P $-7+(-5)+(-3)$ upto 6 terms is:
\begin{quote}\begin{multicols}{2}
\noindent (A) $-12$ \\
(B) 12 \\
(C) 9 \\
(D) $-9$
\end{multicols}\end{quote}

\item An arithmetic progression has a first term 12 and a fifth term 18. The sum of the first 25 terms is:
\begin{quote}\begin{multicols}{2}
\noindent (A) 400 \\
(B) 350 \\
(C) 750 \\
(D) 150
\end{multicols}\end{quote}

\item In an arithmetic series, $S_n$ equals to:
\begin{quote}\begin{multicols}{2}
\noindent (A) $S_{n-1}+a_n$ \\
(B) $S_{n-1}-a_n$ \\
(C) $S_{n-1}\times a_n$ \\
(D) None of these
\end{multicols}\end{quote}

\item The sum of $n$ terms of an A.P is $n^2$, then its common difference is:
\begin{quote}\begin{multicols}{2}
\noindent (A) 4 \\
(B) 2 \\
(C) $-2$ \\
(D) 1
\end{multicols}\end{quote}

\item What is the sum of $n$ terms with $n$th term $a_n=4n+1$:
\begin{quote}\begin{multicols}{2}
\noindent (A) $2n(2n+3)$ \\
(B) $n(2n+3)$ \\
(C) $2n+3$ \\
(D) $4n+6$
\end{multicols}\end{quote}

\item In a geometric sequence, the $n$th term equals:
\begin{quote}\begin{multicols}{2}
\noindent (A) $a_1+(n-1)d$ \\
(B) $\dfrac{n}{2}[2a_1+(n-1)d]$ \\
(C) $a_1r^{n-1}$ \\
(D) $\dfrac{a_1}{1-r}$
\end{multicols}\end{quote}

\item The common ratio of a G.P cannot be:
\begin{quote}\begin{multicols}{2}
\noindent (A) 3 \\
(B) 1 \\
(C) 2 \\
(D) 0
\end{multicols}\end{quote}

\item The sequence $3,6,12$ is:
\begin{quote}\begin{multicols}{2}
\noindent (A) A.P \\
(B) G.P \\
(C) H.P \\
(D) Arithmetic series
\end{multicols}\end{quote}

\item The sequence $4,12,36,\ldots$ is:
\begin{quote}\begin{multicols}{2}
\noindent (A) A.P \\
(B) G.P \\
(C) H.P \\
(D) Series
\end{multicols}\end{quote}

\item The sequence $5,\dfrac{5}{2},\dfrac{5}{4},\ldots$ is:
\begin{quote}\begin{multicols}{2}
\noindent (A) Series \\
(B) A.P \\
(C) G.P \\
(D) H.P
\end{multicols}\end{quote}

\item No term of a geometric sequence can be:
\begin{quote}\begin{multicols}{2}
\noindent (A) 0 \\
(B) 1 \\
(C) 2 \\
(D) 3
\end{multicols}\end{quote}

\item Which of the given cannot be any term of a G.P:
\begin{quote}\begin{multicols}{2}
\noindent (A) $-1$ \\
(B) 0 \\
(C) 1 \\
(D) 5
\end{multicols}\end{quote}

\item If $\dfrac{1}{a},\dfrac{1}{b}$ and $\dfrac{1}{c}$ are in G.P then the common ratio is equal to:
\begin{quote}\begin{multicols}{2}
\noindent (A) $\pm\sqrt{\dfrac{c}{a}}$ \\
(B) $\pm\sqrt{\dfrac{a}{c}}$ \\
(C) $a+b$ \\
(D) None
\end{multicols}\end{quote}

\item If $a_1$ and $r$ are the first term and common ratio respectively, then the $(n+1)$th term of a G.P is:
\begin{quote}\begin{multicols}{2}
\noindent (A) $a_1r^{n-1}$ \\
(B) $a_1r^{n+1}$ \\
(C) $a_1r^{n+2}$ \\
(D) $a_1r^{n}$
\end{multicols}\end{quote}

\item The next term of the G.P $1,2,4,8,16,\ldots$ is:
\begin{quote}\begin{multicols}{2}
\noindent (A) 15 \\
(B) 19 \\
(C) 21 \\
(D) 32
\end{multicols}\end{quote}

\item $a_5$ in the G.P $3,6,12,\ldots$ will be:
\begin{quote}\begin{multicols}{2}
\noindent (A) 48 \\
(B) 64 \\
(C) 128 \\
(D) 224
\end{multicols}\end{quote}

\item If $a,b,c$ are in G.P and $a>0,b>0,c>0$, then the reciprocals of $a,b,c$ form:
\begin{quote}\begin{multicols}{2}
\noindent (A) A.P \\
(B) G.P \\
(C) H.P \\
(D) None of these
\end{multicols}\end{quote}

\item A geometric mean between $a$ and $b$ is:
\begin{quote}\begin{multicols}{2}
\noindent (A) $\dfrac{a+b}{2}$ \\
(B) $\dfrac{2}{a+b}$ \\
(C) $\sqrt{ab}$ \\
(D) $\dfrac{2ab}{a+b}$
\end{multicols}\end{quote}

\item If $a,G,b$ be the G.P where $a,b$ are numbers then G.M is equal to:
\begin{quote}\begin{multicols}{2}
\noindent (A) $\sqrt{ab}$ \\
(B) $-\sqrt{ab}$ \\
(C) $\pm\sqrt{ab}$ \\
(D) $ab$
\end{multicols}\end{quote}

\item A number $G$ is said to be the G.M between $a$ and $b$ if $a,G,b$ are in:
\begin{quote}\begin{multicols}{2}
\noindent (A) A.P \\
(B) H.P \\
(C) G.P \\
(D) G.M
\end{multicols}\end{quote}

\item G.M between $\dfrac{1}{a}$ and $\dfrac{1}{b}$ is:
\begin{quote}\begin{multicols}{2}
\noindent (A) $\dfrac{-1}{ab}$ \\
(B) $\sqrt{\dfrac{-1}{ab}}$ \\
(C) $ab$ \\
(D) None
\end{multicols}\end{quote}

\item The geometric mean between $-2i$ and $8i$ is:
\begin{quote}\begin{multicols}{2}
\noindent (A) $\pm4$ \\
(B) $\pm3$ \\
(C) $\pm2$ \\
(D) $\pm1$
\end{multicols}\end{quote}

\item The G.M between 1 and 16 is:
\begin{quote}\begin{multicols}{2}
\noindent (A) $-4$ \\
(B) 4 \\
(C) $\pm4$ \\
(D) $\dfrac{1}{4}$
\end{multicols}\end{quote}

\item G.M between two numbers $-2$ and $4$ is:
\begin{quote}\begin{multicols}{2}
\noindent (A) $\pm\sqrt{8}$ \\
(B) $\pm2\sqrt{2}$ \\
(C) $\pm2i\sqrt{2}$ \\
(D) $\pm2$
\end{multicols}\end{quote}

\item The geometric means between 2 and 4 are:
\begin{quote}\begin{multicols}{2}
\noindent (A) $\pm8\sqrt{2}$ \\
(B) $\pm4\sqrt{2}$ \\
(C) $\pm2\sqrt{2}$ \\
(D) $\pm3\sqrt{2}$
\end{multicols}\end{quote}

\item If $a=2i,\ b=4i$, then $G=$:
\begin{quote}\begin{multicols}{2}
\noindent (A) $\pm2\sqrt{2}\,i$ \\
(B) $\sqrt{6}\,i$ \\
(C) $\pm2i$ \\
(D) $4i$
\end{multicols}\end{quote}

\item For what value of $n$, $\dfrac{a^n+b^n}{a^{n-1}+b^{n-1}}$ is the positive geometric mean between $a$ and $b$:
\begin{quote}\begin{multicols}{2}
\noindent (A) $n=-1$ \\
(B) 1 \\
(C) $\dfrac{1}{2}$ \\
(D) 0
\end{multicols}\end{quote}

\item Sum of $n$ terms of a G.P with first term ``a'' and common ratio $r$ is:
\begin{quote}\begin{multicols}{2}
\noindent (A) $\dfrac{n}{2}\{2a+(n-1)d\}$ \\
(B) $\dfrac{a(1-r^n)}{1-r}$ \\
(C) $\dfrac{2ab}{a+b}$ \\
(D) $\dfrac{a+b}{2}$
\end{multicols}\end{quote}

\item The formula $S_n=\dfrac{a(r^n-1)}{r-1}$ is used for the sum of $n$ terms of a G.P if:
\begin{quote}\begin{multicols}{2}
\noindent (A) $|r|=1$ \\
(B) $|r|>1$ \\
(C) $|r|<2$ \\
(D) $|r|=0$
\end{multicols}\end{quote}

\item An infinite geometric series is convergent if:
\begin{quote}\begin{multicols}{2}
\noindent (A) $|r|<1$ \\
(B) $|r|>1$ \\
(C) $|r|<2$ \\
(D) $|r|\leq1$
\end{multicols}\end{quote}

\item The series $a+ar+ar^2+\ldots$ is convergent if:
\begin{quote}\begin{multicols}{2}
\noindent (A) $|r|>1$ \\
(B) $|r|<1$ \\
(C) $|r|=1$ \\
(D) None
\end{multicols}\end{quote}

\item For an infinite geometric series for which $|r|<1$, $S_n=$:
\begin{quote}\begin{multicols}{2}
\noindent (A) $\dfrac{a_1(1+r)}{1-r}$ \\
(B) $\dfrac{a_1}{1-r}$ \\
(C) $\dfrac{a_1}{2r}$ \\
(D) $\dfrac{a_1}{2-r}$
\end{multicols}\end{quote}

\item The series $1+2x+4x^2+8x^3+\ldots$ will be convergent only if:
\begin{quote}\begin{multicols}{2}
\noindent (A) $|x|<1/2$ \\
(B) $|x|>1/2$ \\
(C) $|x|>2$ \\
(D) $|x|>2$
\end{multicols}\end{quote}

\item The series $1+\dfrac{x}{2}+\dfrac{x^2}{4}+\dfrac{x^3}{8}+\ldots$ if $0<x<2$ is:
\begin{quote}\begin{multicols}{2}
\noindent (A) Arithmetic \\
(B) Harmonic \\
(C) Convergent \\
(D) Divergent
\end{multicols}\end{quote}

\item Sum of an infinite geometric series is valid if:
\begin{quote}\begin{multicols}{2}
\noindent (A) $r<1$ \\
(B) $|r|<1$ \\
(C) $|r|=1$ \\
(D) $|r|>1$
\end{multicols}\end{quote}

\item What is the sum of the infinite G.P $2,\sqrt{2},1,\ldots$?
\begin{quote}\begin{multicols}{2}
\noindent (A) $4-\sqrt{2}$ \\
(B) $4+2\sqrt{2}$ \\
(C) $2\sqrt{2}$ \\
(D) $2+2\sqrt{2}$
\end{multicols}\end{quote}

\item Sum of the series $1+\dfrac{1}{2}+\dfrac{1}{4}+\ldots$ is:
\begin{quote}\begin{multicols}{2}
\noindent (A) $\dfrac{3}{2}$ \\
(B) $\dfrac{1}{3}$ \\
(C) 2 \\
(D) $\dfrac{2}{3}$
\end{multicols}\end{quote}

\item $1+\dfrac{1}{3}+\dfrac{1}{9}+\ldots=$?
\begin{quote}\begin{multicols}{2}
\noindent (A) $\dfrac{3}{4}$ \\
(B) $\dfrac{3}{2}$ \\
(C) $\dfrac{4}{3}$ \\
(D) 3
\end{multicols}\end{quote}

\item If $\dfrac{1}{5}+\dfrac{1}{25}+\dfrac{1}{125}+\ldots$ is an infinite geometric series, then the sum is:
\begin{quote}\begin{multicols}{2}
\noindent (A) $\dfrac{1}{4}$ \\
(B) $\dfrac{5}{4}$ \\
(C) $\dfrac{4}{5}$ \\
(D) $\dfrac{3}{4}$
\end{multicols}\end{quote}

\item Which one of the following infinite geometric series is convergent:
\begin{quote}\begin{multicols}{2}
\noindent (A) $1+\dfrac{1}{2}+\dfrac{1}{4}+\ldots$ \\
(B) $1+2+4+\ldots$ \\
(C) $1+3+9+\ldots$ \\
(D) $3+6+9+\ldots$
\end{multicols}\end{quote}

\item $S_n=\dfrac{a_1(1-r^n)}{1-r}$ is convergent if:
\begin{quote}\begin{multicols}{2}
\noindent (A) $|r|<1$ \\
(B) $|r|>1$ \\
(C) $r=1$ \\
(D) $r=-1$
\end{multicols}\end{quote}

\item If $y=\dfrac{2}{3}x+\dfrac{4}{9}x^2+\dfrac{8}{27}x^3+\ldots$, then the interval of convergence is:
\begin{quote}\begin{multicols}{2}
\noindent (A) $0<x<-\dfrac{3}{2}$ \\
(B) $-\dfrac{3}{2}<x<0$ \\
(C) $\dfrac{1}{2}<x<-\dfrac{1}{2}$ \\
(D) $0<x<\dfrac{3}{2}$
\end{multicols}\end{quote}

\item The series $1+\dfrac{x}{2}+\dfrac{x^2}{4}+\ldots$ is convergent if:
\begin{quote}\begin{multicols}{2}
\noindent (A) $x\in R$ \\
(B) $x\in[-2,2]$ \\
(C) $x\in[2,2]$ \\
(D) $x\in Z$
\end{multicols}\end{quote}

\item Product of an A.P and a G.P is called:
\begin{quote}\begin{multicols}{2}
\noindent (A) A.P \\
(B) G.P \\
(C) H.P \\
(D) A.G.P
\end{multicols}\end{quote}

\item If $|r|<1$ then $r^n\to$:
\begin{quote}\begin{multicols}{2}
\noindent (A) $-1$ \\
(B) 1 \\
(C) 0 \\
(D) None
\end{multicols}\end{quote}

\item The reciprocal of the term of an H.P is:
\begin{quote}\begin{multicols}{2}
\noindent (A) A.P \\
(B) G.P \\
(C) H.P \\
(D) None of these
\end{multicols}\end{quote}

\item No term of a harmonic sequence can be:
\begin{quote}\begin{multicols}{2}
\noindent (A) 3 \\
(B) 2 \\
(C) 1 \\
(D) 2
\end{multicols}\end{quote}

\item The harmonic mean between two numbers $\dfrac{1}{a}$ and $\dfrac{1}{b}$ is:
\begin{quote}\begin{multicols}{2}
\noindent (A) 3 \\
(B) 4 \\
(C) 2 \\
(D) 1
\end{multicols}\end{quote}

\item H.M between 3 and 7 is:
\begin{quote}\begin{multicols}{2}
\noindent (A) $\dfrac{5}{21}$ \\
(B) $\dfrac{21}{5}$ \\
(C) 5 \\
(D) 51
\end{multicols}\end{quote}

\item $\dfrac{1}{k},\dfrac{1}{2k+1},\dfrac{1}{4k-1}$ are in H.P then $k=$:
\begin{quote}\begin{multicols}{2}
\noindent (A) 3 \\
(B) 4 \\
(C) 2 \\
(D) 1
\end{multicols}\end{quote}

\item 9th term in the sequence $\dfrac{1}{3},\dfrac{1}{5},\dfrac{1}{7},\ldots$ is:
\begin{quote}\begin{multicols}{2}
\noindent (A) $\dfrac{1}{13}$ \\
(B) $\dfrac{1}{15}$ \\
(C) $\dfrac{1}{17}$ \\
(D) $\dfrac{1}{19}$
\end{multicols}\end{quote}

\item The 10th term of $\dfrac{1}{2},\dfrac{1}{5},\dfrac{1}{8},\ldots$ is:
\begin{quote}\begin{multicols}{2}
\noindent (A) 30 \\
(B) 28 \\
(C) $\dfrac{1}{29}$ \\
(D) $\dfrac{1}{32}$
\end{multicols}\end{quote}

\item $1,\dfrac{1}{2},\dfrac{1}{3},\dfrac{1}{4},\ldots$ is:
\begin{quote}\begin{multicols}{2}
\noindent (A) A.P \\
(B) G.P \\
(C) H.P \\
(D) Arithmetic series
\end{multicols}\end{quote}

\item If $\dfrac{4}{7}$ be the third term of an H.P, then the third term of the corresponding A.P is:
\begin{quote}\begin{multicols}{2}
\noindent (A) $\dfrac{7}{4}$ \\
(B) $\dfrac{4}{7}$ \\
(C) 4 \\
(D) 7
\end{multicols}\end{quote}

\item If $a=1,\ b=5$ then $A\times H$ is equal to:
\begin{quote}\begin{multicols}{2}
\noindent (A) 5 \\
(B) $-5$ \\
(C) $\dfrac{-5}{2}$ \\
(D) $\dfrac{2}{5}$
\end{multicols}\end{quote}

\item If $A,G,H$ have their usual meanings and $a,b$ are positive distinct real numbers with $G>0$, then:
\begin{quote}\begin{multicols}{2}
\noindent (A) $A>G>H$ \\
(B) $A>H$ \\
(C) $H>G$ \\
(D) $G>H>A$
\end{multicols}\end{quote}

\item $A,G,H$ are in:
\begin{quote}\begin{multicols}{2}
\noindent (A) G.P \\
(B) A.P \\
(C) H.P \\
(D) All of these
\end{multicols}\end{quote}

\item For two positive real numbers, with usual notation:
\begin{quote}\begin{multicols}{2}
\noindent (A) $A>G$ \\
(B) $A=G$ \\
(C) $A<G$ \\
(D) $A\leq G$
\end{multicols}\end{quote}

\item For two negative real numbers, with usual notation:
\begin{quote}\begin{multicols}{2}
\noindent (A) $A>G$ \\
(B) $A=G$ \\
(C) $A<G$ \\
(D) $A\leq G$
\end{multicols}\end{quote}

\item Let $A,G,H$ be the arithmetic, geometric and harmonic means between $a,b$ respectively, then $G^2=$:
\begin{quote}\begin{multicols}{2}
\noindent (A) $A\times H$ \\
(B) $\sqrt{ab}$ \\
(C) $\dfrac{A}{H}$ \\
(D) $A+H$
\end{multicols}\end{quote}

\item If $A,G,H$ are the arithmetic, geometric and harmonic means, then $G/H=$:
\begin{quote}\begin{multicols}{2}
\noindent (A) $\dfrac{A}{G}$ \\
(B) $\dfrac{G}{A}$ \\
(C) $\dfrac{H}{G}$ \\
(D) $G^2$
\end{multicols}\end{quote}

\item If $x,y,z$ are in a harmonic sequence, then the value of $z$ is:
\begin{quote}\begin{multicols}{2}
\noindent (A) $\dfrac{xy}{2x-y}$ \\
(B) $\dfrac{x+y}{2y}$ \\
(C) $\dfrac{x-y}{2}$ \\
(D) $\dfrac{x+y}{3}$
\end{multicols}\end{quote}

\item Reciprocal of the terms $1/3,\ 1/5,\ldots$ form:
\begin{quote}\begin{multicols}{2}
\noindent (A) Geometric sequence \\
(B) Arithmetic sequence \\
(C) Harmonic sequence \\
(D) None
\end{multicols}\end{quote}

\item If $a,b$ are negative and the G.M is also negative:
\begin{quote}\begin{multicols}{2}
\noindent (A) H \\
(B) A \\
(C) G \\
(D) A
\end{multicols}\end{quote}

\item If $\dfrac{1}{5},\dfrac{1}{8}$ are two H.M's between $\dfrac{1}{2}$ and $b$, then $b$ equals:
\begin{quote}\begin{multicols}{2}
\noindent (A) $\dfrac{1}{11}$ \\
(B) $\dfrac{1}{13}$ \\
(C) $\dfrac{1}{15}$ \\
(D) $\dfrac{1}{3}$
\end{multicols}\end{quote}

\item $S_n=1+2+3+\ldots+n$ can also be shown as:
\begin{quote}\begin{multicols}{2}
\noindent (A) $\displaystyle\sum_{k=1}^n k$ \\
(B) $\displaystyle\sum_{k=1}^n \dfrac{1}{k}$ \\
(C) Both a and b \\
(D) None of these
\end{multicols}\end{quote}

\item If $n$ is any positive integer, then $1^3+2^3+3^3+\ldots+n^3$ equals:
\begin{quote}\begin{multicols}{2}
\noindent (A) $\dfrac{n(n+1)}{2}$ \\
(B) $\left[\dfrac{n(n+1)}{2}\right]^2$ \\
(C) $\dfrac{n^2(n+1)^2}{4}$ \\
(D) $\dfrac{n^2(n+1)}{4}$
\end{multicols}\end{quote}

\item $\displaystyle\sum_{k=1}^n 1$ is equal to:
\begin{quote}\begin{multicols}{2}
\noindent (A) 1 \\
(B) $n$ \\
(C) $n^3$ \\
(D) $n^2$
\end{multicols}\end{quote}

\item $\displaystyle\sum_{k=1}^n k=$ (or sum of first $n$ terms):
\begin{quote}\begin{multicols}{2}
\noindent (A) $\dfrac{n^2(n+1)^2}{4}$ \\
(B) $\dfrac{n(n+1)}{2}$ \\
(C) $\dfrac{n(n+1)(2n+1)}{6}$ \\
(D) $\dfrac{n(n-1)}{2}$
\end{multicols}\end{quote}

\item $\displaystyle\sum_{k=1}^n k^2=1^2+2^2+3^2+\ldots+n^2=$:
\begin{quote}\begin{multicols}{2}
\noindent (A) $\dfrac{n(n+1)}{2}$ \\
(B) $\dfrac{n(n+1)(2n+1)}{6}$ \\
(C) $\dfrac{n(n+1)(n+2)}{6}$ \\
(D) $\left[\dfrac{n(n+1)}{2}\right]^2$
\end{multicols}\end{quote}

\item $\displaystyle\sum_{k=1}^n k^3$ is equal to:
\begin{quote}\begin{multicols}{2}
\noindent (A) $\left[\dfrac{n(n+1)}{2!}\right]$ \\
(B) $\left[\dfrac{n(n+1)}{2}\right]^2$ \\
(C) $\dfrac{n(n-1)}{2}$ \\
(D) $\dfrac{n(n+1)(2n+1)}{6}$
\end{multicols}\end{quote}

\item Riaz borrowed 50{,}000 at an interest of 8\% per year for 11 years; the first year's interest will be:
\begin{quote}\begin{multicols}{2}
\noindent (A) 2000 \\
(B) 8000 \\
(C) 6000 \\
(D) 4000
\end{multicols}\end{quote}

\item If $n$ is a positive integer then the factorial of $n$ is denoted by:
\begin{quote}\begin{multicols}{2}
\noindent (A) $\bar{n}$ \\
(B) $\underline{n}$ \\
(C) $\hat{n}$ \\
(D) $n!$
\end{multicols}\end{quote}

\item $0!=$:
\begin{quote}\begin{multicols}{2}
\noindent (A) 0 \\
(B) 1 \\
(C) $-1$ \\
(D) 2
\end{multicols}\end{quote}

\item $6!+0!$ is equal to:
\begin{quote}\begin{multicols}{2}
\noindent (A) 719 \\
(B) 720 \\
(C) 721 \\
(D) 722
\end{multicols}\end{quote}

\item $\dfrac{3!}{0!}$ equals:
\begin{quote}\begin{multicols}{2}
\noindent (A) 3 \\
(B) 6 \\
(C) 12 \\
(D) None
\end{multicols}\end{quote}

\item The factorial form of $6\cdot5\cdot4$ is:
\begin{quote}\begin{multicols}{2}
\noindent (A) $\dfrac{6!}{3!}$ \\
(B) $6!$ \\
(C) $3!$ \\
(D) None
\end{multicols}\end{quote}

\item The value of $4!\cdot0!\cdot1!$ is:
\begin{quote}\begin{multicols}{2}
\noindent (A) 0 \\
(B) 1 \\
(C) 4 \\
(D) 24
\end{multicols}\end{quote}

\item $(n-1)(n-2)(n-3)\ldots(n-r+1)=$:
\begin{quote}\begin{multicols}{2}
\noindent (A) $\dfrac{(n-1)!}{(n-r)!}$ \\
(B) $\dfrac{n!}{(n-r)!}$ \\
(C) $\dfrac{(n-1)!}{(n-r+2)!}$ \\
(D) $\dfrac{n!}{(n-r+1)!}$
\end{multicols}\end{quote}

\item Factorial form of $n(n-1)(n-2)=$:
\begin{quote}\begin{multicols}{2}
\noindent (A) $\dfrac{n!}{(n-1)!}$ \\
(B) $\dfrac{n!}{(n-2)!}$ \\
(C) $\dfrac{n!}{(n-3)!}$ \\
(D) $\dfrac{n!}{(n-4)!}$
\end{multicols}\end{quote}

\item Factorial form of $\dfrac{(n+1)(n)(n-1)}{3\cdot2\cdot1}$ is:
\begin{quote}\begin{multicols}{2}
\noindent (A) $\dfrac{(n-1)!}{3!2!}$ \\
(B) $\dfrac{(n-1)!}{3!(n-2)!}$ \\
(C) $\dfrac{n!}{3!(n-2)!}$ \\
(D) None of these
\end{multicols}\end{quote}

\item The expression $\dfrac{(n-1)!(n-2)!}{(n!)^2}$ reduces to:
\begin{quote}\begin{multicols}{2}
\noindent (A) $\dfrac{1}{n(n-1)}$ \\
(B) $\dfrac{1}{n^3-n^2}$ \\
(C) $\dfrac{1}{n^2}$ \\
(D) $\dfrac{1}{n}$
\end{multicols}\end{quote}

\item Factorial notation was introduced by:
\begin{quote}\begin{multicols}{2}
\noindent (A) Leibniz \\
(B) Euler \\
(C) Seikkowa \\
(D) Christian Kramp
\end{multicols}\end{quote}

\item $n$ different objects taken all at a time can be arranged in:
\begin{quote}\begin{multicols}{2}
\noindent (A) $n-1$ ways \\
(B) $n+1$ ways \\
(C) $n!$ ways \\
(D) $2n!$ ways
\end{multicols}\end{quote}

\item The value of $^6P_2$ is:
\begin{quote}\begin{multicols}{2}
\noindent (A) 5 \\
(B) 10 \\
(C) 15 \\
(D) 20
\end{multicols}\end{quote}

\item $^nP_n=$:
\begin{quote}\begin{multicols}{2}
\noindent (A) $n$ \\
(B) 0 \\
(C) 1 \\
(D) $n!$
\end{multicols}\end{quote}

\item If $r=1$, then $^nP_r=$:
\begin{quote}\begin{multicols}{2}
\noindent (A) $n!$ \\
(B) $n$ \\
(C) $n-1$ \\
(D) None of these
\end{multicols}\end{quote}

\item The value of $^6P_1$ is:
\begin{quote}\begin{multicols}{2}
\noindent (A) 18 \\
(B) 12 \\
(C) 6 \\
(D) 0
\end{multicols}\end{quote}

\item If $r=n$, then $^nP_r$ equals:
\begin{quote}\begin{multicols}{2}
\noindent (A) 1 \\
(B) $n!$ \\
(C) $(n+1)!$ \\
(D) $(n-1)!$
\end{multicols}\end{quote}

\item The value of $n$ when $^{11}P_n=11\cdot10\cdot9$ is:
\begin{quote}\begin{multicols}{2}
\noindent (A) 1 \\
(B) 3 \\
(C) 2 \\
(D) 0
\end{multicols}\end{quote}

\item If $^nP_2=30$ then $n=$:
\begin{quote}\begin{multicols}{2}
\noindent (A) 6 \\
(B) 4 \\
(C) 5 \\
(D) 8
\end{multicols}\end{quote}

\item $\dfrac{^nP_r}{r!}$ is equal to:
\begin{quote}\begin{multicols}{2}
\noindent (A) $^nC_r$ \\
(B) $^nC_{r-1}$ \\
(C) $^{n+1}C_r$ \\
(D) None of these
\end{multicols}\end{quote}

\item $^6P_4$ is equal to:
\begin{quote}\begin{multicols}{2}
\noindent (A) 160 \\
(B) 260 \\
(C) 360 \\
(D) 340
\end{multicols}\end{quote}

\item The number of words that can be formed from the letters of the word PLANE using all letters at a time is equal to:
\begin{quote}\begin{multicols}{2}
\noindent (A) 5 \\
(B) $4!$ \\
(C) $3!$ \\
(D) 4
\end{multicols}\end{quote}

\item Circular permutation of $n$ non-living things is given by:
\begin{quote}\begin{multicols}{2}
\noindent (A) $n!$ \\
(B) $\dfrac{n!}{2}$ \\
(C) $(n-1)!$ \\
(D) $\dfrac{(n-1)!}{2}$
\end{multicols}\end{quote}

\item The number of ways of arranging 5 keys in a circular ring is:
\begin{quote}\begin{multicols}{2}
\noindent (A) 24 \\
(B) 12 \\
(C) 6 \\
(D) 5
\end{multicols}\end{quote}

\item 5 persons can be seated at a round table in how many ways:
\begin{quote}\begin{multicols}{2}
\noindent (A) 23 \\
(B) 24 \\
(C) 25 \\
(D) 26
\end{multicols}\end{quote}

\item The number of permutations of the word PANAMA in which P must be the first letter:
\begin{quote}\begin{multicols}{2}
\noindent (A) 10 \\
(B) 60 \\
(C) 20 \\
(D) 120
\end{multicols}\end{quote}

\item Number of ways in which a necklace of 6 beads of different colors can be made:
\begin{quote}\begin{multicols}{2}
\noindent (A) $5!$ \\
(B) $\dfrac{5!}{2}$ \\
(C) $6!$ \\
(D) $\dfrac{6!}{2}$
\end{multicols}\end{quote}

\item $^nC_r$ is valid only if:
\begin{quote}\begin{multicols}{2}
\noindent (A) $n\leq r$ \\
(B) $r\leq n$ \\
(C) $r>n$ \\
(D) $r<0$
\end{multicols}\end{quote}

\item $^nC_r\times r!=$:
\begin{quote}\begin{multicols}{2}
\noindent (A) $^{n+1}P_r$ \\
(B) $^nP_{r+1}$ \\
(C) $^{n-1}P_r$ \\
(D) $^nP_r$
\end{multicols}\end{quote}

\item $^nC_{n-r}$ is equal to:
\begin{quote}\begin{multicols}{2}
\noindent (A) $^rP_n$ \\
(B) $^rC_n$ \\
(C) $^nP_r$ \\
(D) $^nC_r$
\end{multicols}\end{quote}

\item $^nC_r=\,^nC_{n-r}$ is useful when:
\begin{quote}\begin{multicols}{2}
\noindent (A) $n=r$ \\
(B) $n\neq r$ \\
(C) $r<\dfrac{n}{2}$ \\
(D) $r>\dfrac{n}{2}$
\end{multicols}\end{quote}

\item If $r=n$, $^nC_r$ is equal to:
\begin{quote}\begin{multicols}{2}
\noindent (A) 0 \\
(B) 1 \\
(C) $n$ \\
(D) $n!$
\end{multicols}\end{quote}

\item $^nC_r+\,^nC_{r-1}=$:
\begin{quote}\begin{multicols}{2}
\noindent (A) $^nC_r$ \\
(B) $^{n+1}C_r$ \\
(C) $^{n+1}C_{r+1}$ \\
(D) None
\end{multicols}\end{quote}

\item $^nC_0$ is equal to:
\begin{quote}\begin{multicols}{2}
\noindent (A) 1 \\
(B) 0 \\
(C) $n$ \\
(D) 2
\end{multicols}\end{quote}

\item $\dfrac{n!}{r!(n-r)!}$ is equal to:
\begin{quote}\begin{multicols}{2}
\noindent (A) $^nC_r$ \\
(B) $^nP_r$ \\
(C) $^{n+1}C_r$ \\
(D) $^{n-1}C_r$
\end{multicols}\end{quote}

\item If $^nC_5=\,^nC_{12}$ then $n$ is equal to:
\begin{quote}\begin{multicols}{2}
\noindent (A) 4 \\
(B) 8 \\
(C) 5 \\
(D) 9
\end{multicols}\end{quote}

\item The value of $^nC_{10}\times\dfrac{12\times11}{2!}$:
\begin{quote}\begin{multicols}{2}
\noindent (A) 12 \\
(B) 8 \\
(C) 11 \\
(D) 13
\end{multicols}\end{quote}

\item In how many ways can a cricket team of 11 players out of 15 be selected if the captain must be included every time:
\begin{quote}\begin{multicols}{2}
\noindent (A) $15!$ \\
(B) $^{14}C_{10}$ \\
(C) $^{15}C_{11}$ \\
(D) $11!$
\end{multicols}\end{quote}

\item The method of induction was devised by:
\begin{quote}\begin{multicols}{2}
\noindent (A) Laplace \\
(B) Francesco Maurolico \\
(C) Aristotle \\
(D) Euclid
\end{multicols}\end{quote}

\item For $n=2$, $\dfrac{n^3+2n}{3}$ represents:
\begin{quote}\begin{multicols}{2}
\noindent (A) Irrational number \\
(B) Negative integer \\
(C) Odd numbers \\
(D) Natural numbers
\end{multicols}\end{quote}

\item If $n^2>n+3$ then it is true for:
\begin{quote}\begin{multicols}{2}
\noindent (A) $n=0$ \\
(B) $n<1$ \\
(C) $n\geq2$ \\
(D) $n\geq3$
\end{multicols}\end{quote}

\item The inequality $n!>2^n+1$ is valid if:
\begin{quote}\begin{multicols}{2}
\noindent (A) $n<4$ \\
(B) $n\geq4$ \\
(C) $n\leq3$ \\
(D) $n=3$
\end{multicols}\end{quote}

\item For what value is $3^n>n!$ untrue if $n\in Z$:
\begin{quote}\begin{multicols}{2}
\noindent (A) $n=6$ \\
(B) $n=7$ \\
(C) $n=2$ \\
(D) $n=3$
\end{multicols}\end{quote}

\item The inequality $4^n>3^n+4$ is true for:
\begin{quote}\begin{multicols}{2}
\noindent (A) $n=1$ \\
(B) $n\geq2$ \\
(C) $n=0$ \\
(D) None
\end{multicols}\end{quote}

\item The inequality $n!>2^{n-1}$ is true for integral values of $n$:
\begin{quote}\begin{multicols}{2}
\noindent (A) $n\leq4$ \\
(B) $n\leq1$ \\
(C) $n\geq2$ \\
(D) $n\geq4$
\end{multicols}\end{quote}

\item If $n$ is a positive integer, then $n!>3^{n-1}$ is true when:
\begin{quote}\begin{multicols}{2}
\noindent (A) $n\leq3$ \\
(B) $n\geq2$ \\
(C) $n\leq5$ \\
(D) None
\end{multicols}\end{quote}

\item The inequality $4^n>3^n+2^{n-1}$ is true for integral values of $n$ if:
\begin{quote}\begin{multicols}{2}
\noindent (A) $n>0$ \\
(B) $n\geq1$ \\
(C) $n\geq2$ \\
(D) $n\leq2$
\end{multicols}\end{quote}

\item If $n$ is any positive integer, then $2^n>2(n+1)$ is true for all:
\begin{quote}\begin{multicols}{2}
\noindent (A) $n\leq3$ \\
(B) $n<3$ \\
(C) $n\geq3$ \\
(D) $n>3$
\end{multicols}\end{quote}

\item Show that $n!>n^2$ is true for integral value of $n$:
\begin{quote}\begin{multicols}{2}
\noindent (A) $n=3$ \\
(B) $n<4$ \\
(C) $n\geq4$ \\
(D) $n\leq4$
\end{multicols}\end{quote}

\item Which one is divisible by 2 for all positive integral values of $n$:
\begin{quote}\begin{multicols}{2}
\noindent (A) $n^3-n$ \\
(B) $5^n-1$ \\
(C) $5^n-2^2$ \\
(D) $n^2+n$
\end{multicols}\end{quote}

\item For positive integral values of $n$, $5^n-2^n$ is divisible by:
\begin{quote}\begin{multicols}{2}
\noindent (A) 2 \\
(B) 3 \\
(C) 4 \\
(D) 5
\end{multicols}\end{quote}

\item If $n$ is positive then $n^2+n$ is divisible by:
\begin{quote}\begin{multicols}{2}
\noindent (A) 2 \\
(B) 3 \\
(C) 4 \\
(D) 5
\end{multicols}\end{quote}

\item If $n$ is any positive integer then $1+3+5+\ldots+(2n+1)=(n+1)^2$ is true for:
\begin{quote}\begin{multicols}{2}
\noindent (A) $n>-1$ \\
(B) $n\geq-2$ \\
(C) $n>1$ \\
(D) $n>2$
\end{multicols}\end{quote}

\item $\dbinom{3}{3}+\dbinom{4}{3}+\dbinom{5}{3}+\ldots+\dbinom{n+2}{3}$ equal to:
\begin{quote}\begin{multicols}{2}
\noindent (A) $n+3$ \\
(B) $\dbinom{n+4}{2}$ \\
(C) $\dbinom{n+3}{2}$ \\
(D) $n+4$
\end{multicols}\end{quote}

\item $1^2+2^2+3^2+\ldots n^2=$:
\begin{quote}\begin{multicols}{2}
\noindent (A) $\dfrac{n(n+1)}{2}$ \\
(B) $\left[\dfrac{n(n+1)}{2}\right]^2$ \\
(C) $\dfrac{n(n+1)(2n+1)}{6}$ \\
(D) $\dfrac{n(n+1)(n+2)}{6}$
\end{multicols}\end{quote}

\item An algebraic expression consisting of two terms is called:
\begin{quote}\begin{multicols}{2}
\noindent (A) Monomial \\
(B) Trinomial \\
(C) Polynomial \\
(D) Binomial
\end{multicols}\end{quote}

\item In the binomial expression of $(a+x)^n$, $n$ is called:
\begin{quote}\begin{multicols}{2}
\noindent (A) Root \\
(B) Element \\
(C) Index \\
(D) None of these
\end{multicols}\end{quote}

\item Number of terms in the expression $(1+x)^{2n+1}$ is:
\begin{quote}\begin{multicols}{2}
\noindent (A) $2n+1$ \\
(B) $2n$ \\
(C) $2n+2$ \\
(D) $3n+1$
\end{multicols}\end{quote}

\item The number of terms in the expansion of $(a+b)^n$ is:
\begin{quote}\begin{multicols}{2}
\noindent (A) $n^2+1$ \\
(B) $n+1$ \\
(C) $n-1$ \\
(D) $n$
\end{multicols}\end{quote}

\item Coefficient of the 8th term in the expansion of $(a+b)^{10}$ is:
\begin{quote}\begin{multicols}{2}
\noindent (A) $^{10}C_7$ \\
(B) $^{10}C_8$ \\
(C) $^8C_{10}$ \\
(D) None of these
\end{multicols}\end{quote}

\item In the expansion of $(x-y)^n$ the terms are alternatively positive and:
\begin{quote}\begin{multicols}{2}
\noindent (A) Negative \\
(B) Undefined \\
(C) Absurd \\
(D) None of these
\end{multicols}\end{quote}

\item If the number of terms in the expansion of $(a+b)^n$ is 16 then $n$ equals:
\begin{quote}\begin{multicols}{2}
\noindent (A) 18 \\
(B) 16 \\
(C) 17 \\
(D) 15
\end{multicols}\end{quote}

\item The number of terms in the expansion of $(x-3)^{10}=$:
\begin{quote}\begin{multicols}{2}
\noindent (A) 10 \\
(B) 11 \\
(C) 12 \\
(D) 13
\end{multicols}\end{quote}

\item The number of terms in the expansion of $(2a+b)^{13}$ is:
\begin{quote}\begin{multicols}{2}
\noindent (A) 13 \\
(B) 14 \\
(C) 15 \\
(D) 12
\end{multicols}\end{quote}

\item The number of terms in the expansion of $(a+x)^{10}$ is:
\begin{quote}\begin{multicols}{2}
\noindent (A) 10 \\
(B) 11 \\
(C) 12 \\
(D) 9
\end{multicols}\end{quote}

\item The number of terms in the expansion of $(a+x)^n$:
\begin{quote}\begin{multicols}{2}
\noindent (A) $n+1$ \\
(B) $n-1$ \\
(C) $n$ \\
(D) $2n$
\end{multicols}\end{quote}

\item The number of terms in the expansion of $(x^2-1)^7$ is:
\begin{quote}\begin{multicols}{2}
\noindent (A) 2 \\
(B) 7 \\
(C) 8 \\
(D) 12
\end{multicols}\end{quote}

\item In the expansion of $(a+b)^7$, the 2nd term is:
\begin{quote}\begin{multicols}{2}
\noindent (A) $a^7$ \\
(B) $7a^6b$ \\
(C) $7ab$ \\
(D) 0
\end{multicols}\end{quote}

\item In the expansion of $(3+x)^4$ the middle term will be:
\begin{quote}\begin{multicols}{2}
\noindent (A) 81 \\
(B) $54x^2$ \\
(C) $26x^2$ \\
(D) $x^4$
\end{multicols}\end{quote}

\item The middle term in the expansion of $(x-y)^{12}$ is:
\begin{quote}\begin{multicols}{2}
\noindent (A) 6th \\
(B) 5th \\
(C) 7th \\
(D) 8th
\end{multicols}\end{quote}

\item Middle term(s) in the expansion of $(x+y)^{21}$:
\begin{quote}\begin{multicols}{2}
\noindent (A) 10th, 11th \\
(B) 9th, 10th \\
(C) 11th, 12th \\
(D) 12th, 13th
\end{multicols}\end{quote}

\item If $n$ is even, then the expansion of $(a+x)^n$ has:
\begin{quote}\begin{multicols}{2}
\noindent (A) Two middle terms \\
(B) No middle term \\
(C) One middle term \\
(D) None of these
\end{multicols}\end{quote}

\item If $n$ is odd, then the expansion of $(a+x)^n$ has:
\begin{quote}\begin{multicols}{2}
\noindent (A) Two middle terms \\
(B) 3 middle terms \\
(C) 4 middle terms \\
(D) 0 middle term
\end{multicols}\end{quote}

\item Using the binomial theorem, $(2.02)^4$ approximated to two decimal places is:
\begin{quote}\begin{multicols}{2}
\noindent (A) 16.46 \\
(B) 16.44 \\
(C) 16.40 \\
(D) 16.60
\end{multicols}\end{quote}

\item The coefficient of the last term in the expansion of $(x-y)^5$ is:
\begin{quote}\begin{multicols}{2}
\noindent (A) $-1$ \\
(B) 1 \\
(C) 5 \\
(D) $-5$
\end{multicols}\end{quote}

\item The sum of the exponents of $a$ and $b$ in every term in the expansion of $(a+b)^n$ is:
\begin{quote}\begin{multicols}{2}
\noindent (A) 1 \\
(B) $n$ \\
(C) $n+1$ \\
(D) $n-1$
\end{multicols}\end{quote}

\item The sum of the odd coefficients in the expansion of $(1+x)^5$ is:
\begin{quote}\begin{multicols}{2}
\noindent (A) 5 \\
(B) 16 \\
(C) 25 \\
(D) 32
\end{multicols}\end{quote}

\item The sum of the even coefficients in the expansion of $(1+x)^5$ is:
\begin{quote}\begin{multicols}{2}
\noindent (A) 32 \\
(B) 64 \\
(C) 16 \\
(D) 8
\end{multicols}\end{quote}

\item The sum of the binomial coefficients in the expansion of $(a+x)^5$ is:
\begin{quote}\begin{multicols}{2}
\noindent (A) 8 \\
(B) 10 \\
(C) 16 \\
(D) 32
\end{multicols}\end{quote}

\item Sum of binomial coefficients is:
\begin{quote}\begin{multicols}{2}
\noindent (A) $n$ \\
(B) $2^n$ \\
(C) $2n$ \\
(D) $n^2$
\end{multicols}\end{quote}

\item The sum of odd coefficients in the binomial expansion of $(a+x)^n$ is equal to:
\begin{quote}\begin{multicols}{2}
\noindent (A) $2^{n-1}$ \\
(B) $2^{n+1}$ \\
(C) $2^n$ \\
(D) $2^n-1$
\end{multicols}\end{quote}

\item The sum of coefficients in the binomial expansion when $n=4$ is:
\begin{quote}\begin{multicols}{2}
\noindent (A) 1 \\
(B) 8 \\
(C) 16 \\
(D) 32
\end{multicols}\end{quote}

\item The sum of the even coefficients in the expansion of $(1+x)^n$ is equal to:
\begin{quote}\begin{multicols}{2}
\noindent (A) $n^2$ \\
(B) $2^{n-2}$ \\
(C) $2^{n-1}$ \\
(D) $2^n$
\end{multicols}\end{quote}

\item The sum of the odd coefficients in the expansion of $(a+b)^4$ is:
\begin{quote}\begin{multicols}{2}
\noindent (A) 4 \\
(B) 6 \\
(C) 8 \\
(D) 16
\end{multicols}\end{quote}

\item The sum of the odd coefficients in the expansion of $(1+x)^4$ is:
\begin{quote}\begin{multicols}{2}
\noindent (A) 8 \\
(B) 1 \\
(C) 12 \\
(D) 16
\end{multicols}\end{quote}

\item The sum of the even coefficients in the expansion of $(1+x)^4$ is:
\begin{quote}\begin{multicols}{2}
\noindent (A) 16 \\
(B) 12 \\
(C) 10 \\
(D) 8
\end{multicols}\end{quote}

\item The sum of all coefficients in the binomial expansion of $(1+x)^n$ is:
\begin{quote}\begin{multicols}{2}
\noindent (A) $2^{n-1}$ \\
(B) $2n+1$ \\
(C) $2^n$ \\
(D) 0
\end{multicols}\end{quote}

\item The number of terms in the expansion of $(2a+b)^{13}$ is:
\begin{quote}\begin{multicols}{2}
\noindent (A) 13 \\
(B) 14 \\
(C) 15 \\
(D) 12
\end{multicols}\end{quote}

\item The number of terms in the expansion of $(1+x)^{1/2}$, $|x|<1$, is:
\begin{quote}\begin{multicols}{2}
\noindent (A) 2 \\
(B) $n$ \\
(C) $\dfrac{n}{2}$ \\
(D) Infinite
\end{multicols}\end{quote}

\item Number of terms in the expansion of $(1+x)^{1/3}$ is:
\begin{quote}\begin{multicols}{2}
\noindent (A) 6 \\
(B) 7 \\
(C) $1/3$ \\
(D) Infinite
\end{multicols}\end{quote}

\item $1-x+x^2-x^3+\ldots$ is the expansion of:
\begin{quote}\begin{multicols}{2}
\noindent (A) $(1+x)^{-2}$ \\
(B) $(1-x)^{-2}$ \\
(C) $(1+x)^{-1}$ \\
(D) $(1-x)^{-1}$
\end{multicols}\end{quote}

\item The 2nd term in the expansion of $(1-x)^{-1/2}$ is:
\begin{quote}\begin{multicols}{2}
\noindent (A) 2 \\
(B) $\dfrac{2}{\sqrt{15}}$ \\
(C) $\dfrac{1}{2}$ \\
(D) $4x$
\end{multicols}\end{quote}

\item The 2nd term in the expansion of $(1+2x)^{-1/3}$ is:
\begin{quote}\begin{multicols}{2}
\noindent (A) $\dfrac{-2}{3}x$ \\
(B) $\dfrac{2}{3}x$ \\
(C) $-6x$ \\
(D) $\dfrac{x}{3}$
\end{multicols}\end{quote}

\item 2nd term in $(1-x)^{-1}$ is:
\begin{quote}\begin{multicols}{2}
\noindent (A) 1 \\
(B) $x$ \\
(C) $2x$ \\
(D) $3x$
\end{multicols}\end{quote}

\item In the expansion of $(1+x)^{-3}$ the 4th term is:
\begin{quote}\begin{multicols}{2}
\noindent (A) $-3x$ \\
(B) $-10x^3$ \\
(C) $6x^2$ \\
(D) $10x^3$
\end{multicols}\end{quote}

\item The second term in the expansion of $\left(1+\dfrac{1}{2}x\right)^{-1}$ is:
\begin{quote}\begin{multicols}{2}
\noindent (A) $\dfrac{-x}{3}$ \\
(B) $2x$ \\
(C) $\dfrac{-x}{2}$ \\
(D) $\dfrac{x}{2}$
\end{multicols}\end{quote}

\item The second term in the expansion of $\left(1-\dfrac{1}{3}x\right)^{-1}$ is:
\begin{quote}\begin{multicols}{2}
\noindent (A) $\dfrac{1}{3}x$ \\
(B) $-\dfrac{1}{3}x$ \\
(C) $3x$ \\
(D) $2x$
\end{multicols}\end{quote}

\item The second term in $(1-x)^{1/3}$ is:
\begin{quote}\begin{multicols}{2}
\noindent (A) $\dfrac{x}{3}$ \\
(B) $\dfrac{3}{x}$ \\
(C) $-\dfrac{x}{3}$ \\
(D) $-\dfrac{3}{x}$
\end{multicols}\end{quote}

\item The second term in the expansion of $(1-2x)^{-1/2}$ is:
\begin{quote}\begin{multicols}{2}
\noindent (A) $x$ \\
(B) $2x$ \\
(C) $3x$ \\
(D) $4x$
\end{multicols}\end{quote}

\item The second term in the expansion of $(1-2x)^{1/3}$ is:
\begin{quote}\begin{multicols}{2}
\noindent (A) $-\dfrac{2x}{3}$ \\
(B) $\dfrac{2x}{3}$ \\
(C) $\dfrac{x}{2}$ \\
(D) $\dfrac{x}{3}$
\end{multicols}\end{quote}

\item The second term in the expansion of $\left(1+\dfrac{2x}{3}\right)^{3/2}$ is:
\begin{quote}\begin{multicols}{2}
\noindent (A) $x$ \\
(B) $\dfrac{2}{3}x$ \\
(C) $\dfrac{9}{4}x$ \\
(D) $\dfrac{4}{9}x$
\end{multicols}\end{quote}

\item If $|x|<1$ then the expansion of $(1-x)^{1/2}$ up to two terms is:
\begin{quote}\begin{multicols}{2}
\noindent (A) $1+\dfrac{1}{2}x$ \\
(B) $1-\dfrac{1}{2}x$ \\
(C) $1-x$ \\
(D) $1+x$
\end{multicols}\end{quote}

\item The expansion of $(1-3x)^{2/3}$ is valid if:
\begin{quote}\begin{multicols}{2}
\noindent (A) $|x|<1$ \\
(B) $|x|<1/3$ \\
(C) $|x|>1$ \\
(D) None of these
\end{multicols}\end{quote}

\item The expansion of $(1-5x)^{1/2}$ is valid if:
\begin{quote}\begin{multicols}{2}
\noindent (A) $|x|<1/2$ \\
(B) $|x|<1/5$ \\
(C) $|x|<1$ \\
(D) None of these
\end{multicols}\end{quote}

\item The expansion of $(1+x)^{-1/4}$ is valid only if:
\begin{quote}\begin{multicols}{2}
\noindent (A) $|x|>1$ \\
(B) $|x|<1$ \\
(C) $|x|<-1$ \\
(D) $|x|>-1$
\end{multicols}\end{quote}

\item The binomial coefficients in Pascal's triangle, in the $n=4$ row, what is the third number:
\begin{quote}\begin{multicols}{2}
\noindent (A) $1$ \\
(B) $4$ \\
(C) $6$ \\
(D) $2$
\end{multicols}\end{quote}

\item Using the binomial expansion, what is the approximate value of $\sqrt{99}$:
\begin{quote}\begin{multicols}{2}
\noindent (A) $9.19$ \\
(B) $9.29$ \\
(C) $9.49$ \\
(D) $9.89$
\end{multicols}\end{quote}

\item $1-x+x^2-x^3+\cdots$ is the expansion of:
\begin{quote}\begin{multicols}{2}
\noindent (A) $(1-x)^{-1}$ \\
(B) $(1+x)^{-1}$ \\
(C) $(1+2x)^{-1/2}$ \\
(D) $(1-2x)^{-1/2}$
\end{multicols}\end{quote}

\item The leading coefficient of the equation $x^3+2x^2+3x+4=0$ is:
\begin{quote}\begin{multicols}{2}
\noindent (A) $1$ \\
(B) $2$ \\
(C) $3$ \\
(D) $4$
\end{multicols}\end{quote}

\item If a polynomial $P(x)=x^3+4x^2-2x+5$ is divided by $x-1$, then the remainder is:
\begin{quote}\begin{multicols}{2}
\noindent (A) $4$ \\
(B) $-2$ \\
(C) $5$ \\
(D) $8$
\end{multicols}\end{quote}

\item Dividend $=$ (divisor)(quotient) $+$ \underline{\hspace{1.5cm}}:
\begin{quote}\begin{multicols}{2}
\noindent (A) Zero \\
(B) Remainder \\
(C) Divisor \\
(D) None of these
\end{multicols}\end{quote}

\item $x^3-3x^2+2x-6$ has factor:
\begin{quote}\begin{multicols}{2}
\noindent (A) $x-4$ \\
(B) $x-3$ \\
(C) $x+3$ \\
(D) $x+2$
\end{multicols}\end{quote}

\item When $3x^4+4x^3+x-5$ is divided by $x+1$, the remainder is:
\begin{quote}\begin{multicols}{2}
\noindent (A) $7$ \\
(B) $6$ \\
(C) $-6$ \\
(D) $-7$
\end{multicols}\end{quote}

\item If $(x-2)$ is a factor of $ax^2-12x+4$, then $a=$:
\begin{quote}\begin{multicols}{2}
\noindent (A) $2$ \\
(B) $5$ \\
(C) $7$ \\
(D) $9$
\end{multicols}\end{quote}

\item If one root of the equation $x^2+ax+2=0$ is $2$, then $a$ is equal to:
\begin{quote}\begin{multicols}{2}
\noindent (A) $-3$ \\
(B) $3$ \\
(C) $4$ \\
(D) $-2$
\end{multicols}\end{quote}

\item If $x-2$ is a factor of the polynomial $x^3+2x^2+kx+4$, then $k$ equals:
\begin{quote}\begin{multicols}{2}
\noindent (A) $10$ \\
(B) $-10$ \\
(C) $2$ \\
(D) $4$
\end{multicols}\end{quote}

\item Synthetic division is a process of:
\begin{quote}\begin{multicols}{2}
\noindent (A) Addition \\
(B) Multiplication \\
(C) Subtraction \\
(D) Division
\end{multicols}\end{quote}

\item A factor of $f(x)=x^n-a^n$, where $n$ is a positive integer, is:
\begin{quote}\begin{multicols}{2}
\noindent (A) $x+a$ \\
(B) $x-a$ \\
(C) $x-1$ \\
(D) $x-2$
\end{multicols}\end{quote}

\item If $x-a$ is a factor of $f(x)$, then $f(a)=0$; $x=a$ is called a:
\begin{quote}\begin{multicols}{2}
\noindent (A) Root \\
(B) Factor \\
(C) Polynomial \\
(D) Degree
\end{multicols}\end{quote}

\item Polynomial regression is used when the data shows a \underline{\hspace{1.5cm}} relationship:
\begin{quote}\begin{multicols}{2}
\noindent (A) Linear \\
(B) Curved \\
(C) Constant \\
(D) Absolute
\end{multicols}\end{quote}

\item Observed $-$ Predicted $=$:
\begin{quote}\begin{multicols}{2}
\noindent (A) Profit \\
(B) Loss \\
(C) Error \\
(D) Absolute
\end{multicols}\end{quote}

\item Angles with a basic angle of measure $\theta$, added to a right angle or its multiple, are called \underline{\hspace{1.5cm}} angles:
\begin{quote}\begin{multicols}{2}
\noindent (A) Supplementary \\
(B) Complementary \\
(C) Obtuse \\
(D) Allied
\end{multicols}\end{quote}

\item The angles $90^{\circ}\pm\theta$, $180^{\circ}\pm\theta$, $270^{\circ}\pm\theta$, $360^{\circ}\pm\theta$ are called \underline{\hspace{1.5cm}} angles:
\begin{quote}\begin{multicols}{2}
\noindent (A) Composite \\
(B) Half \\
(C) Quadrantal \\
(D) Allied
\end{multicols}\end{quote}

\item The distance between $A(3,8)$, $B(5,6)$ is:
\begin{quote}\begin{multicols}{2}
\noindent (A) $\sqrt{2}$ \\
(B) $2\sqrt{2}$ \\
(C) $2$ \\
(D) None of these
\end{multicols}\end{quote}

\item $\sin(-\alpha) = $:
\begin{quote}\begin{multicols}{2}
\noindent (A) $\sec\alpha$ \\
(B) $-\sin\alpha$ \\
(C) $\sin\alpha$ \\
(D) $-\cos\alpha$
\end{multicols}\end{quote}

\item $\dfrac{3\pi}{2}+\theta$ (for small $\theta>0$) lies in:
\begin{quote}\begin{multicols}{2}
\noindent (A) 1st quadrant \\
(B) 2nd quadrant \\
(C) 3rd quadrant \\
(D) 4th quadrant
\end{multicols}\end{quote}

\item $\sin390^{\circ}$ is equal to:
\begin{quote}\begin{multicols}{2}
\noindent (A) $\sin30^{\circ}$ \\
(B) $\cos30^{\circ}$ \\
(C) Zero \\
(D) $\sin60^{\circ}$
\end{multicols}\end{quote}

\item $\cos315^{\circ} = $:
\begin{quote}\begin{multicols}{2}
\noindent (A) $0$ \\
(B) $1$ \\
(C) $\dfrac{\sqrt{3}}{2}$ \\
(D) $\dfrac{1}{\sqrt{2}}$
\end{multicols}\end{quote}

\item $\sin196^{\circ} = $:
\begin{quote}\begin{multicols}{2}
\noindent (A) $\sin16^{\circ}$ \\
(B) $\cos33^{\circ}$ \\
(C) $\cos16^{\circ}$ \\
(D) $-\sin16^{\circ}$
\end{multicols}\end{quote}

\item $\cos(-60^{\circ}) = $:
\begin{quote}\begin{multicols}{2}
\noindent (A) $\dfrac{1}{2}$ \\
(B) $\dfrac{\sqrt{3}}{2}$ \\
(C) $\dfrac{1}{\sqrt{2}}$ \\
(D) None of these
\end{multicols}\end{quote}

\item $\sin(-300^{\circ}) = $:
\begin{quote}\begin{multicols}{2}
\noindent (A) $\dfrac{\sqrt{3}}{2}$ \\
(B) $0$ \\
(C) $2$ \\
(D) $3$
\end{multicols}\end{quote}

\item $\sin\left(\theta-\dfrac{\pi}{2}\right) = $:
\begin{quote}\begin{multicols}{2}
\noindent (A) $\cos\theta$ \\
(B) $-\cos\theta$ \\
(C) $\sin\theta$ \\
(D) $-\sin\theta$
\end{multicols}\end{quote}

\item $\cos\left(\dfrac{3\pi}{2}-\theta\right) = $:
\begin{quote}\begin{multicols}{2}
\noindent (A) $\cos\theta$ \\
(B) $-\cos\theta$ \\
(C) $\sin\theta$ \\
(D) $-\sin\theta$
\end{multicols}\end{quote}

\item $\cot\left(\dfrac{\pi}{2}-\beta\right) = $:
\begin{quote}\begin{multicols}{2}
\noindent (A) $\sin\beta$ \\
(B) $\cot\beta$ \\
(C) $-\cot\beta$ \\
(D) $\tan\beta$
\end{multicols}\end{quote}

\item $\tan(\pi-\alpha) = $:
\begin{quote}\begin{multicols}{2}
\noindent (A) $\tan\alpha$ \\
(B) $-\tan\theta$ \\
(C) $-\tan\alpha$ \\
(D) $\cot\alpha$
\end{multicols}\end{quote}

\item $\tan\left(\dfrac{3\pi}{2}-\theta\right)$ is equal to:
\begin{quote}\begin{multicols}{2}
\noindent (A) $\tan\theta$ \\
(B) $-\cot\theta$ \\
(C) $\cot\theta$ \\
(D) $-\tan\theta$
\end{multicols}\end{quote}

\item $\sin\left(\dfrac{\pi}{2}+\theta\right)$ is equal to:
\begin{quote}\begin{multicols}{2}
\noindent (A) $-\cos\theta$ \\
(B) $\cos\theta$ \\
(C) $\sin\theta$ \\
(D) $-\sin\theta$
\end{multicols}\end{quote}

\item $\sin\left(\dfrac{3\pi}{2}-\theta\right) = $:
\begin{quote}\begin{multicols}{2}
\noindent (A) $\cos\theta$ \\
(B) $\sin\theta$ \\
(C) $-\cos\theta$ \\
(D) $-\sin\theta$
\end{multicols}\end{quote}

\item $\sin(2\pi-\theta) = $:
\begin{quote}\begin{multicols}{2}
\noindent (A) $\sin\theta$ \\
(B) $-\sin\theta$ \\
(C) $\cos\theta$ \\
(D) $\tan\theta$
\end{multicols}\end{quote}

\item $\tan(\pi+\theta) = $:
\begin{quote}\begin{multicols}{2}
\noindent (A) $\cot\theta$ \\
(B) $-\cot\theta$ \\
(C) $\dfrac{1}{\cot\theta}$ \\
(D) None
\end{multicols}\end{quote}

\item If $\alpha,\beta,\gamma$ are the angles of a triangle, then $\tan(\alpha+\beta)+\tan\gamma = $:
\begin{quote}\begin{multicols}{2}
\noindent (A) $0$ \\
(B) $1$ \\
(C) $2$ \\
(D) None of these
\end{multicols}\end{quote}

\item If $\alpha,\beta,\gamma$ are the angles of triangle $ABC$, then $\cos\left(\dfrac{\alpha+\beta}{2}\right)$ will be equal to:
\begin{quote}\begin{multicols}{2}
\noindent (A) $\sin\alpha$ \\
(B) $\sin\gamma$ \\
(C) $\sin\dfrac{\gamma}{2}$ \\
(D) $\sin\beta$
\end{multicols}\end{quote}

\item $\cos(\theta-90^{\circ})-\cos(\theta+90^{\circ}) = $:
\begin{quote}\begin{multicols}{2}
\noindent (A) $-2\cos\theta$ \\
(B) $2\cos\theta$ \\
(C) $2\sin\theta$ \\
(D) $-2\sin\theta$
\end{multicols}\end{quote}

\item The value of $\cos75^{\circ}$ is:
\begin{quote}\begin{multicols}{2}
\noindent (A) $\dfrac{\sqrt{3}-1}{2\sqrt{2}}$ \\
(B) $\dfrac{-\sqrt{3}+1}{2\sqrt{2}}$ \\
(C) $\dfrac{\sqrt{3}+1}{2\sqrt{2}}$ \\
(D) $\dfrac{-\sqrt{3}-1}{2\sqrt{2}}$
\end{multicols}\end{quote}

\item $\sin\left(\theta+\dfrac{\pi}{6}\right)+\cos\left(\theta+\dfrac{\pi}{3}\right) = $:
\begin{quote}\begin{multicols}{2}
\noindent (A) $\cos\theta$ \\
(B) $\sin\theta$ \\
(C) $\sec\theta$ \\
(D) $\csc\theta$
\end{multicols}\end{quote}

\item $\dfrac{\cos11^{\circ}+\sin11^{\circ}}{\cos11^{\circ}-\sin11^{\circ}} = $:
\begin{quote}\begin{multicols}{2}
\noindent (A) $\tan11^{\circ}$ \\
(B) $\cot11^{\circ}$ \\
(C) $\tan56^{\circ}$ \\
(D) $\cot56^{\circ}$
\end{multicols}\end{quote}

\item $\dfrac{\cos8^{\circ}-\sin8^{\circ}}{\cos8^{\circ}+\sin8^{\circ}} = $:
\begin{quote}\begin{multicols}{2}
\noindent (A) $\tan8^{\circ}$ \\
(B) $\cot8^{\circ}$ \\
(C) $\tan37^{\circ}$ \\
(D) $\cot37^{\circ}$
\end{multicols}\end{quote}

\item If $r\cos\theta=3$, $r\sin\theta=4$, then $r=$:
\begin{quote}\begin{multicols}{2}
\noindent (A) $25$ \\
(B) $-25$ \\
(C) $5$ \\
(D) $-5$
\end{multicols}\end{quote}

\item $\sin\theta$ equals:
\begin{quote}\begin{multicols}{2}
\noindent (A) $2\sin^2\dfrac{\theta}{2}$ \\
(B) $2\sin\dfrac{\theta}{2}\cos\dfrac{\theta}{2}$ \\
(C) $2\cos^2\dfrac{\theta}{2}$ \\
(D) $2\tan\dfrac{\theta}{2}$
\end{multicols}\end{quote}

\item $\cos^2\theta = $:
\begin{quote}\begin{multicols}{2}
\noindent (A) $\dfrac{1+\cos\theta}{2}$ \\
(B) $\dfrac{1+\cos2\theta}{2}$ \\
(C) $\dfrac{1-\cos2\theta}{2}$ \\
(D) $\dfrac{\cos2\theta-1}{2}$
\end{multicols}\end{quote}

\item $2\sin^2\dfrac{\theta}{2} = $:
\begin{quote}\begin{multicols}{2}
\noindent (A) $1+\cos\theta$ \\
(B) $1-\sin\theta$ \\
(C) $1+\sin\theta$ \\
(D) $1-\cos\theta$
\end{multicols}\end{quote}

\item If $\sin2\theta=1$, then the value of $\theta$ is:
\begin{quote}\begin{multicols}{2}
\noindent (A) $30^{\circ}$ \\
(B) $45^{\circ}$ \\
(C) $60^{\circ}$ \\
(D) $90^{\circ}$
\end{multicols}\end{quote}

\item $\dfrac{1-\tan^2x}{1+\tan^2x} = $:
\begin{quote}\begin{multicols}{2}
\noindent (A) $\sin2x$ \\
(B) $\cos2x$ \\
(C) $\tan2x$ \\
(D) $\cot2x$
\end{multicols}\end{quote}

\item $\tan\dfrac{\alpha}{2} = $:
\begin{quote}\begin{multicols}{2}
\noindent (A) $\pm\sqrt{\dfrac{1-\cos\alpha}{1-\tan^2\alpha}}$ \\
(B) $\pm\sqrt{\dfrac{1-\cos\alpha}{1+\cos\alpha}}$ \\
(C) $\pm\sqrt{\dfrac{1-\sin\alpha}{1-\cos\alpha}}$ \\
(D) $\pm\sqrt{\dfrac{1-\sin\alpha}{1-\sin\alpha}}$
\end{multicols}\end{quote}

\item $\tan2\theta = $:
\begin{quote}\begin{multicols}{2}
\noindent (A) $\dfrac{2\tan\theta}{1-\tan^2\theta}$ \\
(B) $\dfrac{\tan\theta}{1+\tan^2\theta}$ \\
(C) $\dfrac{2\tan^2\theta}{1-\tan^2\theta}$ \\
(D) $\dfrac{1-\tan^2\theta}{1+\tan^2\theta}$
\end{multicols}\end{quote}

\item $\cos3\alpha = $:
\begin{quote}\begin{multicols}{2}
\noindent (A) $4\sin^3\alpha-3\sin\alpha$ \\
(B) $2\cos\dfrac{3\alpha}{2}\sin\dfrac{3\alpha}{2}$ \\
(C) $\cos^2\dfrac{3\alpha}{2}-\sin^2\dfrac{3\alpha}{2}$ \\
(D) None of these
\end{multicols}\end{quote}

\item $\cos2\theta = $:
\begin{quote}\begin{multicols}{2}
\noindent (A) $\dfrac{2\tan\theta}{1-\tan^2\theta}$ \\
(B) $\dfrac{2\tan\theta}{1+\tan^2\theta}$ \\
(C) $\dfrac{1-\tan^2\theta}{1+\tan^2\theta}$ \\
(D) None of these
\end{multicols}\end{quote}

\item $(\cos2\theta)^2+(\sin2\theta)^2 = $:
\begin{quote}\begin{multicols}{2}
\noindent (A) $0$ \\
(B) $2$ \\
(C) $4$ \\
(D) $1$
\end{multicols}\end{quote}

\item $\dfrac{2\tan\theta}{1-\tan^2\theta}$ is equal to:
\begin{quote}\begin{multicols}{2}
\noindent (A) $\tan2\theta$ \\
(B) $\cot2\theta$ \\
(C) $\sin2\theta$ \\
(D) $\cos2\theta$
\end{multicols}\end{quote}

\item $\pm\sqrt{\dfrac{1-\cos\theta}{2}} = $:
\begin{quote}\begin{multicols}{2}
\noindent (A) $\sin\dfrac{\theta}{2}$ \\
(B) $\cos\dfrac{\theta}{2}$ \\
(C) $\sin\theta$ \\
(D) $\cos\theta$
\end{multicols}\end{quote}

\item On simplifying $\dfrac{\sin2\theta}{1+\cos2\theta}$, the result is:
\begin{quote}\begin{multicols}{2}
\noindent (A) $\sin\theta$ \\
(B) $\cot\theta$ \\
(C) $\tan\theta$ \\
(D) $\sec\theta$
\end{multicols}\end{quote}

\item $\dfrac{1+\cos\theta}{\sin\theta} = $:
\begin{quote}\begin{multicols}{2}
\noindent (A) $\sin\theta$ \\
(B) $\cot\dfrac{\theta}{2}$ \\
(C) $\csc2\theta$ \\
(D) None
\end{multicols}\end{quote}

\item $2\cos^2\dfrac{\theta}{2}$ equals:
\begin{quote}\begin{multicols}{2}
\noindent (A) $1+\cos\theta$ \\
(B) $1-\cos\theta$ \\
(C) $1-\sin\theta$ \\
(D) $1+\sin\theta$
\end{multicols}\end{quote}

\item $\sin\alpha=2/3$, $\cos\alpha=3/4$, $\sin2\alpha = $?
\begin{quote}\begin{multicols}{2}
\noindent (A) $1/2$ \\
(B) $1/3$ \\
(C) $1$ \\
(D) $2$
\end{multicols}\end{quote}

\item $\sin3\theta-\sin5\theta$ equals:
\begin{quote}\begin{multicols}{2}
\noindent (A) $2\sin4\theta\cdot\cos2\theta$ \\
(B) $2\cos4\theta\cdot\sin\theta$ \\
(C) $-2\cos4\theta\cdot\sin\theta$ \\
(D) $2\sin4\theta\cdot\cos\theta$
\end{multicols}\end{quote}

\item $\sin5\theta+\sin3\theta$ is equal to:
\begin{quote}\begin{multicols}{2}
\noindent (A) $2\cos2\theta\sin\theta$ \\
(B) $-2\cos4\theta\sin\theta$ \\
(C) $-2\cos4\theta\cos\theta$ \\
(D) $2\sin4\theta\cos\theta$
\end{multicols}\end{quote}

\item $\cos48^{\circ}+\cos12^{\circ} = $:
\begin{quote}\begin{multicols}{2}
\noindent (A) $2\cos18^{\circ}$ \\
(B) $3\cos18^{\circ}$ \\
(C) $3\cos19^{\circ}$ \\
(D) None of these
\end{multicols}\end{quote}

\item $\cos(\alpha+\beta)+\cos(\alpha-\beta) = $:
\begin{quote}\begin{multicols}{2}
\noindent (A) $2\sin\alpha\cos\beta$ \\
(B) $2\sin\alpha\sin\beta$ \\
(C) $2\cos\alpha\cos\beta$ \\
(D) $-2\sin\alpha\sin\beta$
\end{multicols}\end{quote}

\item $2\sin7\theta\cos3\theta$ is equal to:
\begin{quote}\begin{multicols}{2}
\noindent (A) $\sin10\theta+\sin4\theta$ \\
(B) $\sin10\theta-\sin4\theta$ \\
(C) $4\cos3\theta-3\cos\theta$ \\
(D) $\cos10\theta-\cos4\theta$
\end{multicols}\end{quote}

\item $2\sin\dfrac{P+Q}{2}\cos\dfrac{P-Q}{2} = $?
\begin{quote}\begin{multicols}{2}
\noindent (A) $\sin P+\sin Q$ \\
(B) $\sin P-\sin Q$ \\
(C) $\cos P+\cos Q$ \\
(D) $\cos P-\cos Q$
\end{multicols}\end{quote}

\item Trigonometric functions are:
\begin{quote}\begin{multicols}{2}
\noindent (A) Discontinuous \\
(B) No circular function \\
(C) Periodic \\
(D) Not periodic
\end{multicols}\end{quote}

\item The domain of $\cot x$ is:
\begin{quote}\begin{multicols}{2}
\noindent (A) $[-1,1]$ \\
(B) $(-\pi/2,\pi/2)$ \\
(C) $\mathbb{R}$ but $x\neq n\pi$ \\
(D) $\mathbb{Q}$
\end{multicols}\end{quote}

\item The domain of $\sin x$ is:
\begin{quote}\begin{multicols}{2}
\noindent (A) $[-1,1]$ \\
(B) $[-\pi/2,\pi/2]$ \\
(C) $\mathbb{R}$ \\
(D) $\mathbb{Q}$
\end{multicols}\end{quote}

\item The range of $y=\sec x$ is:
\begin{quote}\begin{multicols}{2}
\noindent (A) $[0,\pi]$ \\
(B) $[-1,1]$ \\
(C) $y\leq-1$ \\
(D) $y\leq-1,\ y\geq1$
\end{multicols}\end{quote}

\item The range of $y=\sin x$ is:
\begin{quote}\begin{multicols}{2}
\noindent (A) $(-1,1)$ \\
(B) $[-1,1)$ \\
(C) $[-1,1]$ \\
(D) $(-1,1]$
\end{multicols}\end{quote}

\item The range of $y=\cos x$ is:
\begin{quote}\begin{multicols}{2}
\noindent (A) $\mathbb{R}$ \\
(B) $-1\leq y\leq1$ \\
(C) $y>1$ or $y<-1$ \\
(D) $-1$
\end{multicols}\end{quote}

\item The range of the function $y=\cot x$ is:
\begin{quote}\begin{multicols}{2}
\noindent (A) $-1\leq y\leq1$ \\
(B) $y\leq-1$ or $y\geq1$ \\
(C) $\mathbb{R}$ \\
(D) None
\end{multicols}\end{quote}

\item The smallest positive number $P$ for which $f(x+P)=f(x)$ is called the:
\begin{quote}\begin{multicols}{2}
\noindent (A) Domain \\
(B) Range \\
(C) Co-domain \\
(D) Period
\end{multicols}\end{quote}

\item Sine is a periodic function and its period is:
\begin{quote}\begin{multicols}{2}
\noindent (A) $\pi$ \\
(B) $\dfrac{\pi}{2}$ \\
(C) $2\pi$ \\
(D) None of these
\end{multicols}\end{quote}

\item Cosine is a periodic function of period:
\begin{quote}\begin{multicols}{2}
\noindent (A) $2\pi$ \\
(B) $\pi$ \\
(C) $\dfrac{\pi}{2}$ \\
(D) None
\end{multicols}\end{quote}

\item The period of $3\cos2x$ is:
\begin{quote}\begin{multicols}{2}
\noindent (A) $\pi$ \\
(B) $2\pi$ \\
(C) $3\pi$ \\
(D) $\pi/2$
\end{multicols}\end{quote}

\item The period of $\sin2x$ is:
\begin{quote}\begin{multicols}{2}
\noindent (A) $\pi/10$ \\
(B) $\pi$ \\
(C) $\pi/2$ \\
(D) None
\end{multicols}\end{quote}

\item The period of $\sec x$ is:
\begin{quote}\begin{multicols}{2}
\noindent (A) $\pi$ \\
(B) $2\pi$ \\
(C) $5\pi$ \\
(D) $6\cdot\dfrac{2\pi}{5}$
\end{multicols}\end{quote}

\item The period of $5\tan2x$ is:
\begin{quote}\begin{multicols}{2}
\noindent (A) $\pi$ \\
(B) $2\pi$ \\
(C) $5\pi$ \\
(D) $\pi/2$
\end{multicols}\end{quote}

\item The period of $\csc10x$ is:
\begin{quote}\begin{multicols}{2}
\noindent (A) $\pi/10$ \\
(B) $2\pi/5$ \\
(C) $\pi/5$ \\
(D) None
\end{multicols}\end{quote}

\item The period of $y=\tan\theta$ is:
\begin{quote}\begin{multicols}{2}
\noindent (A) $\pi$ \\
(B) $2\pi$ \\
(C) $3\pi$ \\
(D) $4\pi$
\end{multicols}\end{quote}

\item The period of $\cot\dfrac{x}{2}$ is equal to:
\begin{quote}\begin{multicols}{2}
\noindent (A) $2\pi$ \\
(B) $4\pi$ \\
(C) $\pi$ \\
(D) None
\end{multicols}\end{quote}

\item The period of $2\cos3x$ is equal to:
\begin{quote}\begin{multicols}{2}
\noindent (A) $\pi$ \\
(B) $5\pi/2$ \\
(C) $2\pi$ \\
(D) None of these
\end{multicols}\end{quote}

\item The period of $2\csc\dfrac{x}{4}$ is:
\begin{quote}\begin{multicols}{2}
\noindent (A) $\pi$ \\
(B) $4\pi$ \\
(C) $2\pi$ \\
(D) $8\pi$
\end{multicols}\end{quote}

\item $\pi$ is the period of:
\begin{quote}\begin{multicols}{2}
\noindent (A) $\sec\theta$ \\
(B) $\csc\theta$ \\
(C) $\cot\theta$ \\
(D) $\sin3\theta$
\end{multicols}\end{quote}

\item The range of $y=3\sin2x$ is:
\begin{quote}\begin{multicols}{2}
\noindent (A) $[-1,1]$ \\
(B) $[-3,3]$ \\
(C) $[-5,5]$ \\
(D) $[-6,6]$
\end{multicols}\end{quote}

\item The graphs of trigonometric functions are:
\begin{quote}\begin{multicols}{2}
\noindent (A) Not continuous \\
(B) Line segments \\
(C) Smooth curves \\
(D) None of these
\end{multicols}\end{quote}

\item Trigonometric functions are periodic, so their curves repeat after a:
\begin{quote}\begin{multicols}{2}
\noindent (A) Fixed interval \\
(B) Not fixed interval \\
(C) Both (A) and (B) \\
(D) None of these
\end{multicols}\end{quote}

\item $a+b\cos(c\theta+d)$ is a type of \underline{\hspace{1.5cm}} function:
\begin{quote}\begin{multicols}{2}
\noindent (A) Trigonometric \\
(B) Exponential \\
(C) Constant \\
(D) Sinusoidal
\end{multicols}\end{quote}

\item The maximum value of $2+3\sin x$ is:
\begin{quote}\begin{multicols}{2}
\noindent (A) $2$ \\
(B) $3$ \\
(C) $4$ \\
(D) $5$
\end{multicols}\end{quote}

\item The maximum value of a function $m=a-|b|$ occurs whenever $\sin\theta=$:
\begin{quote}\begin{multicols}{2}
\noindent (A) $1$ \\
(B) $-1$ \\
(C) $0$ \\
(D) None
\end{multicols}\end{quote}

\item $\displaystyle\lim_{x\to1}\dfrac{x^3-3x^2+3x-1}{x-1}$ equals:
\begin{quote}\begin{multicols}{2}
\noindent (A) $0$ \\
(B) $-1$ \\
(C) $3$ \\
(D) $1$
\end{multicols}\end{quote}

\item $\displaystyle\lim_{x\to3}(2x+4) = $?
\begin{quote}\begin{multicols}{2}
\noindent (A) $12$ \\
(B) $10$ \\
(C) $6$ \\
(D) $3$
\end{multicols}\end{quote}

\item What is the limit of the function $\displaystyle\lim_{x\to5}\dfrac{x^2-25}{x-5}$?
\begin{quote}\begin{multicols}{2}
\noindent (A) $0$ \\
(B) $-10$ \\
(C) $25$ \\
(D) $10$
\end{multicols}\end{quote}

\item $^{\dagger}$ $\displaystyle\lim_{x\to4}\dfrac{x-4}{\sqrt{x}-2}$ equals:
\begin{quote}\begin{multicols}{2}
\noindent (A) $0$ \\
(B) $2$ \\
(C) $4$ \\
(D) $5$
\end{multicols}\end{quote}

\item $\displaystyle\lim_{x\to a}\dfrac{x^3-a^3}{x-a} = $:
\begin{quote}\begin{multicols}{2}
\noindent (A) Undefined \\
(B) $3a^2$ \\
(C) $a^2$ \\
(D) $0$
\end{multicols}\end{quote}

\item $\displaystyle\lim_{x\to2}\dfrac{x^5-32}{x-2} = $?
\begin{quote}\begin{multicols}{2}
\noindent (A) $40$ \\
(B) $60$ \\
(C) $80$ \\
(D) $120$
\end{multicols}\end{quote}

\item $\displaystyle\lim_{x\to2}\dfrac{x-2}{\sqrt{x}-\sqrt{2}} = $?
\begin{quote}\begin{multicols}{2}
\noindent (A) $2\sqrt{2}$ \\
(B) $2$ \\
(C) $0$ \\
(D) $1$
\end{multicols}\end{quote}

\item $\displaystyle\lim_{x\to1}\dfrac{x^3-x}{x+1}$:
\begin{quote}\begin{multicols}{2}
\noindent (A) $0$ \\
(B) $1$ \\
(C) $2$ \\
(D) $3$
\end{multicols}\end{quote}

\item What is the value of $\displaystyle\lim_{h\to0}\dfrac{|h|}{h}$?
\begin{quote}\begin{multicols}{2}
\noindent (A) $1$ \\
(B) $+1$ \\
(C) $0$ \\
(D) None (limit does not exist)
\end{multicols}\end{quote}

\item $\displaystyle\lim_{\theta\to0}\dfrac{\sin\theta}{\theta}=1$ where $\theta$ is measured in:
\begin{quote}\begin{multicols}{2}
\noindent (A) Degree \\
(B) Radian \\
(C) Minutes \\
(D) Second
\end{multicols}\end{quote}

\item $\displaystyle\lim_{x\to0}\dfrac{\sin x^{\circ}}{x^{\circ}} = $:
\begin{quote}\begin{multicols}{2}
\noindent (A) $1$ \\
(B) $2$ \\
(C) $3$ \\
(D) $4$
\end{multicols}\end{quote}

\item $\displaystyle\lim_{x\to0}\dfrac{\sin x^{\circ}}{x}$ is equal to:
\begin{quote}\begin{multicols}{2}
\noindent (A) $\dfrac{\pi}{180}$ \\
(B) $\dfrac{180}{\pi}$ \\
(C) $180$ \\
(D) $1$
\end{multicols}\end{quote}

\item $\displaystyle\lim_{x\to0}\dfrac{\sin ax}{\sin bx} = $:
\begin{quote}\begin{multicols}{2}
\noindent (A) $\dfrac{a}{b}$ \\
(B) $-\dfrac{b}{a}$ \\
(C) $-\dfrac{a}{b}$ \\
(D) $\dfrac{b}{a}$
\end{multicols}\end{quote}

\item $\displaystyle\lim_{x\to0}\dfrac{x}{\sin2x}$ is equal to:
\begin{quote}\begin{multicols}{2}
\noindent (A) $\dfrac{1}{2}$ \\
(B) $2$ \\
(C) $3$ \\
(D) $-2$
\end{multicols}\end{quote}

\item $\displaystyle\lim_{\theta\to0}\dfrac{1-\cos\theta}{\theta}$ equal to:
\begin{quote}\begin{multicols}{2}
\noindent (A) Zero \\
(B) $-1$ \\
(C) $1$ \\
(D) None
\end{multicols}\end{quote}

\item $^{\dagger}$ What is the value of $\displaystyle\lim_{x\to0}\dfrac{\cos\left(\dfrac{3\pi}{2}-x\right)-\cos\dfrac{3\pi}{2}}{x}$?
\begin{quote}\begin{multicols}{2}
\noindent (A) $-1$ \\
(B) $1$ \\
(C) $2$ \\
(D) $-2$
\end{multicols}\end{quote}

\item $\displaystyle\lim_{x\to0}\dfrac{\tan x}{x} = $:
\begin{quote}\begin{multicols}{2}
\noindent (A) $0$ \\
(B) $-1$ \\
(C) $1$ \\
(D) None of these
\end{multicols}\end{quote}

\item $\displaystyle\lim_{\theta\to0}\dfrac{\sin7\theta}{\theta} = $:
\begin{quote}\begin{multicols}{2}
\noindent (A) $7$ \\
(B) $1$ \\
(C) $-7$ \\
(D) $-1$
\end{multicols}\end{quote}

\item $\displaystyle\lim_{x\to0}\dfrac{x}{\tan x} = $?
\begin{quote}\begin{multicols}{2}
\noindent (A) $0$ \\
(B) $1$ \\
(C) $2$ \\
(D) $3$
\end{multicols}\end{quote}

\item $\displaystyle\lim_{\theta\to0}\dfrac{\sin^2\theta}{\theta} = $?
\begin{quote}\begin{multicols}{2}
\noindent (A) $0$ \\
(B) $1$ \\
(C) $-1$ \\
(D) $2$
\end{multicols}\end{quote}

\item $\displaystyle\lim_{n\to\infty}\left(1+\dfrac{1}{n}\right)^{2n} = $?
\begin{quote}\begin{multicols}{2}
\noindent (A) $e^{-1}$ \\
(B) $e^{-1/2}$ \\
(C) $e^2$ \\
(D) $e^3$
\end{multicols}\end{quote}

\item $\displaystyle\lim_{x\to0}(1+x)^{1/x} = $?
\begin{quote}\begin{multicols}{2}
\noindent (A) $e^x$ \\
(B) $e$ \\
(C) $e^1$ \\
(D) $e^2$
\end{multicols}\end{quote}

\item $\displaystyle\lim_{n\to\infty}\left(1+\dfrac{3}{n}\right)^{2n} = $?
\begin{quote}\begin{multicols}{2}
\noindent (A) $e$ \\
(B) $e^2$ \\
(C) $e^3$ \\
(D) $e^6$
\end{multicols}\end{quote}

\item $\displaystyle\lim_{x\to0}\dfrac{e^x-1}{x} = $?
\begin{quote}\begin{multicols}{2}
\noindent (A) $0$ \\
(B) $1$ \\
(C) $e$ \\
(D) $2$
\end{multicols}\end{quote}

\item $^{\dagger}$ $\displaystyle\lim_{x\to\infty}\left(1+\dfrac{2}{x}\right)^{x} = $:
\begin{quote}\begin{multicols}{2}
\noindent (A) $e$ \\
(B) $e^2$ \\
(C) $e^3$ \\
(D) None
\end{multicols}\end{quote}

\item $\displaystyle\lim_{n\to\infty}\left(1-\dfrac{1}{n}\right)^{n} = $:
\begin{quote}\begin{multicols}{2}
\noindent (A) $e$ \\
(B) $e^2$ \\
(C) $\dfrac{1}{e}$ \\
(D) $e^1$
\end{multicols}\end{quote}

\item $\displaystyle\lim_{n\to\infty}\left(\dfrac{5n+1}{5n}\right)^{n} = $:
\begin{quote}\begin{multicols}{2}
\noindent (A) $e^{1/5}$ \\
(B) $e^5$ \\
(C) $e^{-5}$ \\
(D) $e^{1/5}$
\end{multicols}\end{quote}

\item $\displaystyle\lim_{x\to\infty}\left(\dfrac{1}{2}\right)^{x} = $:
\begin{quote}\begin{multicols}{2}
\noindent (A) $0$ \\
(B) $1$ \\
(C) $2$ \\
(D) $3$
\end{multicols}\end{quote}

\item $\displaystyle\lim_{x\to\infty}\dfrac{2-3x}{\sqrt{3+4x^2}} = $:
\begin{quote}\begin{multicols}{2}
\noindent (A) $\dfrac{3}{2}$ \\
(B) $-\dfrac{3}{2}$ \\
(C) $\pm\dfrac{3}{2}$ \\
(D) None of these
\end{multicols}\end{quote}

\item $\displaystyle\lim_{x\to\pi}\dfrac{\sin x}{\pi-x}$ equals:
\begin{quote}\begin{multicols}{2}
\noindent (A) Zero \\
(B) $1$ \\
(C) $2$ \\
(D) $3$
\end{multicols}\end{quote}

\item $\displaystyle\lim_{h\to0}(1-2h)^{1/h}$ is equal to:
\begin{quote}\begin{multicols}{2}
\noindent (A) $e$ \\
(B) $e^{-1}$ \\
(C) $e^2$ \\
(D) $e^{-2}$
\end{multicols}\end{quote}

\item If $f(x)$ is continuous at $x=c$ then:
\begin{quote}\begin{multicols}{2}
\noindent (A) $\displaystyle\lim_{x\to c}f(x)\neq0$ \\
(B) $\displaystyle\lim_{x\to c}f(x)\neq1$ \\
(C) $\displaystyle\lim_{x\to c}f(x)=f(c)$ \\
(D) $\displaystyle\lim_{x\to c}f(x)\neq f(c)$
\end{multicols}\end{quote}

\item The function $f(x)=\dfrac{2+3x}{2x}$ is not continuous at:
\begin{quote}\begin{multicols}{2}
\noindent (A) $x=-3$ \\
(B) $x=0$ \\
(C) $x=1$ \\
(D) $x=-1$
\end{multicols}\end{quote}

\item A function $f(x)$ is not continuous at $x=c$; then it is called \underline{\hspace{1.5cm}} at $x=c$:
\begin{quote}\begin{multicols}{2}
\noindent (A) Discontinuous \\
(B) Defined \\
(C) Stationary \\
(D) None of these
\end{multicols}\end{quote}

\item The function $f(x)=\dfrac{x^2-1}{x-1}$ is discontinuous at:
\begin{quote}\begin{multicols}{2}
\noindent (A) $x=-1$ \\
(B) $x=1$ \\
(C) $x=2$ \\
(D) $x=0$
\end{multicols}\end{quote}

\item Limit and continuity of a \underline{\hspace{1.5cm}} function is a fundamental concept of calculus:
\begin{quote}\begin{multicols}{2}
\noindent (A) Piecewise function \\
(B) Exponential \\
(C) Logarithmic \\
(D) Transcendental
\end{multicols}\end{quote}

\item The notation used by Leibniz for the derivative is:
\begin{quote}\begin{multicols}{2}
\noindent (A) $\dfrac{dy}{dx}$ \\
(B) $f(x)$ \\
(C) $f'(x)$ \\
(D) $Df(x)$
\end{multicols}\end{quote}

\item The notation $f'(x)$ is denoted after:
\begin{quote}\begin{multicols}{2}
\noindent (A) Leibniz \\
(B) Cauchy \\
(C) Lagrange \\
(D) Newton
\end{multicols}\end{quote}

\item The notation for the derivative used by Newton is:
\begin{quote}\begin{multicols}{2}
\noindent (A) $\dfrac{dy}{dx}$ \\
(B) $Df(x)$ \\
(C) $f(x)$ \\
(D) $f'(x)$
\end{multicols}\end{quote}

\item $\displaystyle\lim_{\delta x\to0}\dfrac{f(x+\delta x)-f(x)}{\delta x}$ is equal to:
\begin{quote}\begin{multicols}{2}
\noindent (A) $f(x)$ \\
(B) $f'(a)$ \\
(C) $f'(x)$ \\
(D) $f'(0)$
\end{multicols}\end{quote}

\item $\displaystyle\lim_{x\to a}\dfrac{f(x)-f(a)}{x-a}$ is equal to:
\begin{quote}\begin{multicols}{2}
\noindent (A) $f'(x)$ \\
(B) $f'(a)$ \\
(C) $f'(2)$ \\
(D) None
\end{multicols}\end{quote}

\item $\displaystyle\lim_{h\to0}\dfrac{f(a+h)-f(a)}{h} = $:
\begin{quote}\begin{multicols}{2}
\noindent (A) $f(a)$ \\
(B) $f'(a+h)$ \\
(C) $f'(x)$ \\
(D) $f'(a)$
\end{multicols}\end{quote}

\item If $f(x)=c$ (constant), then $f'(1)=$:
\begin{quote}\begin{multicols}{2}
\noindent (A) $c$ \\
(B) $1$ \\
(C) $0$ \\
(D) None
\end{multicols}\end{quote}

\item $f(x)=x^{2/3}$, then $f'(8)=$:
\begin{quote}\begin{multicols}{2}
\noindent (A) $1/2$ \\
(B) $2/3$ \\
(C) $1/3$ \\
(D) $3$
\end{multicols}\end{quote}

\item $f(x)=(1+x)^n$, $f'(0)$ will be:
\begin{quote}\begin{multicols}{2}
\noindent (A) $0$ \\
(B) $n$ \\
(C) $1$ \\
(D) $n!$
\end{multicols}\end{quote}

\item If $f(x)=x^{100}$, then $f'(1)=$:
\begin{quote}\begin{multicols}{2}
\noindent (A) $100$ \\
(B) $99$ \\
(C) $50$ \\
(D) $50$
\end{multicols}\end{quote}

\item $^{\dagger}$ If $f(x)=a+x$, then $f'(x)=$:
\begin{quote}\begin{multicols}{2}
\noindent (A) $1$ \\
(B) $0$ \\
(C) $a$ \\
(D) $x$
\end{multicols}\end{quote}

\item $\dfrac{d}{dx}(ax+b)^n$ is equal to:
\begin{quote}\begin{multicols}{2}
\noindent (A) $n(ax+b)^{n-1}$ \\
(B) $n(ax+b)$ \\
(C) $nax^{n-1}$ \\
(D) $na(ax+b)^{n-1}$
\end{multicols}\end{quote}

\item $^{\dagger}$ $\dfrac{d}{dx}(4x+7)^9 = $:
\begin{quote}\begin{multicols}{2}
\noindent (A) $36(4x+7)^8$ \\
(B) $9(4x+7)^8$ \\
(C) $28(4x+7)^8$ \\
(D) $63(4x+7)^8$
\end{multicols}\end{quote}

\item $\dfrac{d}{dx}\left(\dfrac{1}{ax+b}\right)$ equals:
\begin{quote}\begin{multicols}{2}
\noindent (A) $\dfrac{1}{(ax+b)^2}$ \\
(B) $\dfrac{a}{(ax+b)^2}$ \\
(C) $\dfrac{-a}{(ax+b)^2}$ \\
(D) $\ln(ax+b)$
\end{multicols}\end{quote}

\item $\dfrac{d}{dx}(x^n)$ is equal to:
\begin{quote}\begin{multicols}{2}
\noindent (A) $nx^{n-1}$ \\
(B) $x^{n-1}$ \\
(C) $x^{n+1}$ \\
(D) None
\end{multicols}\end{quote}

\item $\dfrac{d}{dx}\big[(x-5)(3-x)\big] = $:
\begin{quote}\begin{multicols}{2}
\noindent (A) $2x+8$ \\
(B) $-2x+8$ \\
(C) $2x-8$ \\
(D) $x+8$
\end{multicols}\end{quote}

\item $\dfrac{d}{dx}(x^3+1)^9 = $?
\begin{quote}\begin{multicols}{2}
\noindent (A) $9(x^3+1)^8$ \\
(B) $9x^3(x^3+1)^8$ \\
(C) $27x^2(x^3+1)^8$ \\
(D) $27x(x^3+1)^8$
\end{multicols}\end{quote}

\item Differentiate $(x-1)(x+2)^2$:
\begin{quote}\begin{multicols}{2}
\noindent (A) $2x(x+2)$ \\
(B) $2x(x+2)$ \\
(C) $2(x-1)$ \\
(D) $3x(x+2)$
\end{multicols}\end{quote}

\item $\dfrac{d}{dx}\big(cf(x)\big) = $:
\begin{quote}\begin{multicols}{2}
\noindent (A) $cf'(x)$ \\
(B) $c'f$ \\
(C) $[cf(x)]$ \\
(D) All of these
\end{multicols}\end{quote}

\item The derivative of $x^0$ with respect to $x$ is:
\begin{quote}\begin{multicols}{2}
\noindent (A) $x^{0-1}$ \\
(B) $c$ \\
(C) Zero \\
(D) $1$
\end{multicols}\end{quote}

\item $\dfrac{d}{dx}\left[\dfrac{1}{g(x)}\right] = $:
\begin{quote}\begin{multicols}{2}
\noindent (A) $\dfrac{1}{g^2(x)}$ \\
(B) $\dfrac{-g(x)}{g(x)}$ \\
(C) $\dfrac{1}{[g(x)]^2}$ \\
(D) $\dfrac{-g'(x)}{[g(x)]^2}$
\end{multicols}\end{quote}

\item The derivative of $\dfrac{x-1}{x+1}$ with respect to $x$ is:
\begin{quote}\begin{multicols}{2}
\noindent (A) $\dfrac{-1}{(x+1)^2}$ \\
(B) $\dfrac{2}{(x+1)^2}$ \\
(C) $\dfrac{2x}{(x+1)^2}$ \\
(D) $\dfrac{-2x+1}{(x+1)^2}$
\end{multicols}\end{quote}

\item $\dfrac{d}{dx}\big(f\circ g(x)\big) = $?
\begin{quote}\begin{multicols}{2}
\noindent (A) $f'(g(x))$ \\
(B) $f(g'(x))$ \\
(C) $f[g(x)]g'(x)$ \\
(D) $f'[g(x)]g'(x)$
\end{multicols}\end{quote}

\item $\dfrac{d}{dx}\left(\dfrac{2}{x}\right) = $?
\begin{quote}\begin{multicols}{2}
\noindent (A) $\dfrac{1}{x}$ \\
(B) $\dfrac{-2}{x}$ \\
(C) $\dfrac{2}{x^2}$ \\
(D) $\dfrac{-2}{x^2}$
\end{multicols}\end{quote}

\item $y=x-\dfrac{1}{x}$, then $\dfrac{dy}{dx} = $?
\begin{quote}\begin{multicols}{2}
\noindent (A) $1+\dfrac{1}{x^2}$ \\
(B) $1-\dfrac{1}{x}$ \\
(C) $1+\dfrac{1}{x}$ \\
(D) $\dfrac{x^3-x}{x+1}$
\end{multicols}\end{quote}

\item What is the derivative of $\dfrac{x^3-x}{x+1}$ with respect to $x$:
\begin{quote}\begin{multicols}{2}
\noindent (A) $\dfrac{2x^3+3x^2-1}{(x+1)^2}$ \\
(B) $2x-1$ \\
(C) $1-2x$ \\
(D) $\dfrac{2x^3+3x^2-2x+1}{(x+1)^2}$
\end{multicols}\end{quote}

\item The derivative of $\dfrac{x^3+2x^2}{x^3}$ equals:
\begin{quote}\begin{multicols}{2}
\noindent (A) $\dfrac{2}{x^2}$ \\
(B) $-\dfrac{2}{x^2}$ \\
(C) $\dfrac{1}{2x^2}$ \\
(D) $-\dfrac{1}{2x^2}$
\end{multicols}\end{quote}

\item $\dfrac{dy}{dx} = -\dfrac{1}{x^2}$ if:
\begin{quote}\begin{multicols}{2}
\noindent (A) $y=-\dfrac{1}{x}$ \\
(B) $y=-x$ \\
(C) $y=\ln x$ \\
(D) $y=\dfrac{1}{x}$
\end{multicols}\end{quote}

\item $\dfrac{d}{dx}(ax^n+bx^m) = $:
\begin{quote}\begin{multicols}{2}
\noindent (A) $nax^{n-1}+mbx^{m-1}$ \\
(B) $nax^{n-1}+mbx^{m+1}$ \\
(C) $ax^{n-1}+bx^{m-1}$ \\
(D) None of these
\end{multicols}\end{quote}

\item $\dfrac{d}{dx}\left(\dfrac{x}{a}\right) = $:
\begin{quote}\begin{multicols}{2}
\noindent (A) $\dfrac{x}{a}$ \\
(B) $\dfrac{1}{a}$ \\
(C) $\dfrac{1}{a^2}$ \\
(D) $\dfrac{a}{x}$
\end{multicols}\end{quote}

\item $\dfrac{d}{dx}\left(\dfrac{x^2-4}{x-2}\right)$ equals:
\begin{quote}\begin{multicols}{2}
\noindent (A) $0$ \\
(B) $1$ \\
(C) $x+2$ \\
(D) $x-2$
\end{multicols}\end{quote}

\item $^{\dagger}$ $f(x)=\dfrac{1}{x^m}$, then $f'(x) = $?
\begin{quote}\begin{multicols}{2}
\noindent (A) $-x^{m-1}$ \\
(B) $-mx^{-m-1}$ \\
(C) $-x^{-m+1}$ \\
(D) $-x^{-m-1}$
\end{multicols}\end{quote}

\item If $s(t)=5t^2-2t$, what will be the velocity after two seconds:
\begin{quote}\begin{multicols}{2}
\noindent (A) $14$ \\
(B) $16$ \\
(C) $18$ \\
(D) $20$
\end{multicols}\end{quote}

\item Which of the following is a vector:
\begin{quote}\begin{multicols}{2}
\noindent (A) Time \\
(B) Work \\
(C) Density \\
(D) Electric field
\end{multicols}\end{quote}

\item A vector whose initial and terminal points are the same is called:
\begin{quote}\begin{multicols}{2}
\noindent (A) Null vector \\
(B) Not a null vector \\
(C) Unit vector \\
(D) None of these
\end{multicols}\end{quote}

\item A vector with magnitude $1$ is called:
\begin{quote}\begin{multicols}{2}
\noindent (A) Null vector \\
(B) Unit vector \\
(C) Zero vector \\
(D) Constant vector
\end{multicols}\end{quote}

\item A physical quantity that can be specified by a number along with a unit (no direction) is called a:
\begin{quote}\begin{multicols}{2}
\noindent (A) Vector \\
(B) Scalar \\
(C) Constant \\
(D) Constant vector
\end{multicols}\end{quote}

\item The magnitude of the vector $\vec{v}=a\hat{i}+b\hat{j}+c\hat{k}$ is:
\begin{quote}\begin{multicols}{2}
\noindent (A) $a^2+b^2+c^2$ \\
(B) $a+b+c$ \\
(C) $a^3+b^3+c^3$ \\
(D) $\sqrt{a^2+b^2+c^2}$
\end{multicols}\end{quote}

\item The position vector of any point in the $xy$-plane is:
\begin{quote}\begin{multicols}{2}
\noindent (A) $\vec{r}=x\hat{i}+y\hat{j}$ \\
(B) $\vec{r}=y\hat{i}+z\hat{k}$ \\
(C) $\vec{r}=x\hat{i}+z\hat{k}$ \\
(D) None of these
\end{multicols}\end{quote}

\item Magnitude of the vector $\vec{v}=3\hat{j}+4\hat{k}$ is:
\begin{quote}\begin{multicols}{2}
\noindent (A) $29$ \\
(B) $\sqrt{29}$ \\
(C) $28$ \\
(D) $\sqrt{28}$
\end{multicols}\end{quote}

\item Let $\overrightarrow{OA}=\vec{a}$ and $\overrightarrow{OB}=\vec{b}$ then $\overrightarrow{AB}=$:
\begin{quote}\begin{multicols}{2}
\noindent (A) $\vec{a}-\vec{b}$ \\
(B) $\vec{a}+\vec{b}$ \\
(C) $\vec{b}-\vec{a}$ \\
(D) $\vec{b}+\vec{a}$
\end{multicols}\end{quote}

\item Length of the vector $2\hat{i}-\hat{j}-2\hat{k}$ is:
\begin{quote}\begin{multicols}{2}
\noindent (A) $2$ \\
(B) $4$ \\
(C) $3$ \\
(D) $5$
\end{multicols}\end{quote}

\item If $P=(2,3)$, $Q=(6,-2)$, then the vector $\overrightarrow{PQ}$ is:
\begin{quote}\begin{multicols}{2}
\noindent (A) $\dfrac{4\hat{j}-5\hat{i}}{\sqrt{41}}$ \\
(B) $\dfrac{4\hat{i}-5\hat{j}}{\sqrt{41}}$ \\
(C) $4\hat{i}-5\hat{j}$ \\
(D) $4-5\hat{k}$
\end{multicols}\end{quote}

\item The magnitude of a vector is also called its:
\begin{quote}\begin{multicols}{2}
\noindent (A) Parameter \\
(B) Variable \\
(C) Point \\
(D) Norm
\end{multicols}\end{quote}

\item Magnitude of the vector $\vec{v}=-\hat{i}+\hat{j}$ is:
\begin{quote}\begin{multicols}{2}
\noindent (A) $\sqrt{1}$ \\
(B) $\sqrt{2}$ \\
(C) $5$ \\
(D) $8$
\end{multicols}\end{quote}

\item Magnitude of the vector $\vec{v}=[3,4]$ is:
\begin{quote}\begin{multicols}{2}
\noindent (A) $3$ \\
(B) $4$ \\
(C) $5$ \\
(D) $6$
\end{multicols}\end{quote}

\item The dot product of vectors $\vec{u}$ and $\vec{v}$ is:
\begin{quote}\begin{multicols}{2}
\noindent (A) $\cos\theta\sin\theta$ \\
(B) $|\vec{u}||\vec{v}|=\cos\theta$ \\
(C) $uv\cos\theta$ \\
(D) None
\end{multicols}\end{quote}

\item Commutative law holds in:
\begin{quote}\begin{multicols}{2}
\noindent (A) Vector product \\
(B) Cross product in three vectors \\
(C) Inner product \\
(D) None of these
\end{multicols}\end{quote}

\item $\vec{i}\cdot\vec{i}$ is equal to:
\begin{quote}\begin{multicols}{2}
\noindent (A) $0$ \\
(B) $1$ \\
(C) $2$ \\
(D) $3$
\end{multicols}\end{quote}

\item $\vec{k}\cdot\vec{k}$ is equal to:
\begin{quote}\begin{multicols}{2}
\noindent (A) $-1$ \\
(B) $0$ \\
(C) $1$ \\
(D) $-2$
\end{multicols}\end{quote}

\item $\hat{i}\cdot\hat{j}=$:
\begin{quote}\begin{multicols}{2}
\noindent (A) $0$ \\
(B) $1$ \\
(C) $2$ \\
(D) $3$
\end{multicols}\end{quote}

\item $\vec{a}$ \& $\vec{b}$ are perpendicular if:
\begin{quote}\begin{multicols}{2}
\noindent (A) $\vec{a}\times\vec{b}=0$ \\
(B) $\vec{a}\cdot\vec{b}=0$ \\
(C) $\vec{a}=\vec{b}$ \\
(D) $\vec{a}=-\vec{b}$
\end{multicols}\end{quote}

\item The angle between the vectors $2\hat{i}+3\hat{j}+\hat{k}$ and $2\hat{i}-\hat{j}-\hat{k}$ is:
\begin{quote}\begin{multicols}{2}
\noindent (A) $30^\circ$ \\
(B) $40^\circ$ \\
(C) $60^\circ$ \\
(D) $90^\circ$
\end{multicols}\end{quote}

\item Projection of $\vec{a}$ along $\vec{b}$ is:
\begin{quote}\begin{multicols}{2}
\noindent (A) $\vec{a}-\vec{b}$ \\
(B) $\vec{a}\cdot\hat{b}$ \\
(C) $\hat{a}\cdot\hat{b}$ \\
(D) $\vec{a}\cdot\hat{b}$
\end{multicols}\end{quote}

\item Projection of $\vec{a}=\hat{i}-\hat{k}$ along $\vec{b}=\hat{j}+\hat{k}$ is:
\begin{quote}\begin{multicols}{2}
\noindent (A) $-\dfrac{1}{\sqrt{2}}$ \\
(B) $\dfrac{1}{\sqrt{2}}$ \\
(C) $\dfrac{3}{\sqrt{2}}$ \\
(D) $\dfrac{1}{2}$
\end{multicols}\end{quote}

\item If any two vectors of a scalar product are equal then its value is:
\begin{quote}\begin{multicols}{2}
\noindent (A) $1$ \\
(B) $2$ \\
(C) $-1$ \\
(D) $0$
\end{multicols}\end{quote}

\item Projection of a vector $\vec{v}$ along $\vec{u}$ is:
\begin{quote}\begin{multicols}{2}
\noindent (A) $\dfrac{\vec{u}\times\vec{v}}{|\vec{v}|}$ \\
(B) $\dfrac{\vec{u}\cdot\vec{v}}{|\vec{u}|}$ \\
(C) $\dfrac{\vec{u}\cdot\vec{v}}{\vec{v}}$ \\
(D) $\dfrac{\vec{u}\cdot\vec{v}}{|\vec{v}|}$
\end{multicols}\end{quote}

\item If vectors $2a\hat{i}+\hat{j}-\hat{k}$ and $\hat{i}+a\hat{j}-4\hat{k}$ are perpendicular to each other then the value of $a$ is:
\begin{quote}\begin{multicols}{2}
\noindent (A) $3$ \\
(B) $\dfrac{1}{3}$ \\
(C) $\dfrac{2}{3}$ \\
(D) $\dfrac{4}{3}$
\end{multicols}\end{quote}

\item For what value of $p$, $[2,p,5]$ is perpendicular to $[3,1,p]$:
\begin{quote}\begin{multicols}{2}
\noindent (A) $1$ \\
(B) $-2$ \\
(C) $-1$ \\
(D) $3$
\end{multicols}\end{quote}

\item If the dot product of two vectors is zero then the vectors will be:
\begin{quote}\begin{multicols}{2}
\noindent (A) Perpendicular \\
(B) Parallel \\
(C) Collinear \\
(D) None of these
\end{multicols}\end{quote}

\item If $2\hat{i}+a\hat{j}+5\hat{k}$ and $3\hat{i}+\hat{j}+a\hat{k}$ are perpendicular, then $a=$:
\begin{quote}\begin{multicols}{2}
\noindent (A) $0$ \\
(B) $1$ \\
(C) $-1$ \\
(D) $2$
\end{multicols}\end{quote}

\item For a vector $\vec{v}=a\hat{i}+b\hat{j}+c\hat{k}$, the projection of $\vec{v}$ along $\hat{k}$ is:
\begin{quote}\begin{multicols}{2}
\noindent (A) $a$ \\
(B) $b$ \\
(C) $c$ \\
(D) $a+b+c$
\end{multicols}\end{quote}

\item Zero velocity is perpendicular to:
\begin{quote}\begin{multicols}{2}
\noindent (A) Every vector \\
(B) Unit vector only \\
(C) Position vector only \\
(D) Not any vector
\end{multicols}\end{quote}

\item Angle in a semicircle is:
\begin{quote}\begin{multicols}{2}
\noindent (A) $30^\circ$ \\
(B) $45^\circ$ \\
(C) $90^\circ$ \\
(D) $60^\circ$
\end{multicols}\end{quote}

\item Vector product of two vectors is a:
\begin{quote}\begin{multicols}{2}
\noindent (A) Scalar quantity \\
(B) Unit vector \\
(C) Vector quantity \\
(D) Null vector
\end{multicols}\end{quote}

\item For two non-zero vectors $\vec{a}$ and $\vec{b}$, $|\vec{a}\times\vec{b}|$ is equal to:
\begin{quote}\begin{multicols}{2}
\noindent (A) $ab\cos\theta$ \\
(B) $ab\sin\theta$ \\
(C) $ab\cos\theta\,\hat{n}$ \\
(D) $ab\sin\theta\,\hat{n}$
\end{multicols}\end{quote}

\item The vectors $\vec{u}$ and $\vec{v}$ are parallel if the angle between them is:
\begin{quote}\begin{multicols}{2}
\noindent (A) $\dfrac{\pi}{2}$ \\
(B) $\pi$ \\
(C) $\dfrac{\pi}{4}$ \\
(D) $\dfrac{\pi}{3}$
\end{multicols}\end{quote}

\item $\vec{j}\times\vec{k}=$:
\begin{quote}\begin{multicols}{2}
\noindent (A) $\vec{i}$ \\
(B) $-\vec{i}$ \\
(C) $1$ \\
(D) $0$
\end{multicols}\end{quote}

\item $\hat{i}\times\hat{k}=$:
\begin{quote}\begin{multicols}{2}
\noindent (A) $0$ \\
(B) $\hat{j}$ \\
(C) $-\hat{j}$ \\
(D) $1$
\end{multicols}\end{quote}

\item Area of the triangle whose adjacent sides are $3\hat{i}+4\hat{j}$ and $12\hat{i}+9\hat{j}$ is:
\begin{quote}\begin{multicols}{2}
\noindent (A) $\dfrac{45}{2}$ \\
(B) $\dfrac{55}{2}$ \\
(C) $\dfrac{21}{2}$ \\
(D) $\dfrac{75}{2}$
\end{multicols}\end{quote}

\item $\vec{k}\times\vec{j}=$ is equal to:
\begin{quote}\begin{multicols}{2}
\noindent (A) $\hat{i}$ \\
(B) $1$ \\
(C) $-1$ \\
(D) None
\end{multicols}\end{quote}

\item Non-zero vectors $\vec{a}$ and $\vec{b}$ are parallel if $\vec{a}\times\vec{b}=$:
\begin{quote}\begin{multicols}{2}
\noindent (A) $0$ \\
(B) $1$ \\
(C) $-1$ \\
(D) $\vec{a},\vec{b}$
\end{multicols}\end{quote}

\item Area of a triangle with $\vec{u}$ and $\vec{v}$ as sides is:
\begin{quote}\begin{multicols}{2}
\noindent (A) $\vec{u}\times\vec{v}$ \\
(B) $\dfrac{1}{2}|\vec{u}\times\vec{v}|$ \\
(C) $\vec{u}\cdot\vec{v}$ \\
(D) $uv$
\end{multicols}\end{quote}

\item A unit vector perpendicular to both $\vec{u}$ and $\vec{v}$ is given by:
\begin{quote}\begin{multicols}{2}
\noindent (A) $\vec{u}\times\vec{v}$ \\
(B) $\vec{u}\cdot\vec{v}$ \\
(C) $\dfrac{\vec{u}\times\vec{v}}{|\vec{u}\times\vec{v}|}$ \\
(D) None of these
\end{multicols}\end{quote}

\item If $\vec{a}$ and $\vec{b}$ are two non-zero vectors then the angle between $\vec{a}$ and $\vec{a}\times\vec{b}$ is always:
\begin{quote}\begin{multicols}{2}
\noindent (A) $0$ \\
(B) $30^\circ$ \\
(C) $90^\circ$ \\
(D) $60^\circ$
\end{multicols}\end{quote}

\item $2\hat{i}\times\hat{j}=$:
\begin{quote}\begin{multicols}{2}
\noindent (A) $2\hat{k}$ \\
(B) $-2\hat{k}$ \\
(C) $4\hat{k}$ \\
(D) $0$
\end{multicols}\end{quote}

\item A vector perpendicular to $2\hat{i}-\hat{j}+\hat{k}$ and $4\hat{i}+2\hat{j}+8\hat{k}$ is:
\begin{quote}\begin{multicols}{2}
\noindent (A) $-\hat{i}+6\hat{j}+8\hat{k}$ \\
(B) $-10\hat{i}-12\hat{j}+8\hat{k}$ \\
(C) $\hat{i}+6\hat{j}+8\hat{k}$ \\
(D) None of these
\end{multicols}\end{quote}

\item The number of vectors of unit length perpendicular to $\vec{a}$ and $\vec{b}$ are:
\begin{quote}\begin{multicols}{2}
\noindent (A) $1$ \\
(B) $2$ \\
(C) $3$ \\
(D) Infinite
\end{multicols}\end{quote}

\item The magnitude of the dot and cross products of two vectors are $6$ and $6\sqrt{3}$ respectively; the angle between the vectors is:
\begin{quote}\begin{multicols}{2}
\noindent (A) $90^\circ$ \\
(B) $60^\circ$ \\
(C) $30^\circ$ \\
(D) $45^\circ$
\end{multicols}\end{quote}

\item If $\vec{a}\times\vec{b}=0$ and $\vec{a}\cdot\vec{b}=0$, then:
\begin{quote}\begin{multicols}{2}
\noindent (A) $\vec{a}$ and $\vec{b}$ are parallel \\
(B) $\vec{a}$ and $\vec{b}$ are perpendicular \\
(C) Either $\vec{a}=0$ or $\vec{b}=0$ \\
(D) None of these
\end{multicols}\end{quote}

\item If $\vec{a}$ and $\vec{b}$ are non-zero vectors then $|\vec{a}\times\vec{b}|=$:
\begin{quote}\begin{multicols}{2}
\noindent (A) $ab$ \\
(B) $\vec{a}\cdot\vec{b}$ \\
(C) $\vec{b}\times\vec{a}$ \\
(D) $-\vec{b}\times\vec{a}$
\end{multicols}\end{quote}

\item $\vec{U}\times\vec{V}=$:
\begin{quote}\begin{multicols}{2}
\noindent (A) $\vec{V}\times\vec{U}$ \\
(B) $-\vec{V}\times\vec{U}$ \\
(C) Both a and b \\
(D) None of these
\end{multicols}\end{quote}

\item $\vec{u}\times(\vec{v}\cdot\vec{w})$ is:
\begin{quote}\begin{multicols}{2}
\noindent (A) Scalar product \\
(B) Vector product \\
(C) Cross product \\
(D) Meaningless
\end{multicols}\end{quote}

\item If any two vectors of a scalar triple product are equal, then its value is:
\begin{quote}\begin{multicols}{2}
\noindent (A) $1$ \\
(B) $-1$ \\
(C) $2$ \\
(D) $0$
\end{multicols}\end{quote}

\item Vector triple product of three non-zero vectors $\vec{a},\vec{b},\vec{c}$ is denoted by:
\begin{quote}\begin{multicols}{2}
\noindent (A) $\vec{a}\times(\vec{b}\times\vec{c})$ \\
(B) $\vec{a}\cdot(\vec{b}\times\vec{c})$ \\
(C) $\vec{a}\cdot(\vec{b}+\vec{c})$ \\
(D) None of these
\end{multicols}\end{quote}

\item The vector $\vec{a}$ lies in the plane of the vectors $\vec{b}$ and $\vec{c}$ then $\vec{a}\cdot\vec{b}\times\vec{c}=$ (scalar triple product of coplanar vectors) is:
\begin{quote}\begin{multicols}{2}
\noindent (A) $0$ \\
(B) $1$ \\
(C) $-1$ \\
(D) $2$
\end{multicols}\end{quote}

\item If $\vec{a}=[1,0,1];\ \vec{b}=[0,2,0];\ \vec{c}=[0,0,3]$ then $\vec{a}\cdot(\vec{b}\times\vec{c})$ is:
\begin{quote}\begin{multicols}{2}
\noindent (A) $1$ \\
(B) $2$ \\
(C) $3$ \\
(D) $6$
\end{multicols}\end{quote}

\item $\hat{i}\cdot(\hat{j}\times\hat{k})$ is equal to:
\begin{quote}\begin{multicols}{2}
\noindent (A) $0$ \\
(B) $1$ \\
(C) $-1$ \\
(D) $2$
\end{multicols}\end{quote}

\item $[\hat{k}\ \hat{i}\ \hat{j}]=$:
\begin{quote}\begin{multicols}{2}
\noindent (A) $1$ \\
(B) $-1$ \\
(C) $0$ \\
(D) $3$
\end{multicols}\end{quote}

\item The vectors $\vec{a},\vec{b},\vec{c}$ are said to be coplanar if $(\vec{a}\cdot\vec{b}\times\vec{c})=$:
\begin{quote}\begin{multicols}{2}
\noindent (A) $1$ \\
(B) $0$ \\
(C) $2$ \\
(D) $ab$
\end{multicols}\end{quote}

\item $2\hat{i}\times 3\hat{j}\cdot\hat{k}=$:
\begin{quote}\begin{multicols}{2}
\noindent (A) $-4$ \\
(B) $4$ \\
(C) $1$ \\
(D) $6$
\end{multicols}\end{quote}

\item $\hat{i}\cdot 2\hat{j}\times 3\hat{k}=$:
\begin{quote}\begin{multicols}{2}
\noindent (A) $0$ \\
(B) $6$ \\
(C) $-6$ \\
(D) None
\end{multicols}\end{quote}

\item $(\hat{i}-\hat{j})\cdot(\hat{j}-\hat{k})\times(\hat{k}-\hat{j})=$:
\begin{quote}\begin{multicols}{2}
\noindent (A) $1$ \\
(B) $0$ \\
(C) $\hat{i}-\hat{k}$ \\
(D) $\hat{j}-\hat{k}$
\end{multicols}\end{quote}

\item $[\hat{i}\ \hat{j}\ \hat{k}]$ equals:
\begin{quote}\begin{multicols}{2}
\noindent (A) $1$ \\
(B) $0$ \\
(C) $\hat{i}$ \\
(D) $\hat{j}$
\end{multicols}\end{quote}

\item $(\hat{i}\times\hat{k})\times\hat{j}$ is equal to:
\begin{quote}\begin{multicols}{2}
\noindent (A) $1$ \\
(B) $-\hat{j}$ \\
(C) $0$ \\
(D) $\hat{i}$
\end{multicols}\end{quote}

\item Work done by a constant force $\vec{F}$ during displacement $\vec{d}$ is equal to:
\begin{quote}\begin{multicols}{2}
\noindent (A) $\vec{F}\cdot\vec{d}$ \\
(B) $\vec{F}\times\vec{d}$ \\
(C) $\vec{r}\times\vec{F}$ \\
(D) $\vec{F}\times\vec{r}$
\end{multicols}\end{quote}

\item If $\vec{U}=\vec{V}$, then $\vec{U}\cdot(\vec{V}\times\vec{W})$:
\begin{quote}\begin{multicols}{2}
\noindent (A) $0$ \\
(B) $1$ \\
(C) $-1$ \\
(D) Cannot be calculated
\end{multicols}\end{quote}

\item What is the volume of the parallelepiped if its edges are $[2,-4,5]$, $[2,-3,6]$, $[0,-1,-1]$:
\begin{quote}\begin{multicols}{2}
\noindent (A) $16$ \\
(B) $0$ \\
(C) $5$ \\
(D) $24$
\end{multicols}\end{quote}

\item Moment of force $\vec{F}$ about a point with position vector $\vec{r}$ will be equal to:
\begin{quote}\begin{multicols}{2}
\noindent (A) $\vec{F}\times\vec{r}$ \\
(B) $\vec{r}\times\vec{F}$ \\
(C) $\vec{d}\times\vec{F}$ \\
(D) $\vec{F}\times\vec{d}$
\end{multicols}\end{quote}

\item The triple scalar product of vectors calculates the volume of:
\begin{quote}\begin{multicols}{2}
\noindent (A) Triangle \\
(B) Parallelogram \\
(C) Tetrahedron \\
(D) Parallelepiped
\end{multicols}\end{quote}

\item The work done by the force $2\hat{i}+4\hat{j}$ in moving a body from $A(0,0)$ to $B(3,5)$ is:
\begin{quote}\begin{multicols}{2}
\noindent (A) $6$ \\
(B) $12$ \\
(C) $13$ \\
(D) $26$
\end{multicols}\end{quote}

\item $[2\hat{k}\ \hat{j}\ \hat{i}]$:
\begin{quote}\begin{multicols}{2}
\noindent (A) $1$ \\
(B) $-1$ \\
(C) $-2$ \\
(D) $2$
\end{multicols}\end{quote}

\item $2\hat{j}\cdot(2\hat{k}\times\hat{i})$:
\begin{quote}\begin{multicols}{2}
\noindent (A) $4$ \\
(B) $3$ \\
(C) $2$ \\
(D) $1$
\end{multicols}\end{quote}

\item The moment of force $\vec{F}$ acting at $P$ about $C$ is:
\begin{quote}\begin{multicols}{2}
\noindent (A) $\vec{F}\times\overrightarrow{CP}$ \\
(B) $\overrightarrow{CP}\times\vec{F}$ \\
(C) $\overrightarrow{CP}\cdot\vec{F}$ \\
(D) $\overrightarrow{OP}\times\vec{F}$
\end{multicols}\end{quote}

\end{mcqlist}

\newpage

\section{Short Answer Questions (1--683)}

\section*{Write short answers of the following questions.}

\begin{shortlist}

\item Find the multiplicative inverse of the following complex number: $(-4,7)$.

\item Find the multiplicative inverse of the following complex number: $(\sqrt{2},-\sqrt{5})$.

\item Find the multiplicative inverse of the following complex number: $(1,0)$.

\item Separate into real and imaginary parts (write as a simple complex number): $\dfrac{2-7i}{4+5i}$.

\item Separate into real and imaginary parts (write as a simple complex number): $\dfrac{(-2+3i)^2}{1+i}$.

\item Separate into real and imaginary parts (write as a simple complex number): $\dfrac{i}{1+i}$.

\item Separate into real and imaginary parts (write as a simple complex number): $\dfrac{(4+3i)^2}{4-3i}$.

\item For $z\in\mathbb{C}$, show that: $|z|^2=z\cdot\bar{z}$.

\item If $z_1=2+7i$ and $z_2=-5+3i$, then evaluate the following: $|2z_1-4z_2|$.

\item If $z_1=2+7i$ and $z_2=-5+3i$, then evaluate the following: $|3z_1-2\bar{z_1}|$.

\item If $z_1=2+7i$ and $z_2=-5+3i$, then evaluate the following: $|-7z_2+2\bar{z_2}|$.

\item Find the least positive value of $n$, if $\left(\dfrac{1+i}{1-i}\right)^{2n}=1$.

\item Find the sum, difference and product of the complex numbers $(8,9)$ and $(5,-6)$.

\item If $z_1=(4,2)$ and $z_2=(3,-1)$, then find $\dfrac{z_1}{z_2}$.

\item If $z=\dfrac{(1+2i)^2}{2-i}$ then evaluate $|\bar{z}|$.

\item Find the real values of $x$ and $y$ in the following: $x+iy+2-3i=i(5-i)(3+4i)$.

\item Find the real values of $x$ and $y$ in the following: $(x+iy)(1-i)=(2-3i)(-5+5i)\left(-i\dfrac{3}{5}\right)$.

\item Find the real values of $x$ and $y$ in the following: $\dfrac{x}{2+i}+\dfrac{y}{3-i}=4+5i$.

\item If $z_1=-13+24i$ and $z_2=x+yi$, find the real values of $x$ and $y$ such that $z_1-z_2=-27+15i$.

\item Find the real values of $x$ and $y$ if: $(x+iy)^2=25+60i$.

\item Find the real values of $x$ and $y$ if: $(x+iy)^2=64+48i$.

\item Find the real values of $x$ and $y$ if: $(x+iy)^2=\dfrac{2i-3}{3+i}$.

\item If $z_1=2+3i$ and $z_2=1-\alpha$, find the value of $\alpha$ such that $\operatorname{Im}(z_1z_2)=7$.

\item Show that $\forall\, z_1,z_2\in\mathbb{C}$, $\overline{z_1 z_2}=\bar{z_1}\,\bar{z_2}$.

\item Find the square root of the following complex number: $-7-24i$.

\item Find the square root of the following complex number: $8-6i$.

\item Find the square root of the following complex number: $-15-36i$.

\item Find the square root of the following complex number: $119+120i$.

\item If $z_1=4+5i$ and $z_2=\alpha-2i$, find the real values of $\alpha$ such that $\operatorname{Re}(z_1z_2)=20$.

\item Find the square root of complex number $5+12i$ and also represent the square root on an Argand diagram.

\item Factorize the following: $a^2+4b^2$.

\item Factorize the following: $9a^2+16b^2$.

\item Factorize the following: $3x^2+3y^2$.

\item Factorize the following: $144x^2+225y^2$.

\item Factorize the following: $z^2-2z-1$.

\item Factorize the following: $z^2+6z+13$.

\item Factorize the following: $z^2+4z+5$.

\item Factorize the following: $9a^2+16b^2$.

\item Factorize the following polynomials into its linear factors: $z^3+8$.

\item Factorize the following polynomials into its linear factors: $z^3-2z^2+16z-32$.

\item Factorize the following polynomials into its linear factors: $z^4+21z^2-100$.

\item Factorize the following polynomials into its linear factors: $z^4-16$.

\item Factorize the following polynomials into its linear factors: $z^4+3z^2-4$.

\item Factorize the following polynomials into its linear factors: $z^4+5z^2+6$.

\item Factorize the following polynomials into its linear factors: $z^4-32z^2-3969$.

\item Solve the following complex quadratic equation by completing the square method: $2z^2-3z+4=0$.

\item Solve the following complex quadratic equation by completing the square method: $z^2+4z+13=0$.

\item Solve the following complex quadratic equation by completing the square method: $2z^2+6z+9=0$.

\item Solve the following complex quadratic equation by completing the square method: $3z^2-5z+7=0$.

\item Solve the following equation: $5z^5-5z=0$.

\item Solve the following equation: $z^3-5z^2+z-5=0$.

\item Solve the following equation: $z^3+z^2+z+1=0$.

\item Factorize the polynomial $P(z)=z^2+(i-3)z-3i$.

\item Factorize the polynomial $P(z)=z^2-4iz+12$.

\item Factorize the polynomial $P(z)=z^3+(1+i)z^2+iz$.

\item Find the three cube roots of: $8$.

\item Find the three cube roots of: $-8$.

\item Find the three cube roots of: $-27$.

\item Find the three cube roots of: $-125$.

\item Find the three cube roots of: $64$.

\item Evaluate: $\left(\dfrac{-1+\sqrt{-3}}{2}\right)^7+\left(\dfrac{-1-\sqrt{-3}}{2}\right)^7$.

\item Evaluate: $(-1+\sqrt{-3})^5+(-1-\sqrt{-3})^5$.

\item Evaluate $\displaystyle\sum_{k=0}^{5}\omega^{2k}$, where $\omega$ is an imaginary cube root of unity.

\item If $\omega$ is an imaginary cube root of unity, prove that $\dfrac{a+b\omega^2+c\omega}{a\omega^2+b\omega+c}=\omega$.

\item Prove that: $(x^3+y^3)=(x+y)(x+\omega y)(x+\omega^2 y)$.

\item Plot the following points: $(2,75)$.

\item Plot the following points: $(-3,120)$.

\item Plot the following points: $\left(2,\dfrac{\pi}{6}\right)$.

\item Plot the following points: $\left(5,\dfrac{5\pi}{6}\right)$.

\item Plot the following points: $\left(-\dfrac{5}{2},\dfrac{\pi}{3}\right)$.

\item Plot the following points: $\left(-3,-\dfrac{2\pi}{3}\right)$.

\item Plot the following points: $\left(-\dfrac{9}{2},-\dfrac{19\pi}{12}\right)$.

\item Plot the following points: $\left(-\dfrac{5}{2},\dfrac{5\pi}{12}\right)$.

\item Express the following complex number in polar form: $4+3i$.

\item Express the following complex number in polar form: $1+i$.

\item Express the following complex number in polar form: $\dfrac{1}{2}+\dfrac{\sqrt{3}}{2}i$.

\item Express the following complex number in polar form: $-\dfrac{5}{2}-\dfrac{5\sqrt{3}}{2}i$.

\item Express the following complex number in polar form: $\dfrac{1-i}{1+i}$.

\item Express the following complex number in polar form: $\dfrac{\sqrt{3}+i}{1+\sqrt{3}i}$.

\item Express the following complex number in polar form: $\dfrac{3+4i}{4+3i}$.

\item Convert the complex number $z$ in the rectangular form $x+iy$: $4\left(\cos\dfrac{5\pi}{3}+i\sin\dfrac{5\pi}{3}\right)$.

\item Convert the complex number $z$ in the rectangular form $x+iy$: $\dfrac{3}{2}\left(\cos\dfrac{7\pi}{6}+i\sin\dfrac{7\pi}{6}\right)$.

\item Convert the complex number $z$ in the rectangular form $x+iy$: $|z|=7,\ \arg(z)=\dfrac{23\pi}{3}$.

\item Convert the complex number $z$ in the rectangular form $x+iy$: $|z|=11,\ \arg(z)=-\dfrac{11\pi}{12}$.

\item Convert the complex number $z$ in the rectangular form $x+iy$: $|z|=\dfrac{10}{3},\ \arg(z)=-\dfrac{17\pi}{12}$.

\item Convert the complex number $z$ in the rectangular form $x+iy$: $2\cos(-33)+i2\sin(-33)$.

\item If $z_1=9\left(\cos\dfrac{5\pi}{4}+i\sin\dfrac{5\pi}{4}\right)$ and $z_2=5\left(\cos\dfrac{\pi}{3}+i\sin\dfrac{\pi}{3}\right)$ then find: $z_1+z_2$.

\item If $z_1=9\left(\cos\dfrac{5\pi}{4}+i\sin\dfrac{5\pi}{4}\right)$ and $z_2=5\left(\cos\dfrac{\pi}{3}+i\sin\dfrac{\pi}{3}\right)$ then find: $z_1-z_2$.

\item If $z_1=9\left(\cos\dfrac{5\pi}{4}+i\sin\dfrac{5\pi}{4}\right)$ and $z_2=5\left(\cos\dfrac{\pi}{3}+i\sin\dfrac{\pi}{3}\right)$ then find: $z_1\cdot z_2$.

\item If $z_1=9\left(\cos\dfrac{5\pi}{4}+i\sin\dfrac{5\pi}{4}\right)$ and $z_2=5\left(\cos\dfrac{\pi}{3}+i\sin\dfrac{\pi}{3}\right)$ then find: $\dfrac{z_1}{z_2}$.

\item If $z_1=7\left(\cos\dfrac{23\pi}{12}+i\sin\dfrac{23\pi}{12}\right)$ and $z_2=11\left(\cos\dfrac{11\pi}{2}+i\sin\dfrac{11\pi}{2}\right)$ then find the following and express the result into $x+iy$ form: $z_1+z_2$.

\item If $z_1=7\left(\cos\dfrac{23\pi}{12}+i\sin\dfrac{23\pi}{12}\right)$ and $z_2=11\left(\cos\dfrac{11\pi}{2}+i\sin\dfrac{11\pi}{2}\right)$ then find the following and express the result into $x+iy$ form: $z_1-z_2$.

\item If $z_1=7\left(\cos\dfrac{23\pi}{12}+i\sin\dfrac{23\pi}{12}\right)$ and $z_2=11\left(\cos\dfrac{11\pi}{2}+i\sin\dfrac{11\pi}{2}\right)$ then find the following and express the result into $x+iy$ form: $z_1\cdot z_2$.

\item If $z_1=7\left(\cos\dfrac{23\pi}{12}+i\sin\dfrac{23\pi}{12}\right)$ and $z_2=11\left(\cos\dfrac{11\pi}{2}+i\sin\dfrac{11\pi}{2}\right)$ then find the following and express the result into $x+iy$ form: $\dfrac{z_1}{z_2}$.

\item Divide $z_1=6(\cos 150^\circ+i\sin 150^\circ)$ by $z_2=3(\cos 30^\circ+i\sin 30^\circ)$ and express in $x+iy$ form.

\item Multiply $z_1=2(\cos 60^\circ+i\sin 60^\circ)$ and $z_2=5(\cos 90^\circ+i\sin 90^\circ)$ and express in $x+iy$ form.

\item Find the modulus and argument of $z=-2-2i$.

\item Write the equation $\arg(\bar{z}-2+i)=\dfrac{2\pi}{3}$ in cartesian form, if $z=x+iy$.

\item Solve the equation $|z-2i|=|\bar{z}+2|$ for $z=x+iy$.

\item For $z=x+iy$, solve the equation $|5z+4+i|=|5\bar{z}-3+2i|$.

\item Express the complex number $1+i\sqrt{3}$ in polar form.

\item Given that: (a) $f(x)=x^2-1$ (b) $f(x)=\sqrt{2x+3}$, find (i) $f(-3)$, (ii) $f(0)$.

\item Given that: (a) $f(x)=x^2-1$ (b) $f(x)=\sqrt{2x+3}$, find (i) $f(x-2)$, (ii) $f(x^2+3)$.

\item Find $\dfrac{f(a+h)-f(a)}{h}$ and simplify where: $f(x)=4x+7$.

\item Find $\dfrac{f(a+h)-f(a)}{h}$ and simplify where: $f(x)=\sin x$.

\item Find $\dfrac{f(a+h)-f(a)}{h}$ and simplify where: $f(x)=x^3+x^2-1$.

\item Find $\dfrac{f(a+h)-f(a)}{h}$ and simplify where: $f(x)=\tan x$.

\item Express the following: The area $A$ of a square as a function of its perimeter $P$.

\item Express the following: The circumference $C$ of a circle as a function of its area $A$.

\item Express the following: The surface area $S$ of a cube as a function of its volume $V$.

\item Find the domain and the range of the function $g$ defined below: $g(x)=5-x$.

\item Find the domain and the range of the function $g$ defined below: $g(x)=\sqrt{x+2}$.

\item Find the domain and the range of the function $g$ defined below: $g(x)=|x-5|$.

\item Find the domain and the range of the function $g$ defined below: $g(x)=\dfrac{x+2}{3-x}$.

\item Given $f(x)=x^3-ax^2+bx+1$. If $f(2)=-3$ and $f(-1)=0$, find the value of $a$ and $b$.

\item Consider the function $f(x)=3x-5$. Determine the domain and range of $f(x)$.

\item Consider the function $f(x)=3x-5$. Is the function $f$ one-to-one? Justify your answer.

\item Consider the function $f(x)=3x-5$. Is the function $f$ onto if the co-domain is all real numbers? Explain.

\item Given $f(x)=x^3-2x^2+4x-1$ find: $f(0)$.

\item Given $f(x)=x^3-2x^2+4x-1$ find: $f(1)$.

\item Given $f(x)=x^3-2x^2+4x-1$ find: $f(-2)$.

\item Given $f(x)=x^3-2x^2+4x-1$ find: $f(1+x)$.

\item Given $f(x)=x^3-2x^2+4x-1$ find: $f\left(\dfrac{1}{x}\right),\ x\neq 0$.

\item Find the domain and range of $f(x)=x^2$.

\item Find the domain and range of $f(x)=\dfrac{x}{x^2-4}$.

\item Find the domain and range of $f(x)=\sqrt{x^2-9}$.

\item Find the point of intersection of the coordinate axes and the following linear function graphically: $y=-5x+10$.

\item Find the point of intersection of the coordinate axes and the following linear function graphically: $y=2x-1$.

\item Find the point of intersection of the coordinate axes and the following linear function graphically: $y=\dfrac{1}{2}x-3$.

\item Find the point of intersection of the coordinate axes and the following linear function graphically: $y=3x+\dfrac{3}{2}$.

\item Find the point(s) of intersection of the following function graphically: $f(x)=-3x-4,\ g(x)=\dfrac{1}{2}x+3$.

\item Graph the following function: $y=\sqrt{3x}$.

\item Graph the following function: $y=\sqrt{x+5}$.

\item Graph the following function: $y=-\dfrac{1}{2}\sqrt{x}$.

\item Find the point of intersection of a linear function $y=2x+6$ and coordinate axes graphically.

\item Find the point of intersection of $y=3x+2$ and $y=-x+6$ graphically.

\item Find the maximum and minimum value of the following quadratic function by completing squares: $f(x)=x^2+6x+13$.

\item Find the maximum and minimum value of the following quadratic function by completing squares: $f(x)=x^2+4x$.

\item Find the maximum and minimum value of the following quadratic function by completing squares: $f(x)=-x^2+8x+13$.

\item Find the maximum and minimum value of the following quadratic function by completing squares: $f(x)=-x^2-3x-5$.

\item Find the maximum and minimum value of the following quadratic function by completing squares: $f(x)=3x^2+6x-13$.

\item Solve the following absolute value quadratic equation and inequalities: $|x^2+1|=5$.

\item Solve the following absolute value quadratic equation and inequalities: $|x^2+5x+4|=0$.

\item Solve the following absolute value quadratic equation and inequalities: $|x^2-6x+8|=4$.

\item Solve the following absolute value quadratic equation and inequalities: $|x^2-4|<5$.

\item Find the maximum and minimum value of $f(x)=-2x^2+4x+3$ by completing the square.

\item Find the maximum and minimum value of $f(x)=x^2-2x-3$.

\item Solve: $|x^2-4|=5$.

\item Solve the following: $\dfrac{1}{3x}+\dfrac{4x}{6}=1,\ x\neq 0$.

\item Solve the following: $\dfrac{x}{x+1}+\dfrac{x+1}{x}=\dfrac{5}{2};\ x\neq -1,0$.

\item Solve the following: $\sqrt{x+5}-\sqrt{x-3}=2$.

\item The length of a room is 3 metres greater than its breadth. If the area of the room is 180 square metres, find the length and the breadth of the room.

\item Solve $\sqrt{x+8}+\sqrt{x+3}=\sqrt{12x+13}$.

\item Solve $\dfrac{1}{x}+\dfrac{2}{x+1}=1,\ x\neq 0,\ x\neq -1$.

\item If $A=[a_{ij}]_{3\times 4}$, then show that $I_3A=A$.

\item If $A=[a_{ij}]_{3\times 4}$, then show that $AI_4=A$.

\item If $A=\begin{pmatrix}0 & -1 & 2\\ 3 & 2 & 1\\ -1 & 0 & 4\end{pmatrix}$, $B=\begin{pmatrix}2 & 1 & -1\\ 1 & 2 & 4\\ -1 & 2 & 1\end{pmatrix}$ and $C=\begin{pmatrix}1 & 0 & -2\\ -1 & 5 & 0\\ 3 & 4 & -1\end{pmatrix}$, then find: $A-B$.

\item If $A=\begin{pmatrix}0 & -1 & 2\\ 3 & 2 & 1\\ -1 & 0 & 4\end{pmatrix}$, $B=\begin{pmatrix}2 & 1 & -1\\ 1 & 2 & 4\\ -1 & 2 & 1\end{pmatrix}$ and $C=\begin{pmatrix}1 & 0 & -2\\ -1 & 5 & 0\\ 3 & 4 & -1\end{pmatrix}$, then find: $B-C$.

\item If $A=\begin{pmatrix}0 & -1 & 2\\ 3 & 2 & 1\\ -1 & 0 & 4\end{pmatrix}$, $B=\begin{pmatrix}2 & 1 & -1\\ 1 & 2 & 4\\ -1 & 2 & 1\end{pmatrix}$ and $C=\begin{pmatrix}1 & 0 & -2\\ -1 & 5 & 0\\ 3 & 4 & -1\end{pmatrix}$, then find: $(A-B)-C$.

\item If $A=\begin{pmatrix}0 & -1 & 2\\ 3 & 2 & 1\\ -1 & 0 & 4\end{pmatrix}$, $B=\begin{pmatrix}2 & 1 & -1\\ 1 & 2 & 4\\ -1 & 2 & 1\end{pmatrix}$ and $C=\begin{pmatrix}1 & 0 & -2\\ -1 & 5 & 0\\ 3 & 4 & -1\end{pmatrix}$, then find: $A-(B-C)$.

\item If $A$ and $B$ are square matrices of the same order, then explain why in general $(A+B)^2\neq A^2+2AB+B^2$.

\item If $A$ and $B$ are square matrices of the same order, then explain why in general $(A-B)^2\neq A^2-2AB+B^2$.

\item If $A=\begin{pmatrix}-1 & 2 & 3 & 0\\ 1 & 0 & 2 & -2\\ -3 & 5 & 3 & -1\end{pmatrix}$ then find $AA^t$, $A^tA$ and $(A^t)^t$.

\item Solve the following matrix equation for $X$: $2X-3A=B$ if $A=\begin{pmatrix}2 & 3 & -2\\ -1 & 1 & 5\end{pmatrix}$ and $B=\begin{pmatrix}2 & -3 & 1\\ 5 & 4 & -1\end{pmatrix}$.

\item Find $AB$ and $BA$ if $A=\begin{pmatrix}2 & 0 & 1\\ 1 & 4 & 2\\ 3 & 0 & 6\end{pmatrix}$ and $B=\begin{pmatrix}1 & -1 & 0\\ 2 & 3 & -1\\ 1 & -2 & 3\end{pmatrix}$.

\item Evaluate the following determinant: $\begin{vmatrix}1 & -2 & -4\\ 3 & -1 & -3\\ -2 & 3 & 2\end{vmatrix}$.

\item Evaluate the following determinant: $\begin{vmatrix}a+b & a-b & a\\ a & a+b & a-b\\ a-b & a & a+b\end{vmatrix}$.

\item Evaluate the following determinant: $\begin{vmatrix}2x & x & x\\ y & 2y & y\\ z & z & 2z\end{vmatrix}$.

\item Evaluate the following determinant: $\begin{vmatrix}5 & 2 & -3\\ 3 & 10 & -1\\ -2 & 1 & -2\end{vmatrix}$.

\item Evaluate the following determinant: $\begin{vmatrix}1 & 2 & -3\\ -1 & 3 & 5\\ -2 & 5 & 6\end{vmatrix}$.

\item Evaluate the following determinant: $\begin{vmatrix}1 & 20 & -2\\ -1 & -1 & -3\\ 2 & 4 & -1\end{vmatrix}$.

\item Without expansion show that: $\begin{vmatrix}5 & 6 & -1\\ 2 & 2 & 0\\ 2 & -8 & 10\end{vmatrix}=0$.

\item Without expansion show that: $\begin{vmatrix}-a & 0 & b\\ 0 & a & -c\\ c & -b & 0\end{vmatrix}=0$.

\item Without expansion show that: $\begin{vmatrix}l & m+n & 1\\ m & n+l & 1\\ n & l+m & 1\end{vmatrix}=0$.

\item Without expansion show that: $\begin{vmatrix}2 & 1 & 3x\\ 2 & 3 & 9x\\ 3 & 5 & 15x\end{vmatrix}=0$.

\item Without expansion show that: $\begin{vmatrix}7 & 8 & 9\\ 5 & 6 & 7\\ 2 & 3 & 4\end{vmatrix}=0$.

\item Without expansion show that: $\begin{vmatrix}yz & zx & xy\\ x & y & z\\ \frac{1}{x} & \frac{1}{y} & \frac{1}{z}\end{vmatrix}=0$.

\item Find the value of $x$ if: $\begin{vmatrix}1 & 1 & 1\\ 2 & x & 2\\ 3 & 6 & x\end{vmatrix}=0$.

\item Find $|AA^t|$ and $|A^tA|$: $A=\begin{pmatrix}-3 & 2 & -1\\ 2 & 1 & 3\end{pmatrix}$.

\item Evaluate the determinant if $A=\begin{pmatrix}1 & -2 & 3\\ -2 & 3 & 1\\ 4 & -3 & 2\end{pmatrix}$.

\item Find the cofactor $A_{12},A_{22}$ and $A_{32}$ of $A=\begin{pmatrix}1 & -2 & 3\\ -2 & 3 & 1\\ 4 & -3 & 2\end{pmatrix}$ and find $|A|$.

\item Without expansion, show that $\begin{vmatrix}x & a+x & b+c\\ x & b+x & c+a\\ x & c+x & a+b\end{vmatrix}=0$.

\item Find $A^{-1}$ if $A=\begin{pmatrix}1 & 0 & 2\\ 0 & 2 & 1\\ 1 & -1 & 1\end{pmatrix}$.

\item Find the inverse of the matrix $A=\begin{pmatrix}2 & 5 & -1\\ 3 & 4 & 2\\ 1 & 2 & -2\end{pmatrix}$.

\item Resolve the following into partial fractions: $\dfrac{2}{x^2-1}$.

\item Resolve $\dfrac{7x+25}{(x+3)(x+4)}$ into partial fractions.

\item Find the next four terms of the following sequence: $12,16,20,\ldots$

\item Find the next four terms of the following sequence: $3,1,-1,\ldots$

\item Write down the first three terms of the following sequence: $a_n=3n+5$.

\item Write down the first three terms of the following sequence: $a_{n+1}=4a_n-7$ and $a_1=3$.

\item Write down the first three terms of the following sequence: $a_n=(n-3)(n+1)$.

\item Write down the first three terms of the following sequence: $a_1=-1,\ a_{n+1}=\dfrac{3}{a_n+2}$.

\item Write down the first three terms of the following sequence: $a_n=8-\dfrac{20}{3+n}$.

\item Write down the first three terms of the following sequence: $a_1=1,\ a_{n+1}=(3a_n+2)^2$.

\item Write down the first three terms of the following sequence: $a_n=(-2n)^2$.

\item Write down the first three terms of the following sequence: $a_n=(-1)^n 7n^2$.

\item An expression for the $n^{\text{th}}$ triangular number is $\dfrac{n(n+1)}{2}$. Write down the $15^{\text{th}}$ triangular number. Make a triangle of dots by taking $n=5$.

\item Write down the $n^{\text{th}}$ term of the following sequence: $1,4,9,\ldots$

\item Write down the $n^{\text{th}}$ term of the following sequence: $1,\ 1+2,\ 1+2+3,\ldots$

\item Write down the $n^{\text{th}}$ term of the following sequence: $a_1b_1,\ a_2b_2,\ a_3b_3,\ldots$

\item Write down the $n^{\text{th}}$ term of the following sequence: $x,\ 2x^2,\ 3x^3,\ldots$

\item Write down the $n^{\text{th}}$ term of the following sequence: $a_1,\ a_1+d,\ a_1+2d,\ldots$

\item Write down the $n^{\text{th}}$ term of the following sequence: $a_1,\ a_1r,\ a_1r^2,\ldots$

\item Write down the $n^{\text{th}}$ term of the following sequence: $\dfrac{a_1}{b_1+c_1},\ \dfrac{2a_2}{b_2+c_2},\ \dfrac{3a_3}{b_3+c_3},\ldots$

\item Write down the $n^{\text{th}}$ term of the following sequence: $\dfrac{1}{a_1},\ \dfrac{1}{a_1+d},\ \dfrac{1}{a_1+2d},\ldots$

\item Find the first four terms of the following sequence whose $n^{\text{th}}$ term is given: $a_n=3n+1$.

\item Find the first four terms of the following sequence whose $n^{\text{th}}$ term is given: $a_n=3n^2-3$.

\item Find the common difference and write the next two terms of the following sequence: $9,16,23,\ldots$

\item Find the common difference and write the next two terms of the following sequence: $5,\ 5+\sqrt{2},\ 5+2\sqrt{2},\ldots$

\item Write the first three terms of the following arithmetic sequence, with given information: $a_1=2,\ d=13$.

\item Write the first three terms of the following arithmetic sequence, with given information: $a_1=12,\ d=-13$.

\item Find $a_{n+1}$ and $a_n$ if $a_n=4+3n$.

\item Find the indicated term of the following arithmetic sequence: $a_1=3,\ d=7$, $a_{14}$.

\item Find the indicated term of the following arithmetic sequence: $8,3,-2,\ldots$, $a_{12}$.

\item Is 301 a term of the A.P $5,11,17,\ldots$?

\item If $2x,\ x+8,\ 3x+1$ are in A.P., then find the value of $x$.

\item Which term of the A.P $3,8,13,\ldots$ is $123$?

\item Which term of the A.P $30,29.5,29,\ldots$ is the first negative term?

\item How many numbers of three digits are divisible by 7?

\item Find the $8^{\text{th}}$ term from the end of the A.P $8,11,14,\ldots,185$.

\item If the $p^{\text{th}}$ term of an A.P. is $q$ and the $q^{\text{th}}$ term is $p$, prove that its $n^{\text{th}}$ term is $(p+q-n)$.

\item If $\dfrac{1}{a},\dfrac{1}{b}$ and $\dfrac{1}{c}$ are in A.P, show that $b=\dfrac{2ac}{a+c}$.

\item If the $5^{\text{th}}$ term of an A.P is $13$ and the $17^{\text{th}}$ term is $49$, find $a_n$ and $a_{13}$.

\item Find the number of terms in the A.P; if $a_1=3$, $d=7$ and $a_n=59$.

\item If $a_n=3n-11$ find the $n^{\text{th}}$ term of the sequence.

\item Find the A.M between the given numbers: $2+\sqrt{3}i,\ 2-\sqrt{3}i$.

\item Find the A.M between the given numbers: $(a+b)^2,\ (a-b)^2$.

\item If $6,11,16$ are three A.Ms between $a$ and $b$, find $a$ and $b$.

\item Find three A.Ms between $\sqrt{2}$ and $3\sqrt{2}$.

\item Sum the series: $3+6+9+\cdots a_{20}$.

\item Sum the series: $\dfrac{4}{\sqrt{5}}+\sqrt{5}+\dfrac{6}{\sqrt{5}}+a_n$.

\item Find $S_n$ for the following arithmetic series: $a_1=4,\ n=25,\ a_n=100$.

\item Find $S_n$ for the following arithmetic series: $a_1=40,\ n=20,\ d=-3$.

\item Find $S_n$ for the following arithmetic series: $a_n=52,\ n=21,\ d=-4$.

\item Find $a_1$ for the arithmetic series $d=8$, $n=19$, $S_n=1786$.

\item How many items of the series: $96+93+90+\cdots$ amount to $1071$?

\item If the three sides of a right-angled triangle having perimeter 36 cm are in A.P., find them.

\item Sum the series: $3+5-7+9+11-13+15+17-19+\cdots$ to $3n$ terms.

\item Sum the series: $1+4-7+10+13-16+19+22-25+\cdots$ to $3n$ terms.

\item Find the sum to 20 terms of the series whose $r^{\text{th}}$ term is $3r+1$.

\item Find the sum of the first 100 positive integers.

\item Find the arithmetic series if its fifth term is $19$ and $S_4=a_9+1$.

\item How many terms of the series $-9-6-3+0+\cdots$ amount to $66$?

\item Find the first three terms of an arithmetic series in which $a_1=9$, $a_n=105$ and $S_n=741$.

\item Find the $6^{\text{th}}$ term of the G.P: $-6,-3,\dfrac{-3}{2},\ldots$

\item Find the $8^{\text{th}}$ term of the sequence, $3,3^2,3^3,\ldots$

\item The $n^{\text{th}}$ term of the sequences $1,2,4,8,\ldots$ and $256,128,64,\ldots$ are equal. Find the value of $n$.

\item Find the first five terms of the following sequence described: $a_1=243,\ r=\dfrac{1}{3}$.

\item Find the $12^{\text{th}}$ term of $1+i,\ 2i,-2+2i,\ldots$

\item If $a,b,c,d$ are in G.P., prove that: $a-b,\ b-c,\ c-d$ are in G.P.

\item If $a,b,c,d$ are in G.P., prove that: $a^2-b^2,\ b^2-c^2,\ c^2-d^2$ are in G.P.

\item If $a,b,c,d$ are in G.P., prove that: $a^2+b^2,\ b^2+c^2,\ c^2+d^2$ are in G.P.

\item If $a,\ 2a+2,\ 3a+3,\ldots$ are in G.P then find the fifth term.

\item Find the eighth term of a geometric sequence for which $a_1=-3$ and $r=-2$.

\item Find the $n^{\text{th}}$ term of the G.P., $3,12,48,\ldots$

\item Find the tenth term of the G.P., for which $a_4=108$ and $r=3$.

\item Find the $5^{\text{th}}$ term of the G.P., $3,\ldots,12,\ldots$

\item Find $a_1$ if $a_4=\dfrac{8}{27}$ and $a_7=\dfrac{-64}{729}$ of a G.P.

\item Find G.M. between: $-2$ and $8$.

\item Find G.M. between: $-2i$ and $8i$.

\item Find G.M. between: $6$ and $9$.

\item Insert four real geometric means between $3$ and $96$.

\item Insert three G.Ms. between $2$ and $\dfrac{1}{2}$.

\item Find the sum of first 15 terms of the G.P $1,\dfrac{1}{3},\dfrac{1}{9},\ldots$

\item Sum of $n$ terms of the series: $0.2+0.22+0.222+\cdots$

\item Sum of $n$ terms of the series: $3+33+333+\cdots$

\item Sum of $n$ terms of the series: $r+(1+h)r^2+(1+k+k^2)r^3+\cdots$

\item Sum the series $2+(1-i)+\left(\dfrac{1}{i}\right)+\cdots$ to 8 terms.

\item Find the sum of $n$ terms of the geometric series if $a_n=(-3)\left(\dfrac{2}{5}\right)^n$.

\item Sum of $n$ terms of the following series: $1\cdot 2+3\cdot 4+5\cdot 8+7\cdot 16+\cdots$

\item Sum of $n$ terms of the following series: $2\cdot 3+4\cdot 3^2+6\cdot 3^3+8\cdot 3^4+\cdots$

\item Show that $\sqrt[4]{4}\cdot\sqrt[8]{16}\cdot\sqrt[16]{64}\cdot\sqrt[32]{256}\cdots\infty=16$.

\item Sum of $n$ terms of the series $2+4x+6x^2+8x^3+\cdots,\ x\neq 1$.

\item Sum the series upto $n$ terms: $2\cdot 1+3\cdot 2+4\cdot 4+5\cdot 8+\cdots$

\item Sum the series upto $n$ terms: $2+\dfrac{4}{3}+\dfrac{6}{9}+\dfrac{8}{27}+\cdots$

\item Find the $9^{\text{th}}$ term of the following harmonic sequence: $\dfrac{1}{3},\dfrac{1}{5},\dfrac{1}{7},\ldots$

\item Find the $9^{\text{th}}$ term of the following harmonic sequence: $\dfrac{-1}{5},\dfrac{-1}{3},-1,\ldots$

\item Insert five harmonic means between the following given numbers: $\dfrac{-2}{5}$ and $\dfrac{2}{13}$.

\item Insert five harmonic means between the following given numbers: $\dfrac{1}{4}$ and $\dfrac{1}{24}$.

\item The first term of an H.P is $\dfrac{1}{3}$ and the fifth term is $\dfrac{1}{5}$. Find the $9^{\text{th}}$ term.

\item If 5 is the harmonic mean between 2 and $b$, find $b$.

\item If $a^2,b^2$ and $c^2$ are in A.P., show that $a+b+c$ and $b+c$ are in H.P.

\item If $a,b,c,d$ are in H.P., show that $3(a-b)(c-d)=(b-c)(a-d)$.

\item If $a,b,c$ are in G.P., show that $\log_a x,\ \log_b x,\ \log_c x$ are in H.P.

\item If $a,b,c$ are in H.P., show that $\dfrac{a-b}{b-c}=\dfrac{a}{c}$.

\item If $a,b,c$ are in H.P., show that $(ac)^2=(a+c)(a-2b+c)$.

\item If $2+x,\ 5+x$ and $9+x$ are in H.P., find the value of $x$.

\item Find the $n^{\text{th}}$ and $8^{\text{th}}$ terms of H.P.: $\dfrac{1}{2},\dfrac{1}{5},\dfrac{1}{8},\ldots$

\item Find three harmonic means between $\dfrac{1}{5}$ and $\dfrac{1}{17}$.

\item Sum the following series upto $n$ terms: $1\times 3+2\times 5+3\times 7+\cdots$

\item Sum the following series upto $n$ terms: $1\times 5+2\times 8+3\times 11+\cdots$

\item Sum the following series upto $n$ terms: $1\times 2+2\times 5+3\times 8+\cdots$

\item Sum the following series upto $n$ terms: $1\times 3\times 5+2\times 4\times 6+3\times 5\times 7+\cdots$

\item Sum the following series upto $n$ terms: $1\times 2\times 4+2\times 3\times 7+3\times 4\times 10+\cdots$

\item Sum the following series upto $n$ terms: $2^2+4^2+6^2+\cdots$

\item Sum the following series upto $n$ terms: $3^2+6^2+9^2+\cdots$

\item Sum the following series upto $n$ terms: $4\times 1^2+7\times 2^2+10\times 3^2+\cdots$

\item Sum the following series upto $n$ terms: $3+(3+7)+(3+7+11)+\cdots$

\item Sum the following series upto $n$ terms: $1^2+(1^2+2^2)+(1^2+2^2+3^2)+\cdots$

\item Sum the series: $1^2-2^2+3^2-4^2+\cdots+(2n-1)^2-(2n)^2$.

\item Find the sum to $n$ terms of the series whose $n^{\text{th}}$ term is given: $5n^2+2n+3$.

\item Find the sum of the series $1^3+3^3+5^3+\cdots$ to $n$ terms.

\item Find the sum of $n$ terms of the series whose $n^{\text{th}}$ term is $n^3+\dfrac{3}{2}n^2+\dfrac{1}{2}n+1$.

\item Amna invested Rs.\ 200000 at 5\% interest compounded annually. What total amount will she get after 10 years.

\item To buy furniture for a new apartment, Tayyab borrowed Rs.\ 50{,}000 at 8\% simple interest for 11 years. How much interest will he pay?

\item Evaluate the following: $\dfrac{10!}{0!8!}$.

\item Evaluate the following: $\dfrac{12!}{3!(12-3)!}$.

\item Evaluate the following: $\dfrac{1440}{3!4!}+\dfrac{2400}{5!2!}$.

\item Evaluate the following: $\dfrac{(n+2)!}{(n+a)!}$.

\item Write the following in factorial form: $n^3-n$.

\item Write the following in factorial form: $n(n-1)(n-2)\cdots(n-r+1)$.

\item Find $n$, if $(n+4)!=3024\cdot n!$.

\item If $\dfrac{1}{7!}+\dfrac{1}{8!}=\dfrac{x}{9!}$, find $x$.

\item Evaluate $\dfrac{8!}{6!}$.

\item Write $8\cdot 7\cdot 6\cdot 5$ in the factorial form.

\item Evaluate $\dfrac{9!}{6!3!}$.

\item Evaluate the following: ${}^{10}P_5$.

\item Evaluate the following: ${}^5P_2$.

\item Evaluate the following: ${}^7P_7$.

\item Evaluate the following: ${}^{10}P_3$.

\item Find the value of $n$ when: ${}^nP_3=504$.

\item Find the value of $n$ when: ${}^{15}P_n=15\cdot 14\cdot 13\cdot 12\cdot 11$.

\item Find the value of $n$ when: ${}^nP_5:{}^{n-2}P_2=540:1$.

\item Prove from the first principle that: ${}^nP_r=n\cdot{}^{n-1}P_{r-1}$.

\item How many words can be formed from the letters of the following word using all letters when no letter is to be repeated: NETWORK.

\item How many words can be formed from the letters of the following word using all letters when no letter is to be repeated: PYTHON.

\item How many words can be formed from the letters of the following word using all letters when no letter is to be repeated: COMPUTER.

\item How many signals can be given by 5 flags of different colours, when any number of flags are used at a time?

\item How many 4-digit numbers can be formed, with distinct digits, with each digit odd?

\item How many 5-digit multiples of 5 can be formed from the digits $2,3,5,7,9$, when no digit is repeated?

\item How many different 4-digit numbers can be formed out from the digits $1,2,3,4,5,6$, when no digit is repeated?

\item How many arrangements of the letters of the following word, taken all together, can be made? PAKISTAN.

\item How many arrangements of the letters of the following word, taken all together, can be made? CURRICULUM.

\item How many arrangements of the letters of the following word, taken all together, can be made? PROBABILITY.

\item How many permutations of the letters of the word ``BANANA'' can be made, if B must be the first letter in each arrangement?

\item How many arrangements of the letters of the word TRIGONOMETRY can be made, if each arrangement begins with T and ends with Y?

\item In how many ways can 5 couples sit around a round table if no two women are sitting together?

\item How many 6-digit numbers can be formed from the digits $7,7,8,8,9,9$?

\item In how many ways can 8 keys be arranged in a circular key ring?

\item How many necklaces can be made from 10 beads of different colours?

\item In how many ways can the letters of the word MISSISSIPPS be arranged when all the letters are to be used?

\item Evaluate the following: ${}^{50}C_{50}$.

\item Evaluate the following: ${}^{1000}C_0$.

\item Evaluate the following: ${}^{10}C_7$.

\item Evaluate the following: ${}^{20}C_{17}$.

\item If ${}^{3n}C_2:{}^nC_2=15:1$, find $n$.

\item If ${}^nP_r=120$ and ${}^nC_r=20$, find $r$.

\item Find the value of $n$ and $r$, when: ${}^nC_r=56$, ${}^nP_r=336$.

\item Find the value of $n$ and $r$, when: ${}^{n-1}C_{r-1}:{}^nC_r:{}^nC=1:3:7$.

\item Prove that the product of $r$ consecutive integers is divisible by $r!$.

\item Find the number of possible arrangements of vowel letters from the English alphabet.

\item How many diagonals and triangles can be formed by joining the vertices of the polygon having 15 sides?

\item Find the number of sides of a polygon if the number of its diagonals is 104.

\item How many triangles can be formed by joining 15 points, 6 of which lie on the same straight line?

\item Zain has 8 different fruits. He wants to select 5 fruits out of 8 fruits to make a fruit chat. How many combinations of fruits can he select?

\item Find the number of diagonals of a 6-sided figure.

\item Use mathematical induction to prove the following formula for every positive integer $n$: $\log x^n=n\log x$, where $x$ is positive.

\item Use mathematical induction to prove the following formula for every positive integer $n$: $a+(a+d)+(a+2d)+\cdots+[a+(n-1)d]=\dfrac{n}{2}[2a+(n-1)d]$.

\item Use mathematical induction to prove the following formula for every positive integer $n$: the sum of the first $n$ odd natural numbers is $n^2$.

\item Prove by mathematical induction that for all positive integral values of $n$: $n^2+n$ is divisible by 2.

\item Prove by mathematical induction that for all positive integral values of $n$: $5^n-2^n$ is divisible by 3.

\item $x-y$ is a factor of $x^n-y^n$ for all positive integral values of $n$.

\item $n!>2^n-1$ for integral values of $n\geq 4$.

\item $4^n>3^n+2^{n-1}$ for integral values of $n\geq 2$.

\item $1+nx\leq(1+x)^n$ for $n\geq 2$ and $x>-1$.

\item Use mathematical induction to prove that $3+6+9+\cdots+3n=\dfrac{3n(n+1)}{2}$ for every positive integer $n$.

\item Show that $\dfrac{n^3+2n}{3}$ represents an integer $\forall n\in\mathbb{N}$.

\item Prove that $4^n+6n-1$ is divisible by 9 for all $n\in\mathbb{N}$.

\item Show that $1+3+5+\cdots+(2n+5)=(n+3)^2$ for integral values of $n\geq -2$.

\item Show that the inequality $4^n>4n+4$ is true, for integral values of $n\geq 2$.

\item If a recursively-defined sequence with $n>1$ is given, show that the last digit of $a_n$ is 7.

\item Using the binomial theorem, expand the following: $\left(\dfrac{x}{2}-\dfrac{2}{x^2}\right)^6$.

\item Using the binomial theorem, expand the following: $\left(2a-\dfrac{x}{a}\right)^7$.

\item Using the binomial theorem, expand the following: $\left(\sqrt{\dfrac{a}{x}}-\sqrt{\dfrac{x}{a}}\right)^6$.

\item Calculate the following by means of the binomial theorem: $(0.97)^3$.

\item Calculate the following by means of the binomial theorem: $(2.02)^4$.

\item Calculate the following by means of the binomial theorem: $(9.98)^4$.

\item Calculate the following by means of the binomial theorem: $(2.1)^5$.

\item Expand and simplify the following: $(a+\sqrt{2}x)^4+(a-\sqrt{2}x)^4$.

\item Expand and simplify the following: $(2+\sqrt{3})^5+(2-\sqrt{3})^5$.

\item Expand the following in ascending powers of $x$: $(2+x-x^2)^4$.

\item Expand the following in ascending powers of $x$: $(1-x+x^2)^4$.

\item Find the term involving $x^4$ in the expansion of $(3-2x)^7$.

\item Find the term involving $x^{-2}$ in the expansion of $\left(x-\dfrac{2}{x^2}\right)^{13}$.

\item Find the term involving $y^3$ in the expansion of $(x-\sqrt{y})^{11}$.

\item Find the coefficient of $x^5$ in the expansion of $\left(x^2-\dfrac{3}{2x}\right)^{10}$.

\item Find the coefficient of $x^3$ in the expansion of $\left(x^2-\dfrac{1}{x}\right)^{2n}$.

\item Find the $6^{\text{th}}$ term in the expansion of $\left(x^2-\dfrac{3}{2x}\right)^{10}$.

\item Find the term independent of $x$ in the following expansion: $\left(x-\dfrac{2}{x}\right)^{10}$.

\item Expand $\left(\dfrac{a}{2}-\dfrac{2}{a}\right)^6$ and also find its general term.

\item Evaluate $(9.9)^5$ using the binomial theorem.

\item Find the specified term in the expansion of $\left(\dfrac{3}{2}x-\dfrac{1}{3x}\right)^{11}$.

\item Find the following in the expansion of $\left(\dfrac{x}{2}+\dfrac{2}{x^2}\right)^{12}$.

\item Expand the following upto 4 terms, taking the values of $x$ such that the expansion in the following case is valid: $(1+x)^{-1/3}$.

\item Expand the following upto 4 terms, taking the values of $x$ such that the expansion in the following case is valid: $(4-3x)^{1/2}$.

\item Expand the following upto 4 terms, taking the values of $x$ such that the expansion in the following case is valid: $\sqrt{\dfrac{1+2x}{1-x}}$.

\item Find the general term in the expansion of $(1+x)^{-3}$ when $|x|<1$.

\item Find the coefficient of $x$ in the expansion of $\dfrac{1-x}{(1+x)^2}$.

\item Using the binomial theorem, find the value of the following to three places of decimals: $\sqrt{99}$.

\item Using the binomial theorem, find the value of the following to three places of decimals: $(1.03)^{1/3}$.

\item Using the binomial theorem, find the value of the following to three places of decimals: $\dfrac{1}{\sqrt[5]{252}}$.

\item Using the binomial theorem, find the value of the following to three places of decimals: $\dfrac{\sqrt{7}}{\sqrt{8}}$.

\item Find the remainder when $8^{100}$ is divided by 7.

\item Find the remainder when $2^{100}$ is divided by 3.

\item Using the binomial theorem, find which number is larger: $19^{10}+20^{10}$ or $21^{10}$.

\item Using the binomial theorem, find which number is larger: $29^{15}+30^{15}$ or $31^{15}$.

\item Using the binomial theorem, show that: $5^7+7^5$ is divisible by 36.

\item Using the binomial theorem, show that: $(17)^{15}+(13)^7$ is divisible by 6.

\item Using the binomial theorem, show that: $(21)^9+(19)^{11}$ is divisible by 20.

\item Using the binomial theorem, find the last digit of the number $(32)^{32}$.

\item By using the binomial theorem, show that for each $n\in\mathbb{N}$, $5^n-1$ is divisible by 4.

\item By using the binomial theorem, show that for each $n\in\mathbb{N}$, $5^n-2^n$ is divisible by 3.

\item Expand $(x+y)^4$ using Pascal's triangle.

\item Expand $(x-2)^5$ using the binomial theorem and using Pascal's triangle.

\item Evaluate $\sqrt[3]{30}$ correct to three places of decimal.

\item Using the binomial theorem, show that $11^n-10n$ leaves a remainder 1 when divided by 100 for all positive integers $n$.

\item Using the binomial theorem, find the last two digits of the number $11^{12}$.

\item Show that $(15)^{13}+(13)^{15}$ is divisible by 14.

\item Which number is larger, $51^{25}$ or $49^{25}+50^{25}$?

\item Find the remainder and quotient by simplifying the following: $(3x^2-x+2)\div(x-1)$.

\item Find the remainder and quotient by simplifying the following: $(x^3+12x^2-3x+4)\div(x-2)$.

\item Find the remainder and quotient by simplifying the following: $(x^4-5x^3-8x^2+13x+12)\div(x-6)$.

\item Find the remainder and quotient by simplifying the following: $(5x^4-3x^3+2x^2-1)\div(x^2+4)$.

\item Find the remainder and quotient by simplifying the following: $(3x^4-5x^3+4x-6)\div(x^2-3x+5)$.

\item Use the factor theorem to determine if the first polynomial is a factor of the second polynomial: $x-2$, $x^2-5x+6$.

\item Use the factor theorem to determine if the first polynomial is a factor of the second polynomial: $x+1$, $x^3+x^2+x-3$.

\item Find the remainder without performing division when $f(x)=x^4+x^3+x^2+1$ is divided by $x+1$.

\item Show that $x-2$ is a factor of $f(x)=x^3-7x+6$ without factorizing.

\item If $x+1$ and $x-2$ are factors of $x^3+px^2+qx+2$, find the values of $p$ and $q$.

\item Consider a data set of monthly sales figures. A polynomial regression model $P(x)=x^3+x^2+2x+1$ is fitted to this data. If the observed sales in the $3^{\text{rd}}$ month are 40 units, find the percentage error.

\item Solve a polynomial regression model $P(x)=3x^3-4x^2+2x-5$. If a data point at $x=-1$ is missing, what should be its predicted value?

\item A signal processing system has a transfer function with a denominator $A(z)=z^2-0.25$. Use the factor theorem to find the poles of the system and determine if the system is stable.

\item Without using the table, find the value of: $\cos(-1230^\circ)$.

\item Without using the table, find the value of: $\tan(-1035^\circ)$.

\item Without using the table, find the value of: $\sec(1140^\circ)$.

\item Without using the table, find the value of: $\csc(-690^\circ)$.

\item Without using the table, find the value of: $\cot(1320^\circ)$.

\item Without using the table, find the value of: $\cos(-240^\circ)$.

\item Express the following as a trigonometric function of an angle of positive degree measure less than $45^\circ$: $\cos 168^\circ$.

\item Express the following as a trigonometric function of an angle of positive degree measure less than $45^\circ$: $\sin 192^\circ$.

\item Express the following as a trigonometric function of an angle of positive degree measure less than $45^\circ$: $\cos 333^\circ$.

\item Express the following as a trigonometric function of an angle of positive degree measure less than $45^\circ$: $\tan 213^\circ$.

\item Express the following as a trigonometric function of an angle of positive degree measure less than $45^\circ$: $\cos(-435^\circ)$.

\item Express the following as a trigonometric function of an angle of positive degree measure less than $45^\circ$: $\sin 219^\circ$.

\item Express the following as a trigonometric function of an angle of positive degree measure less than $45^\circ$: $\tan(-597^\circ)$.

\item Express the following as a trigonometric function of an angle of positive degree measure less than $45^\circ$: $\cos(-111^\circ)$.

\item Express the following as a trigonometric function of an angle of positive degree measure less than $45^\circ$: $\sin(-390^\circ)$.

\item Prove the following: $\sin(180^\circ+\alpha)\sin(90^\circ-\alpha)=-\sin\alpha\cos\alpha$.

\item Prove the following: $\sin 810^\circ\sin 630^\circ+\cos 135^\circ\sin 225^\circ=-\dfrac{1}{2}$.

\item Prove the following: $\tan 150^\circ\cot 330^\circ-2\sec 135^\circ\csc 225^\circ=-3$.

\item Find the distance between the following points: $A(3,8)$, $B(5,6)$.

\item Find the distance between the following points: $P(\cos x,\cos y)$, $Q(\sin x,\sin y)$.

\item Find the value of $\sin\dfrac{5\pi}{12}$.

\item Without using the tables, write down the value of: $\sin 225^\circ$.

\item Without using the tables, write down the value of: $\tan 600^\circ$.

\item Without using the tables, write down the value of: $\cot(-225^\circ)$.

\item Without using the tables, write down the value of: $\csc(-420^\circ)$.

\item Without using tables, find the values of the following: $\sin 15^\circ$.

\item Without using tables, find the values of the following: $\cos 15^\circ$.

\item Without using tables, find the values of the following: $\tan 15^\circ$.

\item Without using tables, find the values of the following: $\sin 105^\circ$.

\item Without using tables, find the values of the following: $\cos 105^\circ$.

\item Without using tables, find the values of the following: $\tan 105^\circ$.

\item Prove that: $\sin(45^\circ+\alpha)=\dfrac{1}{\sqrt{2}}(\sin\alpha+\cos\alpha)$.

\item Prove that: $\cos(\alpha+45^\circ)=\dfrac{1}{\sqrt{2}}(\cos\alpha-\sin\alpha)$.

\item Prove that: $\tan(45^\circ+A)\tan(45^\circ-A)=1$.

\item Prove that: $\tan\left(\dfrac{\pi}{4}-\theta\right)+\tan\left(\dfrac{3\pi}{4}+\theta\right)=0$.

\item Prove that: $\sin\left(\theta+\dfrac{\pi}{6}\right)+\cos\left(\theta+\dfrac{\pi}{3}\right)=\cos\theta$.

\item Prove that: $\sin(\alpha+\beta)\sin(\alpha-\beta)=\sin^2\alpha-\sin^2\beta$ (i) $=\cos^2\beta-\cos^2\alpha$ (ii).

\item Without using tables, find the values of all trigonometric functions of $105^\circ$.

\item Prove that: $\dfrac{\cos 11^\circ+\sin 11^\circ}{\cos 11^\circ-\sin 11^\circ}=\tan 56^\circ$.

\item Find the values of $\sin 2\alpha$, $\cos 2\alpha$ and $\tan 2\alpha$, when: $\sin\alpha=\dfrac{3}{5}$.

\item Find the values of $\sin 2\alpha$, $\cos 2\alpha$ and $\tan 2\alpha$, when: $\cos\alpha=\dfrac{4}{5}$, where $0<\alpha<\dfrac{\pi}{2}$.

\item Prove the following: $\cot\alpha-\tan\alpha=2\cot 2\alpha$.

\item Prove the following: $\dfrac{\sin 2\alpha}{1+\cos 2\alpha}=\tan\alpha$.

\item Prove the following: $\dfrac{1-\cos\alpha}{\sin\alpha}=\tan\dfrac{\alpha}{2}$.

\item Prove the following: $1+\tan\alpha\tan2\alpha = \sec2\alpha$

\item Prove the following: $\dfrac{2\sin\theta\sin2\theta}{\cos\theta+\cos3\theta} = \tan2\theta\tan\theta$

\item Prove the following: $\dfrac{1+\sin2\theta}{1-\sin2\theta} = \tan^2\left(\dfrac{\pi}{4}+\theta\right)$

\item Show that: $2\cos\theta = \sqrt{2+\sqrt{2+2\cos4\theta}}$

\item Find the values of $\sin\theta$ and $\cos\theta$ without using a table or calculator, when $\theta = 72^{\circ}$.

\item Prove that: $\dfrac{\sin\theta+\sin2\theta}{1+\cos\theta+\cos2\theta} = \tan\theta$

\item Show that: $\sin2\theta = \dfrac{2\tan\theta}{1+\tan^2\theta}$

\item Show that: $\cos2\theta = \dfrac{1-\tan^2\theta}{1+\tan^2\theta}$

\item Reduce $\cos^4\theta$ to an expression involving only functions of multiples of $\theta$, raised to the first power.

\item Express the following product as a sum or difference: $2\sin3\theta\cos\theta$

\item Express the following product as a sum or difference: $2\cos5\theta\sin\theta$

\item Express the following product as a sum or difference: $\sin5\theta\cos2\theta$

\item Express the following product as a sum or difference: $2\sin7\theta\sin2\theta$

\item Express the following product as a sum or difference: $\cos(x+y)\sin(x-y)$

\item Express the following product as a sum or difference: $\cos(2x+30^{\circ})\cos(2x+30^{\circ})$

\item Express the following product as a sum or difference: $\sin12^{\circ}\sin46^{\circ}$

\item Express the following product as a sum or difference: $\sin(x+45^{\circ})\sin(x-45^{\circ})$

\item Express the following sum as a product: $\sin5\theta+\sin3\theta$

\item Express the following difference as a product: $\sin8\theta-\sin4\theta$

\item Express the following sum as a product: $\cos6\theta+\cos3\theta$

\item Express the following difference as a product: $\cos7\theta-\cos\theta$

\item Express the following sum as a product: $\cos12^{\circ}+\cos48^{\circ}$

\item Express the following sum as a product: $\sin(x+30^{\circ})+\sin(x-30^{\circ})$

\item Prove that $\tan75^{\circ}-\tan15^{\circ} = 2\sqrt{3}$.

\item Prove that: $\cos15^{\circ}-\sin15^{\circ} = \dfrac{1}{\sqrt{2}}$

\item Express $2\sin7\theta\cos3\theta$ as a sum or difference.

\item Express $\sin5x+\sin7x$ as a product.

\item Express $\cos\theta+\cos3\theta+\cos5\theta+\cos7\theta$ as a product.

\item Determine whether the following function is even, odd, or neither: $\sin^2x$

\item Determine whether the following function is even, odd, or neither: $\sin x+\cos x$

\item Determine whether the following function is even, odd, or neither: $\tan x+\sec x$

\item Determine whether the following function is even, odd, or neither: $\dfrac{1}{\csc^3x}$

\item Find the period of the following function: $\sin5x$

\item Find the period of the following function: $\cos7x$

\item Find the period of the following function: $\tan3x$

\item Find the period of the following function: $\cot\dfrac{x}{2}$

\item Find the period of the following function: $19\sin\left(\dfrac{\pi}{20}x\right)$

\item Find the period of the following function: $\csc\left(\dfrac{2x}{5}\right)$

\item Find the period of the following function: $\dfrac{1}{2}\sin\left(\dfrac{3x}{2}-\dfrac{\pi}{2}\right)$

\item Find the period of the following function: $-5-3\sec\left(7\pi x+\dfrac{\pi}{4}\right)$

\item Find the period of the following function: $12+10\tan\left(\dfrac{\pi}{30}x\right)$

\item Find the period of the following function: $6-4\cot\left(\dfrac{7x}{4}+\dfrac{\pi}{4}\right)$

\item Find the period of the following function: $9+30\sec\left(\dfrac{x}{15}+\dfrac{2\pi}{15}\right)$

\item Find the period of $\sin2x$.

\item Find the period of $3+\tan\dfrac{x}{3}$.

\item Solve graphically: $\sin x=\cos x$, $x\in[0,\pi]$

\item Solve graphically: $\sin x=x$, $x\in[0,\pi]$

\item Find the maximum and minimum values of the following function: $3-\sin3x$

\item Find the maximum and minimum values of the following function: $3+\sin2x$

\item Find the maximum and minimum values of the following function: $\dfrac{1}{2}+\sin(5x+\pi)$

\item Find the maximum and minimum values of the following function: $\dfrac{3}{2}+\cos\left(x-\dfrac{\pi}{4}\right)$

\item Find the maximum and minimum values of the following function: $1-3\cos2x$

\item Find the maximum and minimum values of the following function: $1+2\sin\left(x+\dfrac{\pi}{6}\right)$

\item Find the maximum and minimum values of the following function: $\dfrac{1}{10-2\sin3x}$

\item Find the maximum and minimum values of the following function: $\dfrac{1}{7+3\cos(-2x)}$

\item Find the maximum and minimum values of the following function: $\dfrac{1}{5-3\cos(3x-1)}$

\item Find the maximum and minimum values of the following function: $2+3\sin x$

\item Find the maximum and minimum values of the following function: $5-2\cos3x$

\item Find the limit of the following sequence if it exists: $a_n=\dfrac{2n+3}{n+1}$

\item Find the limit of the following sequence if it exists: $b_n=\dfrac{2n+3}{n^2+1}$

\item Find the limit of the following sequence if it exists: $c_n=\dfrac{5n^2}{2n+3}$

\item Find the limit of the following sequence if it exists: $d_n=\dfrac{n^2-3n+1}{2n^2+n+4}$

\item Evaluate the following limit by using theorems of limits: $\displaystyle\lim_{x\to3}(2x+4)$

\item Evaluate the following limit by using theorems of limits: $\displaystyle\lim_{x\to1}(3x^2-2x+4)$

\item Evaluate the following limit by using theorems of limits: $\displaystyle\lim_{x\to3}\sqrt{x^2+x+4}$

\item Evaluate the following limit by using theorems of limits: $\displaystyle\lim_{x\to2}\sqrt{x^2+4}$

\item Evaluate the following limit by using theorems of limits: $\displaystyle\lim_{x\to2}\left(\sqrt{x^3+1}-\sqrt{x^2+5}\right)$

\item Evaluate the following limit by using theorems of limits: $\displaystyle\lim_{x\to2}\dfrac{2x^3+5x}{3x-2}$

\item Evaluate the following limit by using algebraic techniques: $\displaystyle\lim_{x\to-1}\dfrac{x^3-x}{x+1}$

\item Evaluate the following limit by using algebraic techniques: $\displaystyle\lim_{x\to3}\dfrac{x^2-5x+6}{x^2-2x-3}$

\item Evaluate the following limit by using algebraic techniques: $\displaystyle\lim_{x\to2}\dfrac{x^3-8}{x^2-5x+6}$

\item Evaluate the following limit by using algebraic techniques: $\displaystyle\lim_{x\to1}\dfrac{x^3-3x^2+3x-1}{x^3-x}$

\item Evaluate the following limit by using algebraic techniques: $\displaystyle\lim_{x\to2}\dfrac{x^3-6x^2+12x-8}{x^3-4x}$

\item Evaluate the following limit by using algebraic techniques: $\displaystyle\lim_{x\to1}\dfrac{x^4-1}{x^2-3x+2}$

\item Evaluate the following limit by using algebraic techniques: $\displaystyle\lim_{x\to2}\dfrac{x-2}{\sqrt{x+2}-\sqrt{6-x}}$

\item Evaluate the following limit by using algebraic techniques: $\displaystyle\lim_{h\to0}\dfrac{\sqrt{x+h}-\sqrt{x}}{h}$

\item Evaluate the following limit by using algebraic techniques: $\displaystyle\lim_{x\to a}\dfrac{x^n-a^n}{x^m-a^m}$

\item Evaluate the following limit: $\displaystyle\lim_{x\to0}\dfrac{\sin5x}{x}$

\item Evaluate the following limit: $\displaystyle\lim_{x\to0}\dfrac{\sin x^{\circ}}{x}$

\item Evaluate the following limit: $\displaystyle\lim_{\theta\to0}\dfrac{1-\cos\theta}{\sin\theta}$

\item Evaluate the following limit: $\displaystyle\lim_{x\to\pi/4}\dfrac{\sin x-\cos x}{x-\pi/4}$

\item Evaluate the following limit: $\displaystyle\lim_{x\to0}\dfrac{\cos ax-\cos bx}{x^2}$

\item Evaluate the following limit: $\displaystyle\lim_{x\to\pi/4}\dfrac{\tan x-1}{x-\pi/4}$

\item Evaluate the following limit: $\displaystyle\lim_{x\to0}\dfrac{1-\cos2x}{x^2}$

\item Evaluate the following limit: $\displaystyle\lim_{x\to1}\dfrac{x^3-1}{x^2-1}$

\item Evaluate the following limit: $\displaystyle\lim_{x\to0}\dfrac{x(2^x-1)}{1-\cos x}$

\item Express the following limit in terms of $e$: $\displaystyle\lim_{n\to\infty}\left(1+\dfrac{1}{n}\right)^{2n}$

\item Express the following limit in terms of $e$: $\displaystyle\lim_{n\to\infty}\left(1+\dfrac{1}{n}\right)^{n/2}$

\item Express the following limit in terms of $e$: $\displaystyle\lim_{n\to\infty}\left(1-\dfrac{1}{n}\right)^{n}$

\item Express the following limit in terms of $e$: $\displaystyle\lim_{n\to\infty}\left(1+\dfrac{1}{3n}\right)^{n}$

\item Express the following limit in terms of $e$: $\displaystyle\lim_{n\to\infty}\left(1+\dfrac{4}{n}\right)^{n}$

\item Express the following limit in terms of $e$: $\displaystyle\lim_{x\to0}(1+3x)^{2/x}$

\item Express the following limit in terms of $e$: $\displaystyle\lim_{x\to0}(1+2x^2)^{1/x^2}$

\item Express the following limit in terms of $e$: $\displaystyle\lim_{x\to0}\dfrac{e^{ax}-e^{bx}}{abx}$

\item Express the following limit in terms of $e$: $\displaystyle\lim_{x\to\infty}\left(\dfrac{x}{1+x}\right)^{x}$

\item Express the following limit in terms of $e$: $\displaystyle\lim_{x\to0}\dfrac{e^{1/x}-1}{e^{1/x}+1}$, $x<0$

\item Express the following limit in terms of $e$: $\displaystyle\lim_{x\to0}\dfrac{e^{1/x}-1}{e^{1/x}+1}$, $x>0$

\item Express the following limit in terms of $e$: $\displaystyle\lim_{x\to2}\dfrac{e^x-e^2}{x-2}$

\item Evaluate: $\displaystyle\lim_{x\to1}\dfrac{x^2-1}{x^2-x}$

\item Evaluate: $\displaystyle\lim_{x\to3}\dfrac{x-3}{\sqrt{x}-\sqrt{3}}$

\item Find the limit of the sequence $a_n=\dfrac{2n+3}{n+1}$.

\item Evaluate $\displaystyle\lim_{x\to+\infty}\dfrac{5x^4-10x^2+1}{-3x^3+10x^2+50}$.

\item Evaluate $\displaystyle\lim_{x\to-\infty}\dfrac{4x^4-5x^3}{-3x^5+2x^2+1}$

\item Evaluate: $\displaystyle\lim_{x\to-\infty}\dfrac{2-3x}{\sqrt{3+4x^2}}$

\item Evaluate: $\displaystyle\lim_{x\to+\infty}\dfrac{2-3x}{\sqrt{3+4x^2}}$

\item $^{\dagger}$ Express the following limit in terms of $e$: $\displaystyle\lim_{n\to+\infty}\left(1+\dfrac{3}{n}\right)^{2n}$

\item Express the following limit in terms of $e$: $\displaystyle\lim_{n\to0}(1+2n)^{1/n}$

\item Evaluate: $\displaystyle\lim_{\theta\to0}\dfrac{\sin7\theta}{\theta}$

\item $^{\dagger}$ Evaluate: $\displaystyle\lim_{\theta\to0}\dfrac{1-\cos\theta}{\theta^2}$

\item Determine the left-hand limit and the right-hand limit, and then find the limit of the following function as $x\to c$: $f(x)=2x^2+x-5$, $c=1$

\item Determine the left-hand limit and the right-hand limit, and then find the limit of the following function as $x\to c$: $f(x)=\dfrac{x^2-9}{x-3}$, $c=-3$

\item Determine the left-hand limit and the right-hand limit, and then find the limit of the following function as $x\to c$: $f(x)=|x-5|$, $c=5$

\item A stock price grows according to the function $P(t)=50e^{0.05t}$. Find the limit of $P(t)$ as $t\to\infty$.

\item A stock price grows according to the function $P(t)=50e^{0.05t}$. Calculate the price after $10$ years.

\item Find, by definition, the derivative with respect to $x$ of the function defined as: $\dfrac{1}{\sqrt{x}}$

\item Find, by definition, the derivative with respect to $x$ of the function defined as: $x(x-3)$

\item Find the derivative of $x^{2/3}$ at $x=8$ from the first principle.

\item Find the derivative of $x^2+2x+3$ by definition.

\item Find the gradient and equation of the tangent line to $y=3x^2-4x+1$ at $x=2$.

\item For the function $f(x)=2x^3+x$, calculate the equation of the tangent line at $x=-1$.

\item Find the coordinate of the point of tangency and the equation of the tangent line for $f(x)=x^3-2x+1$ at $x=1$.

\item Find the gradient of the curve $f(x)=3x^2+2x$ at $x=1$.

\item Find the gradient and an equation of the tangent line to the graph of $f(x)=\sqrt{x}$ at $x=9$.

\item The position of a car after $t$ hours is given by $s(t)=2t^3-3t^2+t$ (in kilometres). Find the average velocity over the interval $[1,4]$.

\item The position of a car after $t$ hours is given by $s(t)=2t^3-3t^2+t$ (in kilometres). Find the instantaneous velocity at $t=2$.

\item A stone is thrown upwards and its height after $t$ seconds is given by $s(t)=-16t^2+32t+10$ (in feet). Find the instantaneous velocity at $t=1$.

\item Find the gradient and an equation of the tangent line to the graph of $f(x)=x^2-2$ at the point $P(-1,1)$.

\item Find the derivative of the following function by definition: $f(x)=c$

\item Find the derivative of the following function by definition: $f(x)=x^2$

\item Find the derivative of $\sqrt{x}$ at $x=a$ from the first principle.

\item Differentiate with respect to $x$: $x^4+2x^3+x^2$

\item Differentiate with respect to $x$: $x^{-3}+2x^{-3/2}+3$

\item Differentiate with respect to $x$: $\dfrac{2x-3}{2x+1}$

\item Differentiate with respect to $x$: $\dfrac{(1+\sqrt{x})(x-x^{3/2})}{\sqrt{x}}$

\item Differentiate with respect to $x$: $\left(\sqrt{x}-\dfrac{1}{\sqrt{x}}\right)^2$

\item Differentiate with respect to $x$: $(x-5)(3-x)$

\item Differentiate with respect to $x$: $\dfrac{(x^2+1)^2}{x^2-1}$

\item Differentiate with respect to $x$: $\dfrac{x^2+1}{x^2-3}$

\item Calculate $\dfrac{d}{dx}\left(3x^{4/3}\right) = 3\dfrac{d}{dx}\left(x^{4/3}\right)$.

\item Find the derivative of $y=(x^2+5)(x^3+7)$ with respect to $x$.

\item Find the derivative of $y=(2\sqrt{x}+2)(x-\sqrt{x})$.

\item Find the derivative of $y=(2\sqrt{x}+2)(x-\sqrt{x})$ with respect to $x$.

\item Differentiate $\dfrac{2x^3-3x^2+5}{x^2+1}$ with respect to $x$.

\item A car's position at time $t$ is given by $s(t)=5t^3-3t^2+t$. Find the velocity by differentiating the position function with respect to time.

\item Let $\underline{u}=3\underline{i}+2\underline{j}-5\underline{k}$, $\underline{v}=\underline{i}-5\underline{j}-\underline{k}$ and $\underline{w}=-4\underline{i}-\underline{j}+7\underline{k}$. Find: $\underline{u}+2\underline{v}+\underline{w}$

\item Using the same $\underline{u},\underline{v},\underline{w}$, find: $\underline{v}-3\underline{w}$

\item Using the same $\underline{u},\underline{v},\underline{w}$, find: $|3\underline{v}+\underline{w}|$

\item Find the magnitude of the vector $\underline{v}$ and write the direction cosines of $\underline{v}$: $\underline{v}=3\underline{i}-2\underline{j}+6\underline{k}$

\item Find the magnitude of the vector $\underline{v}$ and write the direction cosines of $\underline{v}$: $\underline{v}=-4\underline{i}+4\underline{j}+2\underline{k}$

\item Find the magnitude of the vector $\underline{v}$ and write the direction cosines of $\underline{v}$: $\underline{v}=-6\underline{i}+8\underline{j}$

\item Find $t$ so that (equation as given in the original figure/table).

\item Find a unit vector in the direction of $\underline{v}=-\underline{i}+4\underline{j}-8\underline{k}$.

\item If $\underline{u}=2\underline{i}+\underline{j}-3\underline{k}$, $\underline{v}=-\underline{i}+4\underline{j}+2\underline{k}$ and $\underline{w}=3\underline{i}-2\underline{j}+\underline{k}$, find the unit vector parallel to $4\underline{u}-3\underline{v}+2\underline{w}$.

\item Find the vector whose magnitude is $5$ and is parallel to $3\underline{i}+4\underline{j}-\underline{k}$.

\item Find the vector whose magnitude is $7$ and is parallel to $-\underline{i}+\underline{j}+\underline{k}$.

\item A spacecraft moves from point $(120,240,-50)$ to point $(130,210,80)$ in kilometres. What is the magnitude of the displacement vector in kilometres?

\item Find the direction cosines for the given vector: $\underline{u}=6\underline{i}+3\underline{j}+2\underline{k}$

\item Find the direction cosines for the given vector: $\underline{v}=4\underline{i}+2\underline{j}-5\underline{k}$

\item Find the direction cosines for the given vector $\overrightarrow{PQ}$ where $P(9,3,13)$ and $Q(11,6,19)$.

\item Can the following triple be the direction angles of a single vector? $45^{\circ}, 45^{\circ}, 60^{\circ}$

\item Can the following triple be the direction angles of a single vector? $30^{\circ}, 45^{\circ}, 60^{\circ}$

\item Can the following triple be the direction angles of a single vector? $45^{\circ}, 60^{\circ}, 60^{\circ}$

\item For the vectors $\underline{u}=[1,-2,3]$, $\underline{v}=[2,1,3]$ and $\underline{w}=[-1,4,0]$, find: $\underline{v}+\underline{w}$

\item For the same vectors, find: $2\underline{w}$

\item For the same vectors, find: $|\underline{u}|$

\item For the same vectors, find: $|\underline{v}-2\underline{w}|$

\item For the same vectors, find: $|2\underline{u}-\underline{v}+3\underline{w}|$

\item Find the unit vector of $\underline{u}=2\underline{i}+5\underline{j}-\underline{k}$.

\item Find the cosine of the angle $\theta$ between $\underline{u}$ and $\underline{v}$: $\underline{u}=2\underline{i}+3\underline{j}+\underline{k}$, $\underline{v}=-\underline{i}+2\underline{j}+2\underline{k}$

\item Find the cosine of the angle $\theta$ between $\underline{u}$ and $\underline{v}$: $\underline{u}=[-3,2,5]$, $\underline{v}=[1,6,-2]$

\item If $\underline{a}+\underline{b}+\underline{c}=\underline{0}$ and $|\underline{a}|=3$, $|\underline{b}|=5$ and $|\underline{c}|=7$, find the angle between $\underline{a}$ and $\underline{b}$.

\item If $|\underline{a}|=3$, $|\underline{b}|=4$ and $|\underline{a}+\underline{b}|=5$, find the angle between $\underline{a}$ and $\underline{b}$.

\item Find a real number $a$ so that the vectors $\underline{u}$ and $\underline{v}$ are perpendicular: $\underline{u}=a\underline{i}+3\underline{j}+\underline{k}$, $\underline{v}=\underline{i}-2\underline{j}+a\underline{k}$

\item Find a real number $a$ so that the vectors $\underline{u}$ and $\underline{v}$ are perpendicular: $\underline{u}=a\underline{i}+2a\underline{j}-\underline{k}$, $\underline{v}=\underline{i}+a\underline{j}+3\underline{k}$

\item If $|\underline{a}+\underline{b}|=|\underline{a}-\underline{b}|$, then show that $\underline{a}$ and $\underline{b}$ are perpendicular.

\item Show that the vectors $3\underline{i}-2\underline{j}+\underline{k}$, $\underline{i}-3\underline{j}+5\underline{k}$ and $2\underline{i}+\underline{j}-4\underline{k}$ form a right triangle.

\item Prove that in any triangle $ABC$: $c^2=a^2+b^2-2ab\cos C$

\item If $\underline{u}=3\underline{i}-\underline{j}-2\underline{k}$ and $\underline{v}=\underline{i}+2\underline{j}-\underline{k}$, find $\underline{u}\cdot\underline{v}$.

\item If $\underline{u}=2\underline{i}-4\underline{j}+5\underline{k}$ and $\underline{v}=4\underline{i}-3\underline{j}-4\underline{k}$, then prove that $\underline{u}$ and $\underline{v}$ are orthogonal.

\item Find the angle between the vectors $\underline{u}=2\underline{i}-\underline{j}+\underline{k}$ and $\underline{v}=-\underline{i}+\underline{j}$.

\item Find the area of the triangle formed by the points $P,Q,R$: $P(2,3,5)$; $Q(1,2,3)$; $R(4,1,2)$.

\item Find the area of the triangle formed by the points $P,Q,R$: $P(0,0,1)$; $Q(2,-1,2)$; $R(-1,3,2)$.

\item If the cross product of the vectors $\underline{u}=7\underline{i}-4\underline{j}+5\underline{k}$ and $\underline{v}=a\underline{i}-b\underline{j}+3\underline{k}$ is zero, find the values of $a$ and $b$.

\item Which vectors, if any, are perpendicular or parallel: $\underline{u}=5\underline{i}-\underline{j}+\underline{k}$; $\underline{v}=\underline{j}-5\underline{k}$; $\underline{w}=-15\underline{i}+3\underline{j}-3\underline{k}$

\item $^{\dagger}$ Which vectors, if any, are perpendicular or parallel: $\underline{u}=\underline{i}+2\underline{j}-\underline{k}$; $\underline{v}=-\underline{i}+\underline{j}+\underline{k}$; $\underline{w}=-\dfrac{\pi}{2}\underline{i}-\pi\underline{j}+\dfrac{\pi}{2}\underline{k}$

\item Using the definition of cross product, for any vectors $\underline{u},\underline{v},\underline{w}$ and scalar $k$, prove that: $\underline{u}\times(-\underline{u})=\underline{0}$

\item Prove that: $\underline{a}\times(\underline{b}+\underline{c})+\underline{b}\times(\underline{c}+\underline{a})+\underline{c}\times(\underline{a}+\underline{b})=\underline{0}$

\item If $\underline{a}+\underline{b}+\underline{c}=\underline{0}$, then prove that $\underline{a}\times\underline{b}=\underline{b}\times\underline{c}=\underline{c}\times\underline{a}$.

\item Show that: $|\underline{a}\times\underline{b}|^2 = |\underline{a}|^2|\underline{b}|^2-(\underline{a}\cdot\underline{b})^2$

\item $^{\dagger}$ A force $\underline{F}=-2\underline{i}+\underline{j}-3\underline{k}$ is applied at $P(-1,-3,2)$. Find its moment about the point $Q(4,2,2)$.

\item If $\underline{u}=2\underline{i}-\underline{j}+\underline{k}$ and $\underline{v}=4\underline{i}+2\underline{j}-\underline{k}$, find by the determinant formula: $\underline{u}\times\underline{u}$

\item If $\underline{u}=2\underline{i}-\underline{j}+\underline{k}$ and $\underline{v}=4\underline{i}+2\underline{j}-\underline{k}$, find by the determinant formula: $\underline{u}\times\underline{v}$

\item If $\underline{u}=2\underline{i}-\underline{j}+\underline{k}$ and $\underline{v}=4\underline{i}+2\underline{j}-\underline{k}$, find by the determinant formula: $\underline{v}\times\underline{u}$

\item Prove that the vectors $\underline{i}-2\underline{j}+3\underline{k}$, $-2\underline{i}+3\underline{j}-4\underline{k}$ and $\underline{i}-3\underline{j}+5\underline{k}$ are coplanar.

\item Find the constant $a$ such that the vectors are coplanar: $\underline{i}-\underline{j}+\underline{k}$, $\underline{i}-2\underline{j}-3\underline{k}$ and $3\underline{i}-a\underline{j}+5\underline{k}$

\item Find the value of: $2\underline{i}\times2\underline{j}\cdot\underline{k}$

\item Find the value of: $3\underline{j}\cdot\underline{k}\times\underline{i}$

\item Find the value of: $[\underline{k}\ \underline{i}\ \underline{j}]$

\item Find the value of: $[\underline{i}\ \underline{i}\ \underline{k}]$

\item Prove that: $\underline{u}\cdot(\underline{v}\times\underline{w})+\underline{v}\cdot(\underline{w}\times\underline{u})+\underline{w}\cdot(\underline{u}\times\underline{v})=3\underline{u}\cdot(\underline{v}\times\underline{w})$

\item Find the volume of the tetrahedron with vertices: $(0,1,2)$, $(3,2,1)$, $(1,2,1)$ and $(5,5,6)$.

\item Find the volume of the tetrahedron with vertices: $(2,1,8)$, $(3,2,9)$, $(2,1,4)$ and $(3,3,10)$.

\item Prove that the dot and cross product are interchangeable in the scalar triple product.

\item Find the volume of the parallelepiped determined by: $\underline{u}=\underline{i}+2\underline{j}-\underline{k}$, $\underline{v}=\underline{i}-2\underline{j}+3\underline{k}$, $\underline{w}=\underline{i}-7\underline{j}-4\underline{k}$

\item Find the volume of the tetrahedron whose vertices are $A(2,1,8)$, $B(3,2,9)$, $C(2,1,4)$ and $D(3,3,0)$.

\item Prove that the four points $A(-3,5,-4)$, $B(-1,1,1)$, $C(-1,2,2)$ and $D(-3,4,-5)$ are coplanar.

\item Find the value of $a$ so that $a\underline{i}+\underline{j}$, $\underline{i}+\underline{j}+3\underline{k}$ and $2\underline{i}+\underline{j}-2\underline{k}$ are coplanar.

\end{shortlist}

\newpage

\section{Long Questions (1--668)}

\section*{Write detailed answers of the following questions.}

\begin{longlist}

\item Prove that $\overline{\bar{z}}=z$ if $z$ is real.

\item For $z\in\mathbb{C}$, show that: $\dfrac{z+\bar{z}}{2}=\operatorname{Re}(z)$.

\item For $z\in\mathbb{C}$, show that: $\dfrac{z-\bar{z}}{2i}=\operatorname{Im}(z)$.

\item If $z_1=2+i$, $z_2=3-2i$, $z_3=1+3i$ then express $\dfrac{z_1 z_3}{z_2}$ in the form of $a+ib$.

\item If $z_1=2+7i$ and $z_2=-5+3i$, then evaluate the following: $\left|(z_1+z_2)^3\right|$.

\item Show that $i^{n+1}+i^{n+2}+i^{n+3}+i^{n+4}=0$ for all $n\in\mathbb{N}$.

\item Show that the value of $i^n$ for $n\in\mathbb{N}$ and $n>4$ is $i^r$ where $r$ is the remainder when $n$ is divided by 4.

\item If $z_1=x+yi$ and $z_2=a+bi$, find $x,y,a,b$ such that $z_1+z_2=10+4i$ and $z_1-z_2=6+2i$.

\item Find the square root of $13-20\sqrt{3}\,i$ and represent it on an Argand diagram.

\item Find the real values of $x$ and $y$ if $(-7+i)(x+iy)+(-1-5i)=i(11-i)$.

\item Find the real values of $x$ and $y$ if $(5-2i)(x+iy)+3=i(11-i)-4i$.

\item Find the real values of $u$ and $v$ if $\dfrac{u-2}{2+i}+\dfrac{v-3}{2-i}=4i$.

\item If $(3+2i)(x+iy)=5+12i$ where $x,y\in\mathbb{R}$, then find the value of $x$ and $y$.

\item Factorize the following: $2z^2-22z+65$.

\item Factorize the following polynomial into its linear factors: $x^3+27$.

\item Find the roots of $z^4+7z^2-144=0$ and hence express it as a product of linear factors.

\item Solve the following complex quadratic equation by completing the square method: $2z^2-3z+4=0$.

\item Solve the following complex quadratic equation by completing the square method: $z^2-6z+30=0$.

\item Solve the following complex quadratic equation by completing the square method: $3z^2-18z+50=0$.

\item Solve the following equation: $2z^4-32=0$.

\item Solve the following equation: $3z^5-243z=0$.

\item Solve the following equation: $4z^4-25z^2-21=0$.

\item Find a polynomial $P(z)$ of degree 3 with zeros $3,-2i,2i$ and satisfying $P(1)=20$.

\item Find a polynomial $P(z)$ of degree 4 with zeros $2i,-2i,1,-1$ and satisfying $P(2)=240$.

\item Find a polynomial $P(z)$ of degree 4 with zeros $4,-4,1+i,1-i$ and satisfying $P(2)=72$.

\item Factorize the polynomial $P(z)=z^3-3z^2+z+5$.

\item Solve the equation $2z^2-12z+50=0$ by completing the square method and hence express it as a product of its linear factors.

\item Find the fourth roots of $16$, $81$, $625$. Also show that their sum is zero in each case.

\item If $1,\omega,\omega^2$ are the cube roots of unity, show that $1+\omega^n+\omega^{2n}$ where $n$ is a multiple of 3 (evaluate accordingly).

\item Show that: $(1-\omega+\omega^2)(1-\omega^2+\omega^4)(1-\omega^4+\omega^8)(1-\omega^8+\omega^{16})\cdots$ to $2n$ factors.

\item Prove that: $\left(\dfrac{i+\sqrt{3}}{2}\right)^8+\left(\dfrac{i-\sqrt{3}}{2}\right)^8=-1$.

\item If $\omega$ is a cube root of unity, prove that $\dfrac{a\omega^{12}+b\omega^{17}+c\omega^{19}}{a\omega^{14}+b\omega^{22}+c\omega^{30}}=\omega$.

\item If $z_1$ and $z_2$ are two complex numbers, show that: $\operatorname{Arg}(z_1 z_2)=\operatorname{Arg}\,z_1+\operatorname{Arg}\,z_2$.

\item If $z_1$ and $z_2$ are two complex numbers, show that: $\operatorname{Arg}\left(\dfrac{z_1}{z_2}\right)=\operatorname{Arg}\,z_1-\operatorname{Arg}\,z_2$.

\item If $z=x+iy$ and $\arg\left(\dfrac{\bar{z}-1+2i}{\bar{z}+1-2i}\right)=\dfrac{9\pi}{4}$, show that $x^2+y^2-4x+2y-5=0$.

\item If $z=x+iy$ and $\arg(z-2-3i)-\arg(z+2+3i)=2$, show that $2y=3x$.

\item Determine the set of points $z=x+iy$ that satisfy the equation $|3\bar{z}-2+i|=|3z+i|$.

\item If $z=x+iy$ and $W=\dfrac{1-iz}{z-i}$, show that $|W|=1\Rightarrow z$ is real.

\item If $z_1$ and $z_2$ are different complex numbers with $|z_2|=1$, find $\left|\dfrac{z_2-z_1}{1-z_1\bar{z_2}}\right|$.

\item An AC source supplies a voltage of $V=120\left(\cos\dfrac{\pi}{4}+i\sin\dfrac{\pi}{4}\right)$ volts to a circuit with impedance $Z=\dfrac{1+i\sqrt{3}}{2}$ ohms. Calculate the current in polar form.

\item An AC circuit has an impedance of $z=3-6i$ ohms and is connected to a voltage source of $V=90+30i$ volts. Find the current in both rectangular and polar form.

\item Encrypt the word ``CODE'' by multiplying the complex encryption key $k=2-i$. Then decrypt it back to the original word.

\item Consider the complex encryption $k=3-3i$. Encrypt the word ``QUIZ'', and then recover the original word using the inverse of the key.

\item Encrypt the word ``Class'' by adding the complex number encryption key $k=-3+4i$. Then decrypt it back to the original word.

\item Find the product of $5\left(\cos\dfrac{\pi}{6}+i\sin\dfrac{\pi}{6}\right)$ and $4\left(\cos\dfrac{3\pi}{2}+i\sin\dfrac{3\pi}{2}\right)$.

\item Divide $\dfrac{2}{7}\left(\cos\dfrac{7\pi}{6}+i\sin\dfrac{7\pi}{6}\right)$ by $\dfrac{3}{5}\left(\cos\left(-\dfrac{\pi}{2}\right)+i\sin\left(-\dfrac{\pi}{2}\right)\right)$.

\item If $z=x+iy$, then write the equation $|3\bar{z}-i|=|3z+7|$ in terms of $x$ and $y$.

\item Show that $(x+2)^2+y^2=8$ if $\arg\left(\dfrac{z+2i}{z-2i}\right)=\dfrac{3\pi}{4}$ for $z=x+iy$.

\item If the impedance of a circuit is $11(\cos 55.35^\circ+i\sin 55.35^\circ)$ ohms at a voltage of $25(\cos 30^\circ+i\sin 30^\circ)$ V, find the value of the current in the circuit.

\item Encrypt the word ``MATH'' by multiplying it with a complex number $k=2+3i$ and then decrypt it back to its original form using the concept of multiplicative inverse in complex numbers. Each letter of the alphabet is assigned a numerical value as follows: A=1, B=2, C=3, \ldots, Z=26.

\item Find the domain and the range of the function $g$ defined below:
\[
g(x)=\begin{cases}6x+7, & x\leq -2\\ 4-3x, & x>-2\end{cases}
\]

\item A stone falls from a height of $60$m on the ground, the height $h$ after $x$ seconds is approximately given by $h(x)=40-10x^2$. What is the height of the stone when: (a) $x=1$ sec, (b) $x=1.5$ sec, (c) $x=1.7$ sec.

\item A stone falls from a height of $60$m on the ground, the height $h$ after $x$ seconds is approximately given by $h(x)=40-10x^2$. When does the stone strike the ground?

\item Let $f:\mathbb{R}\to\mathbb{R}$ be defined by $f(x)=\dfrac{2x-3}{x+1}$. Find the domain and range of $f(x)$.

\item Let $f:\mathbb{R}\to\mathbb{R}$ be defined by $f(x)=\dfrac{2x-3}{x+1}$. Determine whether $f(x)$ is into.

\item Let $f:\mathbb{R}\to\mathbb{R}$ be defined by $f(x)=\dfrac{2x-3}{x+1}$. Prove that $f(x)$ is one-to-one.

\item Consider the function $f:\mathbb{R}^+\to\mathbb{R}^+$ defined by $f(x)=e^{-x}$. Show that $f(x)$ is bijective.

\item Let $g:\mathbb{R}\to\mathbb{R}$ be given by $g(x)=x^3-3x$. Determine if $g(x)$ is injective and/or surjective.

\item Show that the function $f(x)=x+1$, where the domain and co-domain are all real numbers, is bijective.

\item Show that the function $f(x)=x^2-2$, where the domain and co-domain are all real numbers, is neither one-to-one nor onto.

\item Find the point(s) of intersection of the following functions graphically: $f(x)=2x+5,\ g(x)=-x+5$.

\item Find the point(s) of intersection of the following functions graphically: $f(x)=3x-2,\ g(x)=10-x$.

\item Find the point(s) of intersection of the following functions graphically: $f(x)=3x-2,\ g(x)=10-x$.

\item Find the point(s) of intersection of the following functions graphically: $f(x)=2x-4,\ g(x)=3x-1$.

\item Find the point(s) of intersection of the following functions graphically: $f(x)=x-1,\ g(x)=x^2-4x+3$.

\item Find the point(s) of intersection of the following functions graphically: $f(x)=3x+4,\ g(x)=x^2+2x-8$.

\item Find the point(s) of intersection of the following functions graphically: $f(x)=-2x-1,\ g(x)=x^2-4x$.

\item Find the point(s) of intersection of the following functions graphically: $f(x)=-x^2-3x+2,\ g(x)=x+6$.

\item Graph the following function: $y=-\sqrt{x+1}+2$.

\item Graph the following function: $y=\sqrt[3]{2x+1}$.

\item Graph the following function: $y=2\sqrt[3]{x-3}$.

\item Graph the following function: $y=\sqrt{x^2+x-2}$.

\item A building's height over time is modeled by $H(t)=100+20t$ (in metres, $t$ in months). The height of a growing tree nearby is given by $T(t)=50+10t+t^2$.
(i) At what time will the building and tree have the same height?
(ii) What will the height be? Sketch the graph of both functions and determine the time when the tree will overtake the height of the building.

\item A radioactive substance has a half-life ($h$) of 2 years. If the initial quantity $Q_0$ is 200 grams and the exponential decay function is $Q(t)=Q_0\left(\dfrac{1}{2}\right)^{1/h}$, then find the remaining quantity after 6 years graphically.

\item Solve the linear function $y=-x+3$ and quadratic function $y=x^2-6x+3$ graphically.

\item Graph the square root function $y=2\sqrt{x+1}$.

\item Graph the cube root function $y=\sqrt[3]{x-1}$.

\item The value of a stock follows the exponential growth model $P(t)=P_0e^{rt}$, where $P_0$ is the initial stock price, $r$ is the growth rate per year, and $t$ is the time in years. Suppose a stock is currently valued at Rs.\ 5{,}000 and is expected to grow at a rate of 5\% per year. (i) Find the value of the stock after 10 years. (ii) After many years will the stock double in value?

\item Find the maximum and minimum value of the following quadratic function by completing squares: $f(x)=-2x^2-x+21$.

\item Find the maximum and minimum point by sketching the following quadratic function. Also find their domain and range: $f(x)=x^2-4x$.

\item Find the maximum and minimum point by sketching the following quadratic function. Also find their domain and range: $f(x)=x^2-5x+6$.

\item Find the maximum and minimum point by sketching the following quadratic function. Also find their domain and range: $f(x)=-x^2+2x-8$.

\item Find the maximum and minimum point by sketching the following quadratic function. Also find their domain and range: $f(x)=-x^2-4x+4$.

\item Find the maximum and minimum point by sketching the following quadratic function. Also find their domain and range: $f(x)=x^2+2x-8.3$.

\item Find the inverse of the following quadratic function. Also find their domain and range: $f(x)=x^2-3,\ x\leq 0$.

\item Find the inverse of the following quadratic function. Also find their domain and range: $f(x)=x^2+6x+4,\ x<-3$.

\item Find the inverse of the following quadratic function. Also find their domain and range: $f(x)=2x^2-8x+11,\ x\geq 2$.

\item Find the inverse of the following quadratic function. Also find their domain and range: $f(x)=3x^2-2x+6,\ x\geq 5$.

\item Find the inverse of the following quadratic function. Also find their domain and range: $f(x)=2(x-3)^2+1,\ x\geq 3$.

\item Find the inverse of the following quadratic function. Also find their domain and range: $f(x)=-3(x+4)^2-5,\ x<-4$.

\item Solve the following absolute value quadratic equation and inequalities: $|3x^2-7x+2|=x^2-x+1$.

\item Solve the following absolute value quadratic equation and inequalities: $|x^2-3x+2|>4$.

\item Solve the following absolute value quadratic equation and inequalities: $|x^2-5x+6|\leq x+2$.

\item Solve the following absolute value quadratic equation and inequalities: $|2x^2-3x-5|<4$.

\item Sketch and analyze: $y=-x^2-2x+3$.

\item Find the inverse of $f(x)=x^2+4x+3,\ x\geq -2$. Also find its domain and range.

\item Solve: $|x^2-6x-4|<3$.

\item Solve the following: $\dfrac{1}{x+1}+\dfrac{2}{x+2}=\dfrac{7}{x+5};\ x\neq -1,-2,-5$.

\item Solve the following: $\dfrac{a}{ax-1}+\dfrac{b}{bx-1}=a+b;\ x\neq \dfrac{1}{a},\dfrac{1}{b}$.

\item Solve the following: $\dfrac{x+1}{x-1}+\dfrac{x-1}{x+1}=2,\ x\neq 1,\ x\neq -1$.

\item Solve the following: $3x^2+15x-2\sqrt{x^2+5x+1}=2$.

\item Solve the following: $\sqrt{2x+8}+\sqrt{x+5}=7$.

\item Solve the following: $\sqrt{3x+4}=2+\sqrt{2x-4}$.

\item Solve the following: $\sqrt{x+7}+\sqrt{x+2}=\sqrt{6x+13}$.

\item A farmer bought some sheep for Rs.\ 9000. If he had paid Rs.\ 100 less for each, he would have got 3 sheep more for the same money. How many sheep did he buy, when the rate in each case is uniform?

\item A man sold his stock of eggs for Rs.\ 2400. If he had 2 dozen more, he would have got the same money by selling the whole for Rs.\ 0.50 per dozen cheaper. How many dozen eggs did he sell?

\item A cyclist travelled 48 km at a uniform speed. If he had travelled 5 km/hour slower, he would have taken 2 hours more to perform the journey. How long did he take to cover 48 km?

\item To do a piece of work, Abdaullah takes 10 days more than Abdul Hadi. Together they finish the work in 12 days. How long would Abdul Hadi take to finish it alone?

\item The braking distance (in metres) of a car is modeled by $d(s)=0.02s^2+0.1s$, where $s$ is the speed of the car in km/h. If the maximum safe braking distance is 50 metres, find the range of speed where braking is safe.

\item A rocket follows the height function $h(t)=5t^2+20t+30$, where $h(t)$ is the height in metres and $t$ is the time in seconds. Find the time interval during which the rocket is at least 40 metres above the ground.

\item A company manufactures laptops and its weekly profit function (in thousands of dollars) is $P(x)=-x^2+2x+3$, where $x$ is the number of laptops produced (in hundreds). Find the range of production levels where the company makes at least \$4{,}000 profit.

\item If $A=\begin{pmatrix}i & 2i\\ 1 & -i\end{pmatrix}$, $B=\begin{pmatrix}-i & 1\\ 2i & 1\end{pmatrix}$ and $C=\begin{pmatrix}2i & 1\\ -i & i\end{pmatrix}$, then show that: $(AB)C=A(BC)$.

\item If $A=\begin{pmatrix}i & 2i\\ 1 & -i\end{pmatrix}$, $B=\begin{pmatrix}-i & 1\\ 2i & 1\end{pmatrix}$ and $C=\begin{pmatrix}2i & 1\\ -i & i\end{pmatrix}$, then show that: $A(B+C)=AB+AC$.

\item If $A$ and $B$ are square matrices of the same order, then explain why in general $(A+B)(A-B)\neq A^2-B^2$.

\item Solve the matrix equation $A^2-5A+4I-X=0$ if $A=\begin{pmatrix}2 & 0 & 1\\ 2 & 1 & 3\\ 1 & -1 & 0\end{pmatrix}$.

\item If $A=\begin{pmatrix}1 & 0 & -1 & 2\\ 3 & 1 & 2 & 5\\ 0 & -2 & 1 & 6\end{pmatrix}$ and $B=\begin{pmatrix}2 & -1 & 3 & 1\\ 1 & 3 & -1 & 4\\ 3 & 1 & 2 & -1\end{pmatrix}$, then show that $(A+B)^t=A^t+B^t$.

\item If $A=\begin{pmatrix}2 & -1 & 0\\ 1 & 2 & -3\\ 1 & 2 & -2\end{pmatrix}$ and $B=\begin{pmatrix}2 & -2 & 3\\ -1 & -4 & 6\\ 0 & -5 & 5\end{pmatrix}$, then compute $A^2B$.

\item Without expansion show that: $\begin{vmatrix}bc & a & a^2\\ ca & b & b^2\\ ab & c & c^2\end{vmatrix}=\begin{vmatrix}1 & a^2 & a^3\\ 1 & b^2 & b^3\\ 1 & c^2 & c^3\end{vmatrix}$.

\item Without expansion show that: $\begin{vmatrix}1 & p^2 & \dfrac{p}{qr}\\ 1 & q^2 & \dfrac{q}{rp}\\ 1 & r^2 & \dfrac{r}{pq}\end{vmatrix}=0$.

\item Without expansion show that: $\begin{vmatrix}p-q & q-r & r-p\\ q-r & r-p & p-q\\ r-p & p-q & q-r\end{vmatrix}=0$.

\item Without expansion show that: $\begin{vmatrix}2a & a+b & a+c\\ 2b & 2b & b+c\\ 2c & b+c & 2c\end{vmatrix}=0$.

\item Using properties of determinants, show that: $\begin{vmatrix}3 & 5 & 0\\ 5 & 25 & 10\\ 7 & 25 & 1\end{vmatrix}=25\begin{vmatrix}3 & 1 & 0\\ 1 & 1 & 2\\ 7 & 5 & 1\end{vmatrix}$.

\item Using properties of determinants, show that: $\begin{vmatrix}a+x & a & a\\ a & a+x & a\\ a & a & a+x\end{vmatrix}=x^2(3a+x)$.

\item Using properties of determinants, show that: $\begin{vmatrix}1 & x & yz\\ 1 & y & zx\\ 1 & z & xy\end{vmatrix}=\begin{vmatrix}1 & x & x^2\\ 1 & y & y^2\\ 1 & z & z^2\end{vmatrix}$.

\item Using properties of determinants, show that: $\begin{vmatrix}1 & x & x^2\\ 1 & y & y^2\\ 1 & z & z^2\end{vmatrix}=(x-y)(y-z)(z-x)$.

\item Using properties of determinants, show that: $\begin{vmatrix}1 & 1 & 1\\ a+1 & b+1 & c+1\\ (a+1)^2 & (b+1)^2 & (c+1)^2\end{vmatrix}=(a-b)(b-c)(c-a)$.

\item Using properties of determinants, show that: $\begin{vmatrix}a^2+b^2 & c^2 & c^2\\ a^2 & b^2+c^2 & a^2\\ b^2 & b^2 & c^2+a^2\end{vmatrix}=4a^2b^2c^2$.

\item Using properties of determinants, show that: $\begin{vmatrix}a & b & c\\ b+c & c+a & a+b\\ a+b & b+c & c+a\end{vmatrix}=a^3+b^3+c^3-3abc$.

\item Using properties of determinants, show that: $\begin{vmatrix}a+t & a & a\\ b & b+t & b\\ c & c & c+t\end{vmatrix}=t^2(a+b+c+t)$.

\item Using properties of determinants, show that: $\begin{vmatrix}a-b-c & 2a & 2a\\ 2b & b-c-a & 2b\\ 2c & 2c & c-a-b\end{vmatrix}=(a+b+c)^3$.

\item Using properties of determinants, show that: $\begin{vmatrix}y+z & z+x & x+y\\ x & y & z\\ x^2 & y^2 & z^2\end{vmatrix}=(x+y+z)(x-y)(y-z)(z-x)$.

\item Using properties of determinants, show that: $\begin{vmatrix}1 & 1 & 1\\ a^2+1 & b^2+1 & c^2+1\\ a^3+a & b^3+b & c^3+c\end{vmatrix}=(a-b)(b-c)(c-a)(ab+bc+ca-1)$.

\item Using properties of determinants, show that: $\begin{vmatrix}1+a & 1 & 1\\ 1 & 1+b & 1\\ 1 & 1 & 1+c\end{vmatrix}=abc+ab+bc+ca$.

\item Using properties of determinants, show that: $\begin{vmatrix}1 & a & a^2-bc\\ 1 & b & b^2-ca\\ 1 & c & c^2-ab\end{vmatrix}=0$.

\item Using properties of determinants, show that: $\begin{vmatrix}2x+3 & x+2 & x+a\\ 2x+5 & x+3 & x+b\\ 2x+7 & x+4 & x+c\end{vmatrix}=0$, where $ab=a+c$.

\item Using properties of determinants, show that: $\begin{vmatrix}a & b & c\\ c & a & b\\ b & c & a\end{vmatrix}=(a+b+c)(a+b\omega+c\omega^2)(a+b\omega^2+c\omega)$, where $\omega$ is an imaginary cube root of unity.

\item If $A=\begin{pmatrix}1 & 2 & -3\\ 0 & -5 & 0\\ -2 & -2 & 7\end{pmatrix}$ and $B=\begin{pmatrix}-5 & -2 & 5\\ -3 & -1 & 4\\ -2 & -1 & 2\end{pmatrix}$, then find: $A_{13},A_{23},A_{33}$ and $|A|$.

\item If $A=\begin{pmatrix}1 & 2 & -3\\ 0 & -5 & 0\\ -2 & -2 & 7\end{pmatrix}$ and $B=\begin{pmatrix}-5 & -2 & 5\\ -3 & -1 & 4\\ -2 & -1 & 2\end{pmatrix}$, then find: $B_{31},B_{32},B_{33}$ and $|B|$.

\item Find the value of $x$ if: $\begin{vmatrix}2 & 1 & x\\ -1 & -4 & -3\\ x & 1 & 0\end{vmatrix}=5$.

\item Find the value of $x$ if: $\begin{vmatrix}1 & x-1 & 3\\ -1 & x+1 & 2\\ 2 & -3 & x\end{vmatrix}=9$.

\item Find $|AA^t|$ and $|A^tA|$: $A=\begin{pmatrix}3 & 1\\ 2 & 2\\ 1 & 3\end{pmatrix}$.

\item If $A$ is a square matrix of order 3, then show that $|KA|=K^3|A|$.

\item Find the values of $\lambda$ and $B$ (where $A$ and $B$ are singular): $A=\begin{pmatrix}4 & 2 & 3\\ 7 & \lambda & 6\\ 2 & 3 & 1\end{pmatrix}$, $B=\begin{pmatrix}-2 & 4 & 5\\ 1 & -2 & 1\\ 2 & \lambda & 0\end{pmatrix}$.

\item Find the inverse of $A=\begin{pmatrix}1 & 2 & 1\\ -5 & 0 & 4\\ 5 & 4 & 0\end{pmatrix}$ and show that $A^{-1}A=I_3$.

\item Verify that $(AB)^t=B^tA^t$ if: $A=\begin{pmatrix}1 & -1 & 2\\ 0 & -3 & 1\end{pmatrix}$ and $B=\begin{pmatrix}1 & 1\\ -3 & -2\\ 0 & 1\end{pmatrix}$.

\item Verify that $(AB)^t=B^tA^t$ if: $A=\begin{pmatrix}1 & 2\\ 1 & 4\\ 2 & 1\end{pmatrix}$ and $B=\begin{pmatrix}1 & -3\\ -2 & 1\end{pmatrix}$.

\item Find the inverse of the following matrix by using row operations: $\begin{pmatrix}2 & 6 & -3\\ 0 & -2 & 0\\ -2 & 5 & 6\end{pmatrix}$.

\item Find the inverse of the following matrix by using row operations: $\begin{pmatrix}1 & 2 & -1\\ 0 & -2 & 8\\ 1 & 0 & 2\end{pmatrix}$.

\item Find the inverse of the following matrix by using row operations: $\begin{pmatrix}1 & 6 & 2\\ 2 & 13 & 0\\ 0 & -1 & -1\end{pmatrix}$.

\item Find the rank of the following matrix: $\begin{pmatrix}1 & -1 & 3 & 1\\ -2 & -6 & 1 & -1\\ 3 & 1 & 4 & -2\end{pmatrix}$.

\item Find the rank of the following matrix: $\begin{pmatrix}1 & -2 & 3\\ 2 & -4 & 6\\ -1 & 0 & 2\\ 0 & 1 & -1\end{pmatrix}$.

\item Find the rank of the following matrix:
\[
\begin{bmatrix} 3 & -1 & 3 & 0 & 1 \\ 1 & 2 & -1 & -3 & -2 \\ 2 & 3 & 4 & 2 & 1 \\ 2 & 5 & -2 & -3 & 3 \end{bmatrix}
\]

\item Solve the following system of linear equations by Cramer's rule:
\[
2x+y-z=1, \qquad x-y+2z=3, \qquad 3x+2y+z=4
\]

\item Solve the following system of linear equations by Cramer's rule:
\[
x_1+2x_2-3x_3=0, \qquad 4x_1-x_2+x_3=5, \qquad -2x_1+3x_2+2x_3=3
\]

\item Solve the following system of linear equations by Cramer's rule:
\[
2x_1-x_2+x_3=1, \qquad x_1+2x_2+2x_3=2, \qquad x_1-2x_2-x_3=1
\]

\item Solve the following system of linear equations by matrix inversion method:
\[
x-2y+z=1, \qquad 3x+y-2z=4, \qquad y-z=1
\]

\item Solve the following system of linear equations by matrix inversion method:
\[
2x_1+x_2+3x_3=3, \qquad x_1+3x_2-2x_3=0, \qquad -3x_1-x_2+2x_3=4
\]

\item Solve the following system of linear equations by matrix inversion method:
\[
x+y=2, \qquad 2x-z=1, \qquad 2y-3z=-1
\]

\item Solve the following system by reducing the augmented matrix to echelon and reduced echelon form:
\[
x_1+2x_2-2x_3=-1, \qquad 2x_1+3x_2+x_3=1, \qquad 5x_1+4x_2-3x_3=1
\]

\item Solve the following system by reducing the augmented matrix to echelon and reduced echelon form:
\[
x+2y+z=2, \qquad 2x+y+2z=3, \qquad 2x+3y-z=7
\]

\item Solve the following system by reducing the augmented matrix to echelon and reduced echelon form:
\[
x_1+4x_2+x_3=2, \qquad 2x_1+x_2-2x_3=9, \qquad 3x_1+x_2-x_3=12
\]

\item Solve the following system by reducing the augmented matrix to echelon and reduced echelon form:
\[
x_1+4x_2+x_3=2, \qquad 2x_1+x_2-2x_3=9, \qquad 3x_1+x_2-x_3=12
\]

\item Solve the following homogeneous system of linear equations by using the Gaussian elimination method:
\[
x_1+4x_2+2x_3=0, \qquad 2x_1+x_2-3x_3=0, \qquad 3x_1+2x_2-4x_3=0
\]

\item Solve the following homogeneous system of linear equations by using the Gaussian elimination method:
\[
x_1+2x_2-x_3=0, \qquad x_1-x_2+5x_3=0, \qquad 2x_1+x_2+4x_3=0
\]

\item A triangle has vertices at $A(4,1)$, $B(-2,5)$ and $C(0,-3)$. Find the vertices of the reflected triangle over the $y$-axis using a transformation matrix.

\item The point $A$ is mapped to $(30,20,-5)$ by the scaling matrix $\begin{bmatrix}-5&0&0\\0&-5&0\\0&0&-5\end{bmatrix}$. Find the coordinates of $A$.

\item Find the equation of the image of the curve with equation $y=x^2$ under the transformation with associated matrix $\begin{bmatrix}1&2\\3&4\end{bmatrix}$.

\item Use the matrix $\begin{bmatrix}1&0&1\\2&-1&1\\0&1&2\end{bmatrix}$ to encode the message: KEEP IT UP, where letters A to Z correspond to the numbers 1 to 26.

\item Decode the message $\begin{bmatrix}11\\20\\43\end{bmatrix}\begin{bmatrix}25\\10\\41\end{bmatrix}\begin{bmatrix}22\\14\\41\end{bmatrix}$ that was encoded using matrix $A=\begin{bmatrix}1&1&-1\\1&0&1\\2&1&1\end{bmatrix}$, where the numbers 1 to 26 correspond to the letters A to Z, and 27 represents space or ``-''.

\item Reduce $\begin{bmatrix}2&3&-1&9\\1&-1&2&-3\\3&1&3&2\end{bmatrix}$ to row echelon and reduced row echelon form.

\item Find the rank of the matrix $\begin{bmatrix}1&-1&2&-3\\2&0&7&-7\\3&1&12&-11\end{bmatrix}$.

\item Solve the following and explain whether the system is consistent or inconsistent:
\[
2x+5y-z=5, \qquad 3x+4y+2z=11, \qquad x+2y-2z=-3
\]

\item Solve the following and explain whether the system is consistent or inconsistent:
\[
x+y+2z=1, \qquad 2x-y+7z=1, \qquad 3x+5y+4z=-3
\]

\item Solve the following and explain whether the system is consistent or inconsistent:
\[
x-y+2z=1, \qquad 2x-6y+5z=1, \qquad 3x+5y+4z=-3
\]

\item Use the matrix inversion method to solve the system:
\[
x_1-2x_2+x_3=-4, \qquad 2x_1-3x_2+2x_3=-6, \qquad 2x_1+2x_2+x_3=5
\]

\item Use Cramer's rule to solve the system:
\[
3x_1+x_2-x_3=-4, \qquad x_1+x_2-2x_3=-4, \qquad -x_1+2x_2-x_3=1
\]

\item Solve the following system of equations by the Gaussian elimination method:
\[
x+2y+z=0, \qquad 2x+3y+4z=0, \qquad 4x+3y+2z=0
\]

\item Solve the following system of equations using the Gaussian elimination method:
\[
x_1+x_2+x_3=0, \qquad x_1-x_2+3x_3=0, \qquad x_1+3x_2-x_3=0
\]

\item A triangle has $A(2,3)$, $B(-1,4)$ and $C(3,-2)$. Find the vertices of the reflected triangle over the $x$-axis using a transformation matrix.

\item Use matrix $A=\begin{bmatrix}1&2\\3&1\end{bmatrix}$ to encode the message: ATTACK, where letters A to Z correspond to the numbers 1 to 26.

\item Resolve the following into partial fractions: $\dfrac{a-b}{(x-a)(x-b)}$, $(a\neq b)$.

\item Resolve the following into partial fractions: $\dfrac{x^2+1}{(x+1)(x-1)}$.

\item Resolve the following into partial fractions: $\dfrac{2x+3}{(x+1)(x+2)(x+3)}$.

\item Resolve the following into partial fractions: $\dfrac{x^2+4x+5}{(x+1)(x^2+5x+6)}$.

\item Resolve the following into partial fractions: $\dfrac{4x^3+5x^3-3x-2}{x^2-1}$.

\item Resolve the following into partial fractions: $\dfrac{3x^2-12x+11}{(x-1)(x-2)(x-3)}$.

\item Resolve the following into partial fractions: $\dfrac{(x-1)(x-2)(x-3)}{(x-4)(x-5)(x-6)}$.

\item Resolve the following into partial fractions: $\dfrac{x^2}{(x^2+a)(x^2+b)(x^2+c)}$.

\item Resolve the following into partial fractions: $\dfrac{x+1}{(x-1)^2}$.

\item Resolve the following into partial fractions: $\dfrac{x^2+x}{(x^2-1)^2}$.

\item Resolve the following into partial fractions: $\dfrac{3x^2+4x-5}{(x-1)^3}$.

\item Resolve the following into partial fractions: $\dfrac{1}{x(x+1)^3}$.

\item Resolve the following into partial fractions: $\dfrac{4x^2-3x+1}{(x+1)(x-1)^2}$.

\item Resolve the following into partial fractions: $\dfrac{12x^2-48}{(x-2)^2(x+2)^2}$.

\item Resolve $\dfrac{x^2-10x+13}{(x-1)(x^2-5x+6)}$ into partial fractions.

\item Resolve $\dfrac{2x^3+x^2-x-3}{x(2x+3)(x-1)}$ into partial fractions.

\item Resolve $\dfrac{x^2+x-1}{(x+2)^3}$ into partial fractions.

\item Resolve $\dfrac{1}{(x+1)^2(x^2-1)}$ into partial fractions.

\item Resolve into partial fractions: $\dfrac{2x^2+3x+3}{(x+1)(x^2+1)}$.

\item Resolve into partial fractions: $\dfrac{2x+1}{(x-2)(x^2+3x+5)}$.

\item Resolve into partial fractions: $\dfrac{2x+32}{(x-2)(x^2+2)^2}$.

\item Resolve into partial fractions: $\dfrac{3x^2+3}{x^3+1}$.

\item Resolve into partial fractions: $\dfrac{x^4}{x^4+2x^2+1}$.

\item Resolve into partial fractions: $\dfrac{6x^4+40x^2}{(4-x^2)(x^2+4)^2}$.

\item Resolve $\dfrac{3x-11}{(x^2+1)(x+3)}$ into partial fractions.

\item Resolve $\dfrac{4x^2+8x}{x^4+2x^2+9}$ into partial fractions.

\item Resolve $\dfrac{4x^2}{(x^2+1)^2(x-1)}$ into partial fractions.

\item The 18th and 30th terms of an arithmetic sequence are 367 and 499 respectively. How many terms of this sequence are less than 1000?

\item The 7th and 21st terms of an A.P. are 37 and 107 respectively. Find the A.P. and its 100th term.

\item If $\dfrac{1}{a-c}, \dfrac{1}{b-c}, \dfrac{1}{b-a}$ are in A.P., show that $\dfrac{a-b}{a-c}=\dfrac{a-c}{b-a}$.

\item Find the $n$th term of the progression $\left(\dfrac{3}{7}\right)^{10}, \left(\dfrac{10}{7}\right)^{10}, \left(\dfrac{17}{7}\right)^{10}, \ldots$. Is the progression an A.P.?

\item If the arithmetic progressions $3,10,17,\ldots$ and $63,65,67,\ldots$ are such that their $n$th terms are equal, find the value of $n$.

\item If $\dfrac{1}{a}, \dfrac{1}{b}$ and $\dfrac{1}{c}$ are in A.P., show that the common difference is $\dfrac{a-c}{2ac}$.

\item If $a_k$ and $a_m$ denote two different terms of an A.P., show that its $n$th term is
\[
a_k+(n-k)\left(\frac{a_k-a_m}{k-m}\right).
\]

\item If $a_1,a_2,a_3,\ldots,a_n$ are positive and in A.P., prove that:
\[
\frac{1}{\sqrt{a_1}+\sqrt{a_2}}+\frac{1}{\sqrt{a_2}+\sqrt{a_3}}+\cdots+\frac{1}{\sqrt{a_{n-1}}+\sqrt{a_n}}=\frac{n-1}{\sqrt{a_1}+\sqrt{a_n}}
\]

\item If the roots of the equation $(b-c)x^2+(c-a)x+(a-b)=0$ are equal, show that $a,b,c$ are in A.P.

\item If the sides of a right-angled triangle are in A.P., find the ratio of its sides.

\item If the $n$th term of a progression is a linear expression in $n$, then prove that this progression is an A.P.

\item Find the general term and the eleventh term of the A.P. whose first term and common difference are 2 and $-3$ respectively. Also write its first four terms.

\item Insert five arithmetic means between $\sqrt2$ and $\dfrac{15}{\sqrt2}$.

\item The A.M. of two numbers is 7 and their product is 45. Find the numbers.

\item If $n$ arithmetic means are inserted between $a$ and $b$, prove that $d=\dfrac{b-a}{n+1}$, where $d$ is the common difference.

\item If $A$ is the A.M. between $a$ and $b$, prove that $(a-A)^2+(A-b)^2=\dfrac{1}{2}(a-b)^2$.

\item For what value of $n$, $\dfrac{a^{n+1}+b^{n+1}}{a^n+b^n}$ is the A.M. between $a$ and $b$, where $a\neq b$?

\item The 5th and 9th terms of an A.P. are 11 and 17 respectively. Find the sum of 20 terms.

\item Obtain the sum of all integers in the first 1000 positive integers which are neither divisible by 5 nor by 2.

\item The sum of 9 terms of an A.P. is 171 and its eighth term is 31. Find the series.

\item The 5th term of an arithmetic progression is 21 and the sum of the first six terms is 90. Find the 18th term.

\item The sum of three numbers in an A.P. is 24 and their product is 440. Find the numbers.

\item The first four terms of an A.P. are 2, 6, 10 and 14. Find the least number of terms needed so that the sum of the terms is greater than 2000.

\item Find four numbers in A.P. whose sum is 32 and the sum of whose squares is 276.

\item Find the five numbers in A.P. whose sum is 25 and the sum of whose squares is 135.

\item If $\dfrac{1}{a+b}, \dfrac{1}{c+a}, \dfrac{1}{b+c}$ are in A.P. then show that $a^2, b^2, c^2$ are in A.P.

\item The sum of the first four terms of an A.P. is 56. The sum of the last four terms is 112. If its first term is 11, then find the number of terms.

\item The first, second and last terms of an A.P. are $a$, $b$ and $c$ respectively. Show that the sum of the A.P. is
\[
\frac{(b+c-2a)(c+a)}{2(b-a)}
\]

\item Show that the sum of $n$ arithmetic means between $a$ and $b$ is $n$ times the single A.M. between them.

\item Find the 19th term and the partial sum of 19 terms of the arithmetic series: $2+\dfrac{7}{2}+5+\dfrac{13}{2}+\cdots$

\item If the 4th and 9th terms of a G.P. are 54 and 13122 respectively, find the G.P. Also find its general term.

\item If the $(p+q)$th term of a G.P. is $m$ and the $(p-q)$th term is $n$, then find the $p$th term.

\item Find three consecutive numbers in G.P. whose sum is 26 and their product is 216.

\item The 3rd term of a G.P. is the square of the 1st term. If the 2nd term is 9, then find the 6th term.

\item If $\dfrac{1}{a}, \dfrac{1}{b}$ and $\dfrac{1}{c}$ are in G.P., show that the common ratio is $\pm\sqrt{\dfrac{a}{c}}$.

\item If the numbers 1, 4 and 3 are subtracted from three consecutive terms of an A.P., the resulting numbers are in G.P. Find the original numbers if their sum is 21.

\item If three consecutive numbers in A.P. are increased by 1, 4, 15 respectively, the resulting numbers are in G.P. Find the original numbers if their sum is 6.

\item If the $p$th, $q$th term of a G.P. are $q$ and $p$ respectively, show that the $(p+q)$th term is
\[
(q^p \div p^q)^{\frac{1}{p-q}}
\]

\item If both $x$ and $y$ are positive distinct real numbers, show that the geometric mean between $x$ and $y$ is less than their arithmetic mean.

\item For what value of $n$ is $\dfrac{a^n+b^n}{a^{n-1}+b^{n-1}}$ the positive geometric mean between $a$ and $b$?

\item The A.M. of two positive integral numbers exceeds their (positive) G.M. by 2, and their sum is 20. Find the numbers.

\item The A.M. between two numbers is 5 and their (positive) G.M. is 4. Find the numbers.

\item The arithmetic mean between two positive numbers $a$ and $b$ is double their geometric mean. Prove that $a:b = (2+\sqrt3):(2-\sqrt3)$.

\item If one geometric mean $G$ and two arithmetic means $p$ and $q$ are inserted between two positive numbers, show that $G^2=(2p-q)(2q-p)$.

\item The 3rd term of a G.P. is 16 and the 6th term is $-128$. Find the first term and the sum of the first seven terms.

\item Sum of $n$ terms of the series: $1+(a+b)+(a^2+ab+b^2)+(a^3+a^2b+ab^2+b^3)+\cdots$

\item Show that the ratio of the sum of the first $n$ terms of a G.P. to the sum of terms from $(n+1)$th to $(2n)$th term is $\dfrac{1}{r^n}$, where $r$ is the common ratio of the G.P.

\item Find the 8th term of the arithmetico-geometric sequence, where the arithmetic part is $1, 4, 7, \ldots$ and the geometric part is $5, 10, 20, \ldots$

\item Find the $n$th term of the arithmetico-geometric sequence, where the arithmetic part is $3, 7, 11, \ldots$ and the geometric part is $2, 6, 18, \ldots$

\item Consider the arithmetico-geometric sequence defined by the arithmetic part $a_{n+1}=2n+5$ and the geometric part $b_{n-2}=\dfrac{1}{9}(-3)^n$. Find the $n$th term and the sum of the first three terms of the arithmetico-geometric sequence.

\item Sum of $n$ terms of the following series: $2+\dfrac{5}{4}+\dfrac{8}{4^2}+\dfrac{11}{4^3}+\cdots$

\item Sum of $n$ terms of the following series: $1+\dfrac{3}{5}+\dfrac{5}{5^2}+\dfrac{7}{5^3}+\cdots$

\item Sum of $n$ terms of the following series: $1+\dfrac{4}{3}+\dfrac{7}{9}+\dfrac{10}{27}+\cdots$

\item Sum the following infinite series: $1+\dfrac{3}{2}+\dfrac{5}{4}+\dfrac{7}{8}+\cdots$

\item Sum the following infinite series: $2+\dfrac{5}{3}+\dfrac{8}{9}+\dfrac{11}{27}+\cdots$

\item Show that $2^{1/2}\cdot4^{1/4}\cdot8^{1/8}\cdot16^{1/16}\cdots\infty = 4$.

\item Find the sum of $n$ terms of the series
\[
\frac{2n+1}{2n-1}+3\left(\frac{2n+1}{2n-1}\right)^2+5\left(\frac{2n+1}{2n-1}\right)^3+\cdots
\]

\item Prove that $1+2\left(1+\dfrac{1}{n}\right)+3\left(1+\dfrac{1}{n}\right)^2+\cdots$ to $n$ terms $=n^2$.

\item Sum the series to $n$ terms $2+5x+8x^2+11x^3+\cdots$ and deduce the sum to infinity if $|x|<1$.

\item Find the sum of $n$ terms of the series $1+2x+3x^2+4x^3+\cdots$, where $x\neq1$. If $|x|<1$, sum the series to infinity.

\item If the numbers $\dfrac{1}{k}, \dfrac{1}{2k+1}$ and $\dfrac{1}{4k-1}$ are in harmonic sequence, find $k$.

\item Find $n$ so that $\dfrac{a^{n+1}+b^{n+1}}{a^n+b^n}$ may be the H.M. between $a$ and $b$.

\item If the H.M. and A.M. between two numbers are 4 and $\dfrac{9}{2}$ respectively, find the numbers.

\item If the (positive) G.M. and H.M. between two numbers are 4 and $\dfrac{16}{5}$, find the numbers.

\item If $\dfrac{b+c-a}{a}, \dfrac{c+a-b}{b}, \dfrac{a+b-c}{c}$ are in A.P., show that $a,b,c$ are in H.P.

\item If between any two numbers there are inserted two A.Ms. $A_1,A_2$, two G.Ms. $G_1,G_2$ and two H.Ms. $H_1,H_2$, show that
\[
\frac{A_1+A_2}{G_1G_2}=\frac{H_1+H_2}{H_1H_2}
\]

\item The H.M. of two numbers is 4. The A.M., $A$, and the G.M., $G$, satisfy the relation $2A+G^2=27$. Find the numbers.

\item The first three of four numbers $a,b,c,d$ are in A.P., and the next three (last three) are in H.P. Show that $ad=bc$.

\item If the roots of the equation $a(b-c)x^2+b(c-a)x+c(a-b)=0$ are equal, prove that $a,b,c$ are in H.P.

\item If the 4th and 7th terms of an H.P. are $\dfrac{2}{13}$ and $\dfrac{2}{25}$ respectively, find the sequence.

\item Sum the series: $\dfrac{1^2}{1}+\dfrac{1^2+2^2}{2}+\dfrac{1^2+2^2+3^2}{3}+\cdots$ to $n$ terms.

\item Find the sum to $n$ terms of the series whose $n$th term is given: $5n^2+2n+3$.

\item Find the sum to $n$ terms of the series whose $n$th term is given: $n^2+2n-3$.

\item Given the $n$th terms of the series, find the sum to $2n$ terms: $3n^2+5n+2$.

\item Given the $n$th terms of the series, find the sum to $2n$ terms: $n^2+n-2$.

\item A sum of Rs.~10400 is paid off in 40 installments such that each installment is Rs.~10 more than the preceding installment. Calculate the value of the first installment.

\item An investor invests Rs.~150000 at an annual compound interest rate of 6\% for 8 years. Find the total amount he will get after 8 years.

\item The population of a town is 4084101 at present and five years ago it was 3200000. Find its rate of increase if it increased geometrically.

\item Determine the total worth of a yearly Rs.~5000 investment after 20 years if the interest rate is 5\% compounded annually.

\item A water tank has a leakage. Each week, the tank loses 5 gallons of water due to the leakage. Initially, the tank is full and contains 2000 gallons. (i) How many gallons are in the tank 20 weeks later? (ii) How many weeks until the tank is half-full? (iii) How many weeks until the tank is empty?

\item A drug company has manufactured 7 million doses of a vaccine to date. They promise additional production at a rate of 1.4 million doses/month over the next year. (i) How many doses of the vaccine, in total, will have been produced after a year? (ii) Give the general term $a_n$ describing the total number of doses of the vaccine produced, and describe the meaning of the variable $n$ in the context of this problem. Find the general term $a_n$. (iii) Find the value of $a_{10}$ and interpret its meaning in words.

\item At a toll booth, the number of vehicles passing through during the first minute is 100. Due to road congestion, each minute only 80\% of the vehicles from the previous minute manage to pass. (i) Represent the number of vehicles passing each minute as a sequence. (ii) Find the total number of vehicles that pass through in 15 minutes. (iii) What is the maximum number of vehicles that can pass in the long run (as $t\to\infty$)?

\item A sum of Rs.~5000 is invested at 8\% simple interest per year. Calculate the interest at the end of each year. Do these interests form an A.P.? If so, find the interest at the end of 20 years making use of this fact.

\item A machine is purchased for Rs.~20,000. It depreciates at 6\% per annum for the first four years and after that 8\% per annum for the next six years, depreciation being calculated on diminishing value. Find the value of the machine after a period of 10 years.

\item Two cars start together in the same direction from the same place. The first goes with a uniform speed of 20~km/h. The second goes at a speed of 12~km/h in the first hour and increases its speed by 1~km/h each succeeding hour. After how many hours will the second car overtake the first car if both cars go non-stop?

\item 150 workers were engaged to finish a piece of work in a certain number of days. Five workers dropped out on the second day, five more dropped out on the third day, and so on. It takes 10 more days to finish the work now. Find the number of days in which the work was completed.

\item A radioactive product has a half-life of 5 years. If the radioactivity level is 68 microcuries after 20 years, determine the original level of radioactivity.

\item An object moving in a line is given an initial velocity of 4.5~m/s and a constant acceleration of $2.5$~m/s$^2$. How long will it take the object to reach a velocity of 20~m/s?

\item In an integrated circuit with an initial current of 1080~mA, the temperature in the components decreases from 20\% to 17\%. Assuming that each temperature decrease is caused by a decrease in the initial current, what is the value of the current at the fourth measurement?

\item Show that the maximum amount of a certain sum of money at compound interest of $r\%$ per year for $n$ years forms a G.P.

\item Vehicles arrive at a toll booth at a rate of 4 cars every 5 minutes. Represent the number of cars arriving over time as a sequence and predict the total number of cars after an hour.

\item A web designer is using a 12-column grid system where each column increases in width by 10px from the previous one. The first column width is 50px wide. Find the total width occupied by all 12 columns.

\item A company leases a motor vehicle with the following terms: (i) The first monthly payment is Rs.~15,000. (ii) Each subsequent payment increases by Rs.~500 due to inflation adjustments. (iii) The lease term is 24 months. Find: (i) the payment on the 24th month; (ii) the total amount paid over 24 months; (iii) if the company can only afford to pay a total of Rs.~400,000, can they complete the 24-month lease? (iv) Find the maximum months $n$ such that the total payment $S_n \le 400{,}000$.

\item Prove that:
\[
\frac{(2n+1)!}{n!}=[1\cdot3\cdot5\cdots(2n-1)(2n+1)]\,2^n
\]

\item Express as a single fraction: $\dfrac{(n+2)!}{(r+2)!}+\dfrac{(n+1)!}{(r+1)!}$

\item There are four distinct colored balls and four boxes of the same colors as those of the balls. Determine the number of possible ways the balls, one each in a box, can be placed such that a ball does not go to a box of its own colour.

\item Prove from first principles that: ${}^{n}P_r={}^{n-1}P_r+r\cdot{}^{n-1}P_{r-1}$

\item How many numbers between 100 and 1000 can be formed using the digits $0,1,2,3,4,5$ without repetition? How many of them are divisible by 5?

\item Find the number of numbers greater than 35000 that can be formed from the digits $1,2,3,4,5,6$ without repeating any digit.

\item Find the number of 5-digit numbers that can be formed from the digits $1,2,4,6,8$ (when no digit is repeated), but (i) the digits 2 and 8 are next to each other; (ii) the digits 2 and 8 are not next to each other.

\item How many 6-digit numbers can be formed, without repeating any digit, from the digits $0,1,2,3,4,5$? In how many of them will 0 be at the tens place?

\item In how many ways can 8 different books, including 2 on English, be arranged on a shelf in such a way that the English books are never together?

\item Find the number of arrangements of 3 different books on English and 5 different books on Urdu for placing them on a shelf such that the books on the same subject are together.

\item In how many ways can 5 boys and 4 girls be seated on a bench so that the girls and the boys occupy alternate seats?

\item In how many ways can a set of 4 different mathematics books, 3 different physics books and 2 different chemistry books be placed on a shelf with a space for 9 books if: (a) all the books are kept without any restriction; (b) all the books of the same subject are kept together; (c) only the mathematics books are kept together?

\item In how many ways can 5 people be seated on a bench if: (a) there are no restrictions; (b) two particular people must sit next to each other; (c) two particular people cannot sit next to each other?

\item Abdullah has a collection of 9 marbles consisting of 4 identical red marbles, 3 identical blue marbles and 2 identical green marbles. If he wants to arrange all of them in a straight row, how many distinct arrangements are possible?

\item In how many different ways can the following persons sit around a round table? (a) 8 persons (b) 7 persons (c) 6 persons.

\item 15 members of a club are to form 4 committees of 3, 5, 4 and 3 members so that no member is a member of more than one committee. Find the number of ways of forming the committees.

\item The D.C.Os of 11 districts meet to discuss the law-and-order situation in their districts. In how many ways can they be seated at a round table, when two particular D.C.Os insist on sitting together?

\item The Governor of the Punjab calls a meeting of 14 officers. In how many ways can they be seated at a round table?

\item Fatima invites 14 people for dinner. There are 9 males and 5 females who are seated at two different tables. Guests of one sex sit at one round table and the guests of the other sex sit at the second table. Find the number of ways in which all guests can be seated.

\item Find the number of ways in which 5 men and 5 women can be seated at a round table in such a way that no two persons of the same sex sit together.

\item In how many ways can 4 persons be seated at a round table, while (a) clockwise and anticlockwise orders are different; (b) clockwise and anticlockwise orders are identical?

\item Prove that: ${}^{n}C_r + {}^{n}C_{r-1} = {}^{n+1}C_r$

\item Prove that: $r\cdot{}^{n}C_r = (n-r+1)\,{}^{n}C_{r-1}$

\item In how many ways can five subjects be selected out of eight subjects to select a course programme?

\item In how many ways can 3 dishes of Desi food and 2 dishes of Chinese food be selected from 6 dishes of Desi food and 8 dishes of Chinese food?

\item From a standard deck of 52 playing cards, there are 26 black cards and 26 red cards. How many different ways can eight cards be selected if 3 are black and the remaining 5 are red?

\item A bag contains 8 red balls and 7 green balls. Find the total number of possible ways in which five balls are selected such that: (i) 3 red and 2 green; (ii) 1 red and 4 green; (iii) 4 red and 1 green; (iv) all the red balls.

\item The members of a club are 10 boys and 8 girls. In how many ways can a committee of 6 boys and 3 girls be formed?

\item How many committees of 7 members can be chosen from a group of 10 persons when each committee must include 2 particular persons?

\item In how many ways can a cricket team of 11 players be selected out of 17 players? How many of them will include a particular player?

\item There are 6 men and 8 women members of a club. How many committees of seven can be formed: (i) with 3 women; (ii) with at most 3 women; (iii) with at least 5 women?

\item There are 3 sections in a question paper, each section has 3 questions. A student has to solve 5 questions in total, choosing at least one question from each section. In how many ways can the student make this choice?

\item Consider a cryptographic system that generates an 8-character password. Each character can be either a lowercase letter (a--f) or a digit (0--9). How many passwords can be generated if each password must contain exactly 5 lowercase letters and 3 digits: (a) with repetition allowed (b) without repetition.

\item On Defence Day, Teacher I compiles a list of 10 distinct national songs, while Teacher II prepares a separate list of 10 different national songs (with no overlap between the two lists). The principal needs to select 3 songs from Teacher I's list and 3 songs from Teacher II's list. Determine the number of possible selection methods when: (i) the order/sequence of the selected songs is important; (ii) the order/sequence of the selected songs is not important.

\item In a school, a class consists of 12 girls and 8 boys. The teacher wants to select 5 students for an activity. In how many ways can the students be selected including: (i) 2 girls; (ii) 5 boys; (iii) 2 boys.

\item Zain wants to generate a password for his laptop. He can use only 6 characters, each of which can be either an uppercase letter (A--Z) or a digit (0--9). How many passwords can be generated: (i) if repetition of characters is not allowed; (ii) if repetition of characters is allowed.

\item A box contains 15 cards numbered 1--15. Danish is to select 5 cards. If all the selected cards are the first five multiples of 2, Danish wins the game. Find Danish's chance of winning the game, when (i) order is important; (ii) order is not important.

\item On Independence Day, a DJ has a list of ten different national songs. He wants to select any five national songs for the day. Find how many ways he can select and play the songs: (i) if the order of playing the songs matters; (ii) if the order of playing the songs does not matter.

\item Use mathematical induction to prove the following formula for every positive integer $n$: $2+5+8+\cdots+(3n-1)=\dfrac{n}{2}(3n+1)$

\item Use mathematical induction to prove the following formula for every positive integer $n$: $2+(2+5)+(2+5+8)+\cdots+\dfrac{n}{2}(3n+1)=\dfrac{n}{4}(n+1)^2$

\item Use mathematical induction to prove the following formula for every positive integer $n$: $2+6+18+\cdots+2\times 3^{n-1}=3^n-1$

\item Use mathematical induction to prove the following formula for every positive integer $n$: $1\times3+2\times5+3\times7+\cdots+n\times(2n+1)=\dfrac{n(n+1)(4n+5)}{6}$

\item Use mathematical induction to prove the following formula for every positive integer $n$: $\dfrac{1}{1\times2}+\dfrac{1}{2\times3}+\dfrac{1}{3\times4}+\cdots+\dfrac{1}{n(n+1)}=1-\dfrac{1}{n+1}$

\item Use mathematical induction to prove the following formula for every positive integer $n$: $r+r^2+r^3+\cdots+r^n=\dfrac{r(1-r^n)}{1-r}$, $(r\neq 1)$

\item Use mathematical induction to prove the following formula for every positive integer $n$: $a_n=a_1+(n-1)d$ when $a_1,\ a_1+d,\ a_1+2d,\ldots$ form an A.P.

\item Use mathematical induction to prove the following formula for every positive integer $n$: $a_n=a_1r^{n-1}$ when $a_1,\ a_1r,\ a_1r^2,\ldots$ form a G.P.

\item Use mathematical induction to prove the following formula for every positive integer $n$: $\dbinom{3}{3}+\dbinom{4}{3}+\dbinom{5}{3}+\cdots+\dbinom{n+2}{3}=\dbinom{n+3}{4}$

\item Prove by mathematical induction that for all positive integral values of $n$: $8\times 10^{n-1}-2$ is divisible by 6.

\item Prove that $\displaystyle\sum_{k=1}^{n} r^k=\dfrac{r^{n+1}-1}{r-1}$, whenever $n$ is a positive integer.

\item Aliza invests Rs.\ 1,000,000 in a business that promises a 6\% return compounded annually. Prove by mathematical induction that the amount of money after $n$ years is $1{,}000{,}000(1.06)^n$.

\item A bank offers an investment with an annual interest rate $r$. If $P$ rupees are invested, the amount after $n$ years is given by $A(n)=P(1+r)^n$. Prove by induction that this formula holds for all $n\geq 0$.

\item Sikandar saves Rs.\ 500 in the first month and increases his savings by Rs.\ 500 every subsequent month. Using mathematical induction, determine whether his total savings will reach at least Rs.\ 12,000 after 24 months.

\item Prove by mathematical induction that if Ali takes a loan of Rs.\ 2,000,000 and pays Rs.\ 50,000 at the end of each year, the remaining balance after $n$ years is $R_n=2{,}000{,}000-50{,}000n$.

\item If Salman starts saving with Rs.\ 5,000 and saves an additional Rs.\ 1,000 at the end of every month, derive a formula $S(n)$ for his total savings after $n$ months. Prove the correctness of the formula using mathematical induction.

\item Use mathematical induction to prove that for any positive integer $n$: $1^2+2^2+3^2+\cdots+n^2=\dfrac{n(n+1)(2n+1)}{6}$

\item Use mathematical induction to prove that $3+3.5+3.5^2+\cdots+3.5^n=\dfrac{3(5^{n+1}-1)}{4}$, whenever $n$ is a non-negative integer.

\item Use mathematical induction to prove that $\displaystyle\sum_{k=1}^{n}\dfrac{1}{(2k-1)(2k+1)}=\dfrac{n}{2n+1}$, whenever $n$ is a positive integer.

\item Faris starts a saving plan where he deposits 1,000 rupees into his bank account every month. Using mathematical induction, prove that the total amount saved after $n$ months is given by $S(n)=1000n$ rupees.

\item Ali starts a daily exercise routine where each day he increases the number of push-ups he does by 2. On the first day, he does 10 push-ups. Prove that after $n$ days, the total number of push-ups Ali has done is $n^2+9n$.

\item Suppose you aim to lose weight by reducing your calorie intake by 50 calories each week. If you start at 2500 calories, prove that after $n$ weeks, your daily intake is $2500-50n$ calories.

\item Find the term independent of $x$ in the following expansion: $\left(\sqrt{x}+\dfrac{1}{2x^2}\right)^{10}$

\item Find the term independent of $x$ in the following expansion: $(1+x^2)^3\left(1+\dfrac{1}{x^2}\right)^4$

\item Determine the middle term in the following expansion: $\left(\dfrac{1}{x}-\dfrac{x^2}{2}\right)^{12}$

\item Determine the middle term in the following expansion: $\left(\dfrac{3}{2}x-\dfrac{1}{3x}\right)^{11}$

\item Determine the middle term in the following expansion: $\left(2x-\dfrac{1}{2x}\right)^{2m+1}$

\item Show that: $\dbinom{n}{1}+\dbinom{n}{2}+\dbinom{n}{3}+\cdots+\dbinom{n}{n}=2^n-1$

\item Show that: $\dbinom{n}{1}+2\dbinom{n}{2}+3\dbinom{n}{3}+\cdots+n\dbinom{n}{n}=n\cdot 2^{n-1}$

\item Find the term involving $a^4$ in the expansion of $\left(\dfrac{2}{x}-a\right)^9$

\item Expand the following up to 4 terms, taking values of $x$ for which the expansion is valid: $\dfrac{(1-x)^{-1}}{(1+x)^2}$

\item If $x$ is so small that its square and higher powers can be neglected, then show that: $\sqrt{\dfrac{1-x}{1+x}}\approx 1-\dfrac{3}{2}x$

\item If $x$ is so small that its square and higher powers can be neglected, then show that: $\sqrt{\dfrac{1+2x}{1-x}}\approx 1+\dfrac{3}{2}x$

\item If $x$ is so small that its square and higher powers can be neglected, then show that: $\dfrac{(9+7x)^{1/2}-(16+3x)^{1/4}}{4+5x}\approx \dfrac{1}{4}-\dfrac{17}{384}x$

\item If $x$ is so small that its square and higher powers can be neglected, then show that: $\dfrac{\sqrt{4+x}}{(1-x)^3}\approx 2+\dfrac{25}{4}x$

\item If $x$ is so small that its cube and higher powers can be neglected, then show that: $\sqrt{1-x-2x^2}\approx 1-\dfrac{1}{2}x-\dfrac{9}{8}x^2$

\item If $x$ is so small that its cube and higher powers can be neglected, then show that: $\sqrt{\dfrac{1+x}{1-x}}\approx 1+x+\dfrac{1}{2}x^2$

\item If $x$ is very nearly equal to 1, then prove that $px^p-qx^q\approx (p-q)x^{p+q}$

\item Identify the following series as a binomial expansion and find the sum: $1-\dfrac{1}{2}\left(\dfrac{1}{4}\right)+\dfrac{1.3}{2!4}\left(\dfrac{1}{4}\right)^2-\dfrac{1.3.5}{3!8}\left(\dfrac{1}{4}\right)^3+\cdots$

\item Use the binomial theorem to show that $1+\dfrac{1}{4}+\dfrac{1.3}{4.8}+\dfrac{1.3.5}{4.8.12}+\cdots=\sqrt{2}$

\item If $y=\dfrac{1}{3}+\dfrac{1.3}{2!}\left(\dfrac{1}{3}\right)^2+\dfrac{1.3.5}{3!}\left(\dfrac{1}{3}\right)^3+\cdots$, prove that $y^2+2y-2=0$

\item If $2y=\dfrac{1}{2^2}+\dfrac{1.3}{2!}\cdot\dfrac{1}{2^4}+\dfrac{1.3.5}{3!}\cdot\dfrac{1}{2^6}+\cdots$, prove that $4y^2+4y-1=0$

\item Show that the coefficient of $x^r$ in $\dfrac{x}{(1-px)(1-qx)}$ is $\dfrac{p^r-q^r}{p-q}$

\item If $x$ is so small that its cube and higher powers can be neglected, show that $\sqrt{\dfrac{1-x}{1+x}}\approx 1-x+\dfrac{1}{2}x^2$.

\item For $y=\dfrac{1}{2}\left(\dfrac{4}{9}\right)+\dfrac{1.3}{2^2\cdot 2!}\left(\dfrac{4}{9}\right)^2+\dfrac{1.3.5}{2^3\cdot 3!}\left(\dfrac{4}{9}\right)^3+\cdots$, show that $5y^2+10y-4=0$

\item Using the binomial theorem, show that: $(31)^4+(29)^6$ is divisible by 30.

\item Using the binomial theorem, show that: $(101)^5+(99)^7$ is divisible by 100.

\item Using the binomial theorem, find the remainder when $3^{101}$ is divided by 8.

\item Using the binomial theorem, show that: $(31)^4+(29)^6$ is divisible by 30.

\item Using the binomial theorem, show that $7^n-6n$ leaves remainder 1 when divided by 6, for all positive integers $n$.

\item Show that $a^2+(a+2)^2+(a+4)^2+1$ is divisible by 12, whenever ``$a$'' is an odd integer.

\item A company expects its annual revenue to grow at a fixed rate of 6\% per year. The revenue in year 1 is Rs.\ 10,000,000. Estimate the company's revenue after 4 years using the binomial theorem for small growth rates.

\item A bank offers a compound interest rate of 10\% per year. Zafar invests Rs.\ 2,000,000 for 4 years. How much will his investment be worth at the end of 4 years?

\item Zaid is organizing a sports competition with 8 teams. Every team plays against every other team exactly once. How many matches will be played in total? Use Pascal's triangle to solve this.

\item Expand $(1-2x)^{1/3}$ to four terms and apply it to evaluate $(0.8)^{1/3}$ correct to three decimal places.

\item A bank offers a compound interest rate of 5\% per year. Sumaira invests Rs.\ 100,000 for 3 years. How much will her investment be worth at the end of 3 years?

\item Use the remainder theorem to find the remainder when the first polynomial is divided by the second polynomial: $x^2+5x+6$, $x-2$

\item Use the remainder theorem to find the remainder when the first polynomial is divided by the second polynomial: $x^3+5x^2+6$, $x+1$

\item Use the remainder theorem to find the remainder when the first polynomial is divided by the second polynomial: $x^4+x^3+x^2+x+1$, $x-1$

\item Use the remainder theorem to find the remainder when the first polynomial is divided by the second polynomial: $x^4+x^2+1$, $x+3$

\item Use the remainder theorem to find the remainder when the first polynomial is divided by the second polynomial: $x^4+x^3+2$, $x+2$

\item Use the factor theorem to determine whether the first polynomial is a factor of the second polynomial: $x+1$, $x^2-1$

\item Use the factor theorem to determine whether the first polynomial is a factor of the second polynomial: $x-2$, $x^3+x^2-7x+2$

\item Use the factor theorem to determine whether the first polynomial is a factor of the second polynomial: $x-3$, $x^4-3x^3+x^2-x+1$

\item Use synthetic division to show that $x=2$ is a zero of the polynomial $x^3-7x+6$, and use the result to factorize the polynomial completely.

\item Use synthetic division to show that $x=-4$ is a zero of the polynomial $x^3-28x-48$, and use the result to factorize the polynomial completely.

\item Use synthetic division to show that $x=2$ and $x=-3$ are zeros of the polynomial $2x^4+7x^3-4x^2-27x-18$, and use the result to factorize the polynomial completely.

\item Use synthetic division to find the quotient and the remainder when the polynomial $x^4-10x^2-2x+4$ is divided by $x+3$.

\item If $x+1$ and $x-2$ are factors of $x^3-px^2+qx+2$. Using synthetic division, find the values of $p$ and $q$.

\item When the polynomial $4x^4+2x^3+kx^2+13$ is divided by $x+1$, the remainder is 16. Find the value of $k$.

\item When the polynomial $x^3+x^2+x+k$ is divided by $x+1$, the remainder is 7. Find the value of $k$.

\item Use the factor theorem to find the values of $p$ and $q$ if $x+1$ and $x-2$ are factors of the polynomial $x^3+px^2+qx+3$.

\item Use the factor theorem to find the values of $a$ and $b$ if $-2$ and $2$ are the roots of the polynomial $2x^3+4x^2+ax+b$.

\item Divide the cubic polynomial $3x^3-10x^2+13x-6$ by the linear polynomial $x-2$. Also find the quotient and remainder.

\item Divide the polynomial $x^4-3x^3+5x^2-7x+2$ by $x^2-x+1$. Also find the quotient and remainder.

\item Find the value of $k$ if the polynomial $x^3+kx^2-7x+6$ has a remainder $-4$ when divided by $x+2$.

\item If $(x-2)$ and $(x+2)$ are factors of $x^4-13x^2+36$. Using synthetic division, find the other two factors.

\item Consider a data set of monthly sales figures. A polynomial regression model $P(x)=x^3+2x^2+x-3$ is fitted to this data. If the observed sales in the 5th month are 240 units, find the percentage error.

\item A retailer company has developed a polynomial regression model to predict weekly product demand: $D(w)=w^3-2w^2+5w-4$, where $D(w)$ represents predicted demand (in units) and $w$ is the week number. Given the actual demand is 22 units, calculate the prediction error.

\item A digital processing system has a transfer function with numerator $B(z)=z^2-z-2$. Use the factor theorem to find the zeros of the system.

\item A signal processing system has a transfer function $H(z)=\dfrac{z^2+3z+2}{z^2-0.2z+0.9}$. Find the zero(s) of the transfer function using the factor theorem.

\item A signal processing system has a transfer function $H(z)=\dfrac{z^2-0.5z-0.5}{z^3+1}$. Find the zero(s) of the transfer function using the factor theorem.

\item A signal processing system has a transfer function with denominator $A(z)=z^2-0.3z-0.4$. Use the factor theorem to find the poles of the system and determine if the system is stable.

\item The denominator of a signal processing system's transfer function equals $A(z)=z^2+1.2z+0.35$. Use the factor theorem to determine the location of the corresponding poles and assess the stability of the system.

\item Prove the following: $\sin210^\circ+\cos240^\circ+\tan225^\circ+\cot225^\circ=1$

\item Prove that: $\dfrac{\tan(180^\circ+\alpha)\cot(90^\circ-\alpha)}{\sin(360^\circ-\alpha)\cos(270^\circ+\alpha)}=-\sec^2\alpha$

\item Prove that: $\dfrac{\sin^2(\pi+\theta)\tan\left(\dfrac{3\pi}{2}+\theta\right)}{\cot^2\left(\dfrac{3\pi}{2}-\theta\right)\cos^2(\pi-\theta)\csc(2\pi-\theta)}=\cos\theta$

\item Prove that: $\dfrac{\cos(90^\circ+\theta)\sec(-\theta)\tan(180^\circ-\theta)}{\sec(360^\circ-\theta)\sin(180^\circ+\theta)\cot(90^\circ-\theta)}=-1$

\item Show that: $\sec\left(\dfrac{3\pi}{2}-\theta\right)\sec\left(\dfrac{5\pi}{2}-\theta\right)-\tan\left(\dfrac{3\pi}{2}-\theta\right)\tan\left(\dfrac{5\pi}{2}+\theta\right)=-1$

\item If $\alpha,\beta,\gamma$ are the angles of a triangle $ABC$, then prove that: $\sin(\alpha+\beta)=\sin\gamma$

\item If $\alpha,\beta,\gamma$ are the angles of a triangle $ABC$, then prove that: $\sec\left(\dfrac{\alpha+\beta}{2}\right)=\csc\dfrac{\gamma}{2}$

\item If $\alpha,\beta,\gamma$ are the angles of a triangle $ABC$, then prove that: $\csc\alpha=\dfrac{1}{\sin(\beta+\gamma)}$

\item If $\alpha,\beta,\gamma$ are the angles of a triangle $ABC$, then prove that: $\tan(\alpha+\beta)+\tan\gamma=0$

\item Simplify: $\dfrac{\sin(180^\circ-\theta)\cos(360^\circ-\theta)\tan(90^\circ+\theta)}{\sin(90^\circ-\theta)\cos(180^\circ+\theta)\tan(270^\circ-\theta)}$

\item Prove that: $\dfrac{1-\tan\theta\tan\phi}{1+\tan\theta\tan\phi}=\dfrac{\cos(\theta+\phi)}{\cos(\theta-\phi)}$

\item Show that $\cos(\alpha+\beta)\cos(\alpha-\beta)=\cos^2\alpha-\sin^2\beta=\cos^2\beta-\sin^2\alpha$

\item Show that $\dfrac{\sin(\alpha+\beta)+\sin(\alpha-\beta)}{\cos(\alpha+\beta)+\cos(\alpha-\beta)}=\tan\alpha$

\item Show that $\cot(\alpha+\beta)=\dfrac{\cot\alpha\cot\beta-1}{\cot\alpha+\cot\beta}$

\item Show that $\cot(\alpha-\beta)=\dfrac{\cot\alpha\cot\beta+1}{\cot\beta-\cot\alpha}$

\item Show that $\dfrac{\tan\alpha+\tan\beta}{\tan\alpha-\tan\beta}=\dfrac{\sin(\alpha+\beta)}{\sin(\alpha-\beta)}$

\item Show that: $\cos(\alpha-\beta)=\dfrac{1+\tan\alpha\tan\beta}{\sec\alpha\sec\beta}$

\item Show that: $\sin(\alpha+\beta)=\dfrac{1+\cot\alpha\tan\beta}{\csc\alpha\sec\beta}$

\item Show that: $\cot(\alpha-\beta)=\dfrac{\cot\alpha\cot\beta+1}{\cot\beta-\cot\alpha}$

\item Show that: $\dfrac{\tan\alpha+\tan\beta}{\tan\alpha-\tan\beta}=\dfrac{\sin(\alpha+\beta)}{\sin(\alpha-\beta)}$

\item Show that: $\cot(\alpha+\beta)=\dfrac{\cot\alpha\cot\beta-1}{\cot\alpha+\cot\beta}$

\item If $\sin\alpha=\dfrac{24}{25}$ and $\cos\beta=\dfrac{20}{29}$, where $0<\alpha<\dfrac{\pi}{2}$ and $0<\beta<\dfrac{\pi}{2}$, show that $\sin(\alpha-\beta)=\dfrac{333}{725}$.

\item If $\sin\alpha=-\dfrac{8}{17}$ and $\cos\beta=-\dfrac{4}{5}$, where $\dfrac{3\pi}{2}<\alpha<2\pi$ and $\pi<\beta<\dfrac{3\pi}{2}$, find $\sin(\alpha+\beta)$; in which quadrants do the terminal sides of the angles of measure $(\alpha+\beta)$ and $(\alpha-\beta)$ lie?

\item Using the same data as above, find $\cos(\alpha+\beta)$; in which quadrants do the terminal sides of $(\alpha+\beta)$ and $(\alpha-\beta)$ lie?

\item Using the same data as above, find $\tan(\alpha+\beta)$; in which quadrants do the terminal sides of $(\alpha+\beta)$ and $(\alpha-\beta)$ lie?

\item Using the same data as above, find $\sin(\alpha-\beta)$; in which quadrants do the terminal sides of $(\alpha+\beta)$ and $(\alpha-\beta)$ lie?

\item Using the same data as above, find $\cos(\alpha-\beta)$; in which quadrants do the terminal sides of $(\alpha+\beta)$ and $(\alpha-\beta)$ lie?

\item Using the same data as above, find $\tan(\alpha-\beta)$; in which quadrants do the terminal sides of $(\alpha+\beta)$ and $(\alpha-\beta)$ lie?

\item Find $\sin(\alpha+\beta)$ and $\cos(\alpha+\beta)$, given that $\tan\alpha=\dfrac{3}{4}$, $\cos\beta=\dfrac{5}{13}$ and neither the terminal side of the angle of measure $\alpha$ nor that of $\beta$ is in quadrant I.

\item Find $\sin(\alpha+\beta)$ and $\cos(\alpha+\beta)$, given that $\tan\alpha=-\dfrac{15}{8}$ and $\sin\beta=-\dfrac{7}{25}$ and neither the terminal side of the angle of measure $\alpha$ nor that of $\beta$ is in quadrant IV.

\item Prove that: $\dfrac{\cos19^\circ+\sin19^\circ}{\cos19^\circ-\sin19^\circ}=\tan64^\circ$

\item Prove that: $\cos(60^\circ+\theta)\cos(60^\circ-\theta)+\sin(60^\circ+\theta)\sin(60^\circ-\theta)=\cos2\theta$

\item If $\alpha,\beta,\gamma$ are the angles of a triangle $ABC$, show that: $\cot\dfrac{\alpha}{2}+\cot\dfrac{\beta}{2}+\cot\dfrac{\gamma}{2}=\cot\dfrac{\alpha}{2}\cot\dfrac{\beta}{2}\cot\dfrac{\gamma}{2}$

\item If $\alpha+\beta+\gamma=180^\circ$, show that $\cot\alpha\cot\beta+\cot\beta\cot\gamma+\cot\gamma\cot\alpha=1$.

\item Express the following in the form $r\sin(\theta+\phi)$ or $r\sin(\theta-\phi)$, where the terminal sides of the angles of measure $\theta$ and $\phi$ are in the first quadrant: $24\sin\theta+7\cos\theta$

\item Express the following in the form $r\sin(\theta+\phi)$ or $r\sin(\theta-\phi)$: $12\sin\theta-5\cos\theta$

\item Express the following in the form $r\sin(\theta+\phi)$ or $r\sin(\theta-\phi)$: $\sin\theta-\cos\theta$

\item Express the following in the form $r\sin(\theta+\phi)$ or $r\sin(\theta-\phi)$: $8\sin\theta-6\cos\theta$

\item Express the following in the form $r\sin(\theta+\phi)$ or $r\sin(\theta-\phi)$: $\dfrac{1}{2}\sin\theta+\dfrac{\sqrt{3}}{2}\cos\theta$

\item Express the following in the form $r\sin(\theta+\phi)$ or $r\sin(\theta-\phi)$: $13\sin\theta-84\cos\theta$

\item Prove the following: $\dfrac{\cos\alpha-\sin\alpha}{\cos\alpha+\sin\alpha}=\sec2\alpha-\tan2\alpha$

\item If $\cos\alpha=-\dfrac{7}{25}$, $\tan\beta=\dfrac{12}{5}$, the terminal side of the angle of measure $\alpha$ is in quadrant II and that of $\beta$ is in quadrant III, find $\sin(\alpha+\beta)$; in which quadrant does the terminal side of $(\alpha+\beta)$ lie?

\item Using the same data, find $\cos(\alpha+\beta)$; in which quadrant does the terminal side of $(\alpha+\beta)$ lie?

\item If $\alpha,\beta,\gamma$ are the angles of $\triangle ABC$, prove that: $\tan\alpha+\tan\beta+\tan\gamma=\tan\alpha\tan\beta\tan\gamma$

\item If $\alpha,\beta,\gamma$ are the angles of $\triangle ABC$, prove that: $\tan\dfrac{\alpha}{2}\tan\dfrac{\beta}{2}+\tan\dfrac{\beta}{2}\tan\dfrac{\gamma}{2}+\tan\dfrac{\gamma}{2}\tan\dfrac{\alpha}{2}=1$

\item Express $3\sin\theta+4\cos\theta$ in the form $r\sin(\theta+\phi)$, where the terminal side of the angle of measure $\phi$ is in quadrant I.

\item Prove that: $\dfrac{\sin\theta-\cos\theta\tan\dfrac{\theta}{2}}{\cos\theta+\sin\theta\tan\dfrac{\theta}{2}}=\tan\dfrac{\theta}{2}$

\item Prove the following: $\sqrt{\dfrac{1+\sin\alpha}{1-\sin\alpha}}=\dfrac{\sin\dfrac{\alpha}{2}+\cos\dfrac{\alpha}{2}}{\sin\dfrac{\alpha}{2}-\cos\dfrac{\alpha}{2}}$

\item Prove the following: $\dfrac{\csc\theta+2\csc2\theta}{\sec\theta}=\cot\dfrac{\theta}{2}$

\item Prove the following: $\dfrac{\sin3\theta}{\sin\theta}-\dfrac{\cos3\theta}{\cos\theta}=2$

\item Prove the following: $\dfrac{\cos3\theta}{\cos\theta}+\dfrac{\sin3\theta}{\sin\theta}=4\cos2\theta$

\item Prove the following: $\dfrac{\tan\dfrac{\theta}{2}+\cot\dfrac{\theta}{2}}{\cot\dfrac{\theta}{2}-\tan\dfrac{\theta}{2}}=\sec\theta$

\item Prove the following: $\dfrac{\sin3\theta}{\cos\theta}+\dfrac{\cos3\theta}{\sin\theta}=2\cot2\theta$

\item Prove the following: $\dfrac{3+\cos4\theta}{1-\cos4\theta}=\dfrac{1}{2}\left(\tan^2\theta+\cot^2\theta\right)$

\item Prove the following: $\cos^2\dfrac{\pi}{8}+\cos^2\dfrac{3\pi}{8}+\cos^2\dfrac{5\pi}{8}+\cos^2\dfrac{7\pi}{8}=2$

\item Reduce $\sin^4\theta$ to an expression involving only functions of multiples of $\theta$, raised to the first power.

\item Find the values of $\sin\theta$ and $\cos\theta$ without using a table or calculator, when $\theta=18^\circ$.

\item Find the values of $\sin\theta$ and $\cos\theta$ without using a table or calculator, when $\theta=36^\circ$.

\item Find the values of $\sin\theta$ and $\cos\theta$ without using a table or calculator, when $\theta=54^\circ$.

\item Prove the following identity: $\dfrac{\sin3x-\sin x}{\cos x-\cos3x}=\cot2x$

\item Prove the following identity: $\dfrac{\sin8x-\sin2x}{\cos8x+\cos2x}=\tan5x$

\item Prove the following identity: $\dfrac{\sin A-\sin B}{\sin A+\sin B}=\tan\dfrac{A-B}{2}\cot\dfrac{A+B}{2}$

\item Prove the following identity: $\dfrac{\sin80^\circ+\sin40^\circ}{\cos80^\circ+\cos40^\circ}=\sqrt{3}$

\item Prove that: $\cos15^\circ+\cos105^\circ+\cos195^\circ+\cos165^\circ+\cos285^\circ=0$

\item Prove that: $\dfrac{\sin2\theta+\sin4\theta+\sin6\theta+\sin8\theta}{\cos2\theta+\cos4\theta+\cos6\theta+\cos8\theta}=\tan5\theta$

\item Prove that: $\cos^2\left(\dfrac{\pi}{4}-\dfrac{\alpha}{2}\right)-\cos^2\left(\dfrac{\pi}{4}+\dfrac{\alpha}{2}\right)=\sin\alpha$

\item Prove that: $\sin\left(\dfrac{\pi}{4}-\theta\right)\sin\left(\dfrac{\pi}{4}+\theta\right)=\dfrac{1}{2}\cos2\theta$

\item Prove that: $\dfrac{\sin\theta+\sin3\theta+\sin5\theta+\sin7\theta}{\cos\theta+\cos3\theta+\cos5\theta+\cos7\theta}=\tan4\theta$

\item Prove that: $\cos20^\circ\cos40^\circ\cos60^\circ\cos80^\circ=\dfrac{1}{16}$

\item Prove that: $\sin\dfrac{\pi}{9}\sin\dfrac{2\pi}{9}\sin\dfrac{\pi}{3}\sin\dfrac{4\pi}{9}=\dfrac{3}{16}$

\item Prove that: $\sin10^\circ\sin30^\circ\sin50^\circ\sin70^\circ=\dfrac{1}{16}$

\item Prove that: $\dfrac{\sin3\theta}{1+2\cos2\theta}=\sin\theta$; deduce the value of $\sin15^\circ$.

\item Prove that: $\dfrac{\sin^2\alpha-\sin^2\beta}{\sin\alpha\cos\alpha-\sin\beta\cos\beta}=\tan(\alpha+\beta)$

\item Prove that: $\sin\alpha\sin\beta+\sin\gamma-\sin(\alpha+\beta+\gamma)=4\sin\left(\dfrac{\alpha+\beta}{2}\right)\sin\left(\dfrac{\beta+\gamma}{2}\right)\sin\left(\dfrac{\gamma+\alpha}{2}\right)$

\item Prove without using a table or calculator, that $\sin19^\circ\cos11^\circ+\sin71^\circ\sin11^\circ=\dfrac{1}{2}$

\item Show that $\cos20^\circ\cos40^\circ\cos80^\circ=\dfrac{1}{8}$.

\item Determine whether the following function is even, odd, or neither: $\sin^4x+\cos^4x$

\item Determine whether the following function is even, odd, or neither: $\dfrac{\sin x+\sin3x}{\cos x+\cos3x}$

\item Determine whether the following function is even, odd, or neither: $\dfrac{1}{\sec x+\sec^3x}$

\item Determine whether the following function is even, odd, or neither: $\dfrac{1}{\sec x+\cot^2x}$

\item Draw the graph of the following function for the interval mentioned: $y=-\sin2x$, $x\in[-2\pi,2\pi]$

\item Draw the graph of the following function for the interval mentioned: $y=2\cos2x$, $x\in[-2\pi,2\pi]$

\item Draw the graph of the following function for the interval mentioned: $y=\tan2x$, $x\in[-\pi,\pi]$

\item Draw the graph of the following function for the interval mentioned: $y=\tan\dfrac{x}{2}$, $x\in[-2\pi,2\pi]$

\item Draw the graph of the following function for the interval mentioned: $y=\sin\dfrac{\pi}{2}x$, $x\in[0,2\pi]$

\item Draw the graph of the following function for the interval mentioned: $y=\cos\dfrac{\pi}{2}x$, $x\in[-\pi,\pi]$

\item On the same axes and to the same scale, draw graphs of $y=\sin x$ and $y=\sin2x$ for their complete period.

\item On the same axes and to the same scale, draw graphs of $y=\cos x$ and $y=\cos2x$ for their complete period.

\item The temperature $T$ in degrees Celsius of a certain city varies throughout the day according to the equation $T(t)=\dfrac{13}{2}\sin\left(\dfrac{\pi}{6}t-\dfrac{\pi}{9}\right)+15$, where $t$ is the time in hours, with $t=0$ corresponding to midnight. Find the maximum and minimum temperature during the day.

\item Using the same equation as above, find the temperature at $t=9$ hours (9:00 a.m.).

\item A man on top of a 100 m high lighthouse is in line with two ships on the same side of it, whose angles of depression from the man are $17^\circ$ and $19^\circ$ respectively. Find the distance between the ships.

\item P and Q are two points in line with a tree. If the distance between P and Q be 30 m and the angles of elevation of the top of the tree at P and Q are $12^\circ$ and $15^\circ$ respectively, find the height of the tree.

\item A giant Ferris wheel has a diameter of 60 feet. The lowest point of the wheel is located 6 feet above the ground. The wheel completes one full revolution every 80 seconds. Make an equation that represents the height $h(t)$ of a rider on the Ferris wheel at any given time $t$.

\item A giant Ferris wheel has a diameter of 60 feet. The lowest point of the wheel is located 6 feet above the ground. The wheel completes one full revolution every 80 seconds. Find the maximum height of the rider.

\item A giant Ferris wheel has a diameter of 60 feet. The lowest point of the wheel is located 6 feet above the ground. The wheel completes one full revolution every 80 seconds. Find the height of the rider from the ground after 35 seconds.

\item A child is playing on the swing in a playground. The height $h(t)$ of the swing seat above the ground (in metres) at time $t$ (in seconds) is modeled by the function $h(t)=1.5+1.2\sin(3\pi t)$. What is the maximum height reached by the swing seat?

\item A child is playing on the swing in a playground. The height $h(t)$ of the swing seat above the ground (in metres) at time $t$ (in seconds) is modeled by the function $h(t)=1.5+1.2\sin(3\pi t)$. What is the minimum height reached by the swing seat?

\item A child is playing on the swing in a playground. The height $h(t)$ of the swing seat above the ground (in metres) at time $t$ (in seconds) is modeled by the function $h(t)=1.5+1.2\sin(3\pi t)$. How long does it take for the swing to complete one full back-and-forth motion (period)?

\item A child is playing on the swing in a playground. The height $h(t)$ of the swing seat above the ground (in metres) at time $t$ (in seconds) is modeled by the function $h(t)=1.5+1.2\sin(3\pi t)$. At what time(s) does the swing seat first reach a height of 2.12 metres?

\item A carnival ride consists of a vertical wheel with a diameter of 40 feet. The centre of the wheel is 28 feet above the ground. The wheel rotates at a constant speed and takes 120 seconds to make one complete revolution. Model an equation that describes the height $h(t)$ of a rider on the wheel as a function of time $t$. How high is the rider from the ground after 90 seconds? At what times will the rider be 36 feet above the ground?

\item Suppose the temperature $T$ in degrees Fahrenheit of Lahore city in the month of December throughout the day can be modeled by the equation $T = 64+8\sin\left(\dfrac{\pi}{12}(t-8)\right)$, where $t$ is the time in hours. The temperature oscillates 8 degrees above and below an average temperature of 64 degrees. Find the temperature at $t=9$ hours.

\item Suppose the temperature $T$ in degrees Fahrenheit of Lahore city in the month of December throughout the day can be modeled by the equation $T = 64+8\sin\left(\dfrac{\pi}{12}(t-8)\right)$, where $t$ is the time in hours. The temperature oscillates 8 degrees above and below an average temperature of 64 degrees. At what time will the temperature be maximum?

\item Suppose the temperature $T$ in degrees Fahrenheit of Lahore city in the month of December throughout the day can be modeled by the equation $T = 64+8\sin\left(\dfrac{\pi}{12}(t-8)\right)$, where $t$ is the time in hours. The temperature oscillates 8 degrees above and below an average temperature of 64 degrees. Calculate the maximum temperature.

\item Suppose the population of a coastal city follows a sinusoidal pattern due to seasonal migration. The population of the city over the course of a year can be modeled by the equation $P(t) = 70000 +10000\cos\left(\dfrac{\pi}{6}t-\dfrac{\pi}{2}\right)$, where $P(t)$ is the population at time $t$ ($t$ is the time in months, with $t=0$ corresponding to January 1st) and $t$ denotes the months in a year. Find the population of the city at $t=7$ months.

\item Suppose the population of a coastal city follows a sinusoidal pattern due to seasonal migration. The population of the city over the course of a year can be modeled by the equation $P(t) = 70000 +10000\cos\left(\dfrac{\pi}{6}t-\dfrac{\pi}{2}\right)$, where $P(t)$ is the population at time $t$ ($t$ is the time in months, with $t=0$ corresponding to January 1st) and $t$ denotes the months in a year. Find the maximum population.

\item A Ferris wheel with a radius of 45 feet has its lower point located 5 feet above the ground. It completes one full revolution every 60 seconds in a counter-clockwise direction. Model an equation that describes the height of a rider on the Ferris wheel as a function of time $t$. How high is the rider from the ground after 40 seconds? Also graph the model equation.

\item The water level $L$ (in feet) of a tidal river varies throughout the day. Suppose the level of the tidal river can be modeled by the equation $L(t) = 8+4\sin\left(\dfrac{\pi}{6}t\right)$, where $t$ denotes the time (in hours). The water level oscillates 4 feet above and below an average level of 8 feet. Find the water level at $t=3$ hours.

\item The water level $L$ (in feet) of a tidal river varies throughout the day. Suppose the level of the tidal river can be modeled by the equation $L(t) = 8+4\sin\left(\dfrac{\pi}{6}t\right)$, where $t$ denotes the time (in hours). The water level oscillates 4 feet above and below an average level of 8 feet. What is the maximum water level?

\item From a point 100 m above the surface of a lake, the angle of elevation of a peak of a cliff is found to be $15^\circ$ and the angle of depression of the image of the peak is $30^\circ$. Find the height of the peak.

\item Evaluate the following limit: $\displaystyle\lim_{x\to0}\dfrac{\cos ax - \cos bx}{\cos cx - \cos dx}$

\item Evaluate the following limit: $\displaystyle\lim_{x\to3}\dfrac{x^2-x\log x + 3\log x -9}{x-3}$

\item Discuss the continuity of $f(x)$ at $x=c$: $f(x) = \begin{cases} 2x+5 & \text{if } x\leq 2 \\ 4x+1 & \text{if } x>2 \end{cases}$, $c=2$

\item Discuss the continuity of $f(x)$ at $x=c$: $f(x) = \begin{cases} 3x-1 & \text{if } x<1 \\ 4 & \text{if } x=1 \\ 2x & \text{if } x>1 \end{cases}$, $c=1$

\item If $f(x) = \begin{cases} 3x & \text{if } x\leq -2 \\ x^2-1 & \text{if } -2<x<2 \\ 3 & \text{if } x\geq 2 \end{cases}$, discuss continuity at $x=2$ and $x=-2$.

\item If $f(x) = \begin{cases} x+2 & x\leq -1 \\ c+2 & x>-1 \end{cases}$, find $c$ so that $\displaystyle\lim_{x\to-1} f(x)$ exists.

\item Find the values of $m$ and $n$, so that the given function $f$ is continuous at $x=3$: $f(x) = \begin{cases} mx & \text{if } x<3 \\ n & \text{if } x=3 \\ -2x+9 & \text{if } x>3 \end{cases}$

\item Find the values of $m$ and $n$, so that the given function $f$ is continuous at $x=3$: $f(x) = \begin{cases} mx & \text{if } x<3 \\ x^2 & \text{if } x\geq 3 \end{cases}$

\item $f(x)=\begin{cases} \dfrac{\sqrt{2x+5}-\sqrt{x+7}}{x-2} & x\neq 2 \\ k & x=2 \end{cases}$. Find the value of $k$ so that $f$ is continuous at $x=2$.

\item Given the function $f(x) = \begin{cases} 2x+3 & x\leq 1 \\ -x+4 & x>1 \end{cases}$. Discuss the limit and continuity at $x=1$.

\item Determine whether $\displaystyle\lim_{x\to2} f(x)$ and $\displaystyle\lim_{x\to4} f(x)$ exist, when $f(x) = \begin{cases} 2x+1 & \text{if } 0\leq x\leq 2 \\ 7-x & \text{if } 2<x<4 \\ x & \text{if } 4\leq x\leq 6 \end{cases}$

\item Consider the function $f(x) = \dfrac{x^2-1}{x-1}$, discuss the continuity of $f$ at $x=1$.

\item For $f(x)=3x^2-5x+4$, discuss continuity of $f$ at $x=1$.

\item Discuss the continuity of the function $f(x)$ and $g(x)$ at $x=3$: $f(x) = \begin{cases} \dfrac{x^2-9}{x-3} & \text{if } x\neq 3 \\ 6 & \text{if } x=3 \end{cases}$

\item Discuss the continuity of the function $f(x)$ and $g(x)$ at $x=3$: $g(x) = \dfrac{x^2-9}{x-3}$ if $x\neq 3$.

\item Discuss continuity of $f(x)$ at $x=3$, when $f(x) = \begin{cases} x-1 & x<3 \\ 2x+1 & x\geq 3 \end{cases}$

\item A substance decays exponentially following the formula $A(t)=A_0e^{-0.1t}$, where $A_0$ is the initial amount. Find the limit of $A(t)$ as $t\to\infty$.

\item A town's population is modeled by $P(t) = \dfrac{100{,}000}{1+9e^{-0.5t}}$. What is the long-term population as $t\to\infty$?

\item A company's weekly sales (in thousands) follow the function $S(t) = \dfrac{500t}{t+10}$. What is the limit of $S(t)$ as $t\to\infty$ and what does it represent?

\item Signal strength $S(d)$ at a distance $d$ from a tower is modeled as $S(d) = \dfrac{1000}{d^2}$. What is the signal strength at $d=10$?

\item Signal strength $S(d)$ at a distance $d$ from a tower is modeled as $S(d) = \dfrac{1000}{d^2}$. What happens to signal strength as $d\to\infty$?

\item The factory's cost function is given as $C(x) = \begin{cases} 10x+500 & \text{if } x\leq100 \\ 12x+300 & \text{if } x>100 \end{cases}$. Is the cost function continuous at $x=100$?

\item Inflation is modeled by $I(t)=I_0e^{0.03t}$, where $I_0$ is the initial price index and $t$ is the number of years. Find the inflation rate after 5 years if $I_0=100$.

\item Inflation is modeled by $I(t)=I_0e^{0.03t}$, where $I_0$ is the initial price index and $t$ is the number of years. What is the expected price index after 10 years?

\item The cost of producing $x$ units is $C(x) = \begin{cases} 5x+20 & \text{if } x\leq 10 \\ 6x+10 & \text{if } x>10 \end{cases}$. Is the cost function continuous at $x=10$?

\item The radioactive decay of a substance is given by the function $A(t)=A_0e^{-kt}$ where $A_0$ is the initial amount of the substance, $k$ is the decay constant, and $t$ is the time in years. Find the limit of the amount of substance as $t\to\infty$.

\item The value of an investment grows according to the formula for continuous compounding $A(t) = P_0e^{rt}$, where $P_0$ is the initial principal, $r$ is the annual interest rate, and $t$ is the time in years. What happens to the value of the investment as $t\to\infty$?

\item In economics, the demand function $D(p)$ decreases as the price $p$ increases. Suppose the demand function is given by $D(p) = \dfrac{100}{p+1}$, where $p$ is the price in dollars. Find the limit of the demand as the price becomes very large, i.e., $\displaystyle\lim_{p\to\infty} D(p)$.

\item The apparent brightness $B(d)$ of a star decreases as the distance from Earth increases following the inverse square law $B(d) = \dfrac{L}{d^2}$, where $L$ is the star's luminosity. Find the limit of the brightness as $d\to\infty$.

\item Find by definition, the derivative w.r.t $x$ of the following function defined as: $2x^2+1$

\item Find by definition, the derivative w.r.t $x$ of the following function defined as: $2-\sqrt{x}$

\item Find (by equation) from the first principle the final gradient of the curve at the given point: $\sqrt{x+2}$ at $x=6$.

\item Find (by equation) from the first principle the final gradient of the curve at the given point: $\dfrac{1}{\sqrt{x+a}}$ at $x=a$.

\item Find from first principle the derivatives of the following expressions w.r.t.\ their respective independent variables: $(3x-2)^{-2}$

\item Find from first principle the derivatives of the following expressions w.r.t.\ their respective independent variables: $(2t+3)^5$

\item Find from first principle the derivatives of the following expressions w.r.t.\ their respective independent variables: $(aw+b)^7$

\item The outdoor temperature (in $^\circ$C) over time is modeled by $T(t)=-t^2+12t+10$, where $t$ is the time in hours. Find the instantaneous rate of change at $t=2$.

\item Differentiate w.r.t.\ $x$: $\dfrac{2x-1}{\sqrt{x^2+1}}$

\item Differentiate w.r.t.\ $x$: $\sqrt{\dfrac{a-x}{a+x}}$

\item Differentiate w.r.t.\ $x$: $\dfrac{\sqrt{x^2+1}}{\sqrt{x^2-1}}$

\item Find $\dfrac{dy}{dx}$ if $y = \dfrac{(\sqrt{x}+1)\left(x^{\frac{3}{2}}-1\right)}{x^{\frac{1}{2}}}$, $(x\neq 1)$

\item Differentiate $\dfrac{(\sqrt{x}+1)\left(x^{\frac{3}{2}}-1\right)}{x^{\frac{3}{2}}-x^{\frac{1}{2}}}$ with respect to $x$.

\item If $y = \sqrt{x}-\dfrac{1}{\sqrt{x}}$, show that $2x\dfrac{dy}{dx}+y = 2\sqrt{x}$.

\item If $y = x^4+2x^2+2$, prove that $\dfrac{dy}{dx}=4x\sqrt{y-1}$.

\item A particle moves along a line such that its position after $t$ hours is given by $s(t)=4t^2+2t+1$ (in miles). Find the average velocity over the interval $[2,5]$.

\item A particle moves along a line such that its position after $t$ hours is given by $s(t)=4t^2+2t+1$ (in miles). Find the instantaneous velocity at $t=3$.

\item If $y = \dfrac{1}{x^2}$, then find $\dfrac{dy}{dx}$ at $x=-1$ by the ab-initio method.

\item Find the derivative of $y = \dfrac{3}{4}x^4+\dfrac{2}{3}x^3+\dfrac{1}{2}x^2+2x+5$ w.r.t.\ $x$.

\item Structural stress on a bridge is modeled by the function $S(x)=100-5x^2$, where $x$ is the distance from the centre of the bridge. Calculate the rate of change of stress at that point.

\item A company's revenue function is given by $R(x)=100x-10x^2$, where $x$ is the number of units produced. The cost function is $C(x)=300x+2000$. Find the profit function $P(x)$.

\item A company's revenue function is given by $R(x)=100x-10x^2$, where $x$ is the number of units produced. The cost function is $C(x)=300x+2000$. Determine the marginal profit when $x=15$.

\item An investment grows according to the function $A(t)=1000(1+0.5t)^3$, where $A(t)$ is the value of the investment and $t$ is the time in years. Find the rate of change of the investment after 8 years.

\item An investment grows according to the function $A(t)=1000(1+0.5t)^3$, where $A(t)$ is the value of the investment and $t$ is the time in years. What is the investment value after 8 years?

\item The position of a particle moving along a line is given by $s(t)=5t^3-12t^2+8t$, where $s(t)$ is the position in metres and $t$ is the time in seconds. Determine the velocity of the particle at $t=4$ seconds.

\item The position of a particle moving along a line is given by $s(t)=5t^3-12t^2+8t$, where $s(t)$ is the position in metres and $t$ is the time in seconds. Find the acceleration at $t=4$ seconds.

\item The position of a particle moving along a line is given by $s(t)=5t^3-12t^2+8t$, where $s(t)$ is the position in metres and $t$ is the time in seconds. When is the particle at rest?

\item The position of a car traveling along a straight highway is given by $x(t)=30t^2-4t^3$, where $x(t)$ is the distance traveled in kilometres and $t$ is the time in hours. Find the car's velocity at $t=3$ hours.

\item The position of a car traveling along a straight highway is given by $x(t)=30t^2-4t^3$, where $x(t)$ is the distance traveled in kilometres and $t$ is the time in hours. Determine the car's acceleration at $t=3$ hours.

\item The stress on a beam under a varying load is given by $S(x)=400x-x^2$, where $S(x)$ is the stress in pascals (Pa) and $x$ is the distance from the fixed end in metres. Calculate the rate of change of stress at 6 metres.

\item The cost $C(r)$ to construct a cylindrical tank depends on the radius of the base, and is given by $C(r) =8000\pi r^2+\dfrac{150000}{r}$, where the first term represents the cost of the base and the second term represents the cost of the walls. Determine the rate of change of the cost at $r=4$ metres.

\item A company's profit function is given by $P(x)=100x-5x^2$, where $x$ is the number of units produced. Determine the marginal profit when $x=8$ units.

\item A particle moves along a line according to the position function $s(t)=4t^3-3t^2+2t$, where $s(t)$ is the position in metres and $t$ is the time in seconds. Find the velocity and acceleration at $t=2$ seconds.

\item A bank offers a compound interest rate on an investment, and the value of the investment after $t$ years is given by $V(t)=5000(1+0.04t)^2$. Find the rate of change of the investment value after 10 years.

\item The stress on a beam under a varying load is modeled by $S(x)=500x-2x^3$, where $S(x)$ is the stress in pascals (Pa) and $x$ is the distance (in metres) from the beam's fixed end. Find the rate of change of stress at $x=5$ metres.

\item If $\underline{u}= x\underline{i} + 2\underline{j} + 3\underline{k}$, $\underline{v}=\underline{i}+y\underline{j}-3\underline{k}$ and $\underline{w}=2\underline{i}-3\underline{j}$ represent the sides of a triangle, find the values of $x$ and $y$.

\item The position vectors of the points $A,B,C$ and $D$ are $\underline{u}=\underline{i}+2\underline{j}+\underline{k}$, $\underline{v}=7\underline{i}+8\underline{j}+4\underline{k}$, $\underline{w}=-\underline{i}+\underline{k}$ and $\underline{z}=\underline{i}+2\underline{j}+2\underline{k}$ respectively. Show that $\overrightarrow{AB}$ is parallel to $\overrightarrow{CD}$.

\item We say that two vectors $\underline{v}$ and $\underline{w}$ in space are parallel if there is a scalar $c$ such that $\underline{v}=c\underline{w}$. The vectors point in the same direction if $c>0$ and the vectors point in the opposite direction if $c<0$. Find two vectors of length 2 parallel to the vector $\underline{v}=2\underline{i}-4\underline{j}+4\underline{k}$.

\item We say that two vectors $\underline{v}$ and $\underline{w}$ in space are parallel if there is a scalar $c$ such that $\underline{v}=c\underline{w}$. The vectors point in the same direction if $c>0$ and the vectors point in the opposite direction if $c<0$. Find the constant $a$ so that the vectors $\underline{v}=\underline{i}-3\underline{j}+4\underline{k}$ and $\underline{w}=a\underline{i}+9\underline{j}-12\underline{k}$ are parallel.

\item We say that two vectors $\underline{v}$ and $\underline{w}$ in space are parallel if there is a scalar $c$ such that $\underline{v}=c\underline{w}$. The vectors point in the same direction if $c>0$ and the vectors point in the opposite direction if $c<0$. Find a vector of length 5 in the direction opposite that of $\underline{v}=\underline{i}-2\underline{j}+3\underline{k}$.

\item We say that two vectors $\underline{v}$ and $\underline{w}$ in space are parallel if there is a scalar $c$ such that $\underline{v}=c\underline{w}$. The vectors point in the same direction if $c>0$ and the vectors point in the opposite direction if $c<0$. Find $a$ and $b$ so that the vectors $3\underline{i}-\underline{j}+4\underline{k}$ and $a\underline{i}+b\underline{j}-2\underline{k}$ are parallel.

\item If $\underline{u}=2\underline{i}+3\underline{j}+\underline{k}$, $\underline{v}=4\underline{i}+6\underline{j}+2\underline{k}$ and $\underline{w}=-6\underline{i}-9\underline{j}-3\underline{k}$, then show that $\underline{u}, \underline{v}$ and $\underline{w}$ are parallel to each other.

\item Suppose a butterfly's flight path passed through points $(2,4,7)$ and $(6,1,3)$, where each unit represents a metre. What is the magnitude of the displacement the butterfly experienced in traveling between these two points?

\item Calculate the projection of $\underline{a}$ along $\underline{b}$ and the projection of $\underline{b}$ along $\underline{a}$ when: $\underline{a}=2\underline{i}+3\underline{j}-\underline{k}$, $\underline{b}=\underline{i}-2\underline{j}+4\underline{k}$

\item Calculate the projection of $\underline{a}$ along $\underline{b}$ and the projection of $\underline{b}$ along $\underline{a}$ when: $\underline{a}=4\underline{i}-2\underline{j}+3\underline{k}$, $\underline{b}=\underline{i}+\underline{j}+\underline{k}$

\item Find the number $z$ so that the triangle with vertices $A(3,0,-2)$, $B(0,3,1)$ and $C(1,1,z)$ is a right triangle with the right angle at $C$.

\item If $\hat{a}$ and $\hat{b}$ are unit vectors and $2\theta$ is the angle between them, show that $\sin\theta = \dfrac{1}{2}|\hat{a}-\hat{b}|$.

\item Show that the set of points $P(4,-1,2)$, $Q(1,3,-1)$ and $R(-2,4,6)$ form a right triangle.

\item Prove that $\cos(\alpha+\beta)=\cos\alpha\cos\beta - \sin\alpha\sin\beta$.

\item Prove that in any triangle $ABC$: $b=c\cos A + a\cos C$.

\item Prove that in any triangle $ABC$: $c=a\cos B + b\cos A$.

\item Show that for any vectors $\underline{a}$ and $\underline{b}$: $\big||\underline{a}|-|\underline{b}|\big| \leq |\underline{a}+\underline{b}| \leq |\underline{a}|+|\underline{b}|$

\item Find the work done, if the point at which the constant force $\underline{F}=2\underline{i}+5\underline{j}+3\underline{k}$ is applied to an object moves it from $P_1(2,-3,0)$ to $P_2(7,5,3)$.

\item A particle, acted on by constant forces $\underline{F}_1=3\underline{i}+4\underline{j}-3\underline{k}$ and $\underline{F}_2=\underline{i}+4\underline{j}-\underline{k}$, is displaced from $A(2,1,3)$ to $B(5,4,4)$. Find the work done.

\item A particle is displaced from the point $A(5,-5,-7)$ to the point $B(6,2,-2)$ under the action of constant forces defined by $10\underline{i}-\underline{j}+11\underline{k}$, $4\underline{i}+5\underline{j}+9\underline{k}$ and $-2\underline{i}+\underline{j}-9\underline{k}$. Show that the total work done by the force is 102 units.

\item A force of magnitude 6 units acting parallel to $4\underline{i}+3\underline{j}-\underline{k}$ displaces the point of application from $A(2,-1,3)$ to $B(7,3,2)$. Find the work done.

\item Show that the components of a vector are the projections of that vector along $\underline{i}$, $\underline{j}$ and $\underline{k}$ respectively.

\item Prove that in any triangle $ABC$: $a^2=b^2+c^2-2bc\cos A$.

\item Prove that in any triangle $ABC$: $a^2=b^2+c^2-2bc\cos A$.

\item Prove that in any triangle $ABC$: $a=b\cos C + c\cos B$.

\item Prove that: $\cos(\alpha-\beta)=\cos\alpha\cos\beta + \sin\alpha\sin\beta$.

\item Find a scalar $a$ so that the vectors $2\underline{i}+a\underline{j}+5\underline{k}$ and $3\underline{i}+\underline{j}+a\underline{k}$ are orthogonal.

\item Show that the vectors $\overrightarrow{AB} =2\underline{i}-\underline{j}+\underline{k}$, $\overrightarrow{BC} = \underline{i}-3\underline{j}-5\underline{k}$ and $\overrightarrow{AC} =3\underline{i}-4\underline{j}-5\underline{k}$ are the sides of a right triangle.

\item The constant forces $2\underline{i}+5\underline{j}+6\underline{k}$ and $-\underline{i}-2\underline{j}-\underline{k}$ act on a body displaced from the position $P(4,-3,-2)$ to $Q(6,1,-3)$. Find the total work done.

\item Compute the cross products $\underline{a}\times\underline{b}$ and $\underline{b}\times\underline{a}$. Check your answer by showing that each of $\underline{a}$ and $\underline{b}$ is perpendicular to $\underline{a}\times\underline{b}$ and $\underline{b}\times\underline{a}$. $\underline{a}=2\underline{i}+\underline{j}-\underline{k}$, $\underline{b}=\underline{i}-\underline{j}+\underline{k}$

\item Compute the cross products $\underline{a}\times\underline{b}$ and $\underline{b}\times\underline{a}$. Check your answer by showing that each of $\underline{a}$ and $\underline{b}$ is perpendicular to $\underline{a}\times\underline{b}$ and $\underline{b}\times\underline{a}$. $\underline{a}=\underline{i}+3\underline{j}+2\underline{k}$, $\underline{b}=2\underline{i}-\underline{j}+\underline{k}$

\item Compute the cross products $\underline{a}\times\underline{b}$ and $\underline{b}\times\underline{a}$. Check your answer by showing that each of $\underline{a}$ and $\underline{b}$ is perpendicular to $\underline{a}\times\underline{b}$ and $\underline{b}\times\underline{a}$. $\underline{a}=2\underline{i}-2\underline{j}+\underline{k}$, $\underline{b}=-\underline{i}+\underline{j}+3\underline{k}$

\item Compute the cross products $\underline{a}\times\underline{b}$ and $\underline{b}\times\underline{a}$. Check your answer by showing that each of $\underline{a}$ and $\underline{b}$ is perpendicular to $\underline{a}\times\underline{b}$ and $\underline{b}\times\underline{a}$. $\underline{a}=-4\underline{i}+\underline{j}-2\underline{k}$, $\underline{b}=2\underline{i}+\underline{j}+\underline{k}$

\item Find a unit vector perpendicular to the plane containing $\underline{a}$ and $\underline{b}$. Also find the sine of the angle between them. $\underline{a}=\underline{i}+6\underline{j}-3\underline{k}$, $\underline{b}=2\underline{i}+\underline{j}+3\underline{k}$

\item Find a unit vector perpendicular to the plane containing $\underline{a}$ and $\underline{b}$. Also find the sine of the angle between them. $\underline{a}=-\underline{i}-\underline{j}-\underline{k}$, $\underline{b}=2\underline{i}-3\underline{j}+4\underline{k}$

\item Find a unit vector perpendicular to the plane containing $\underline{a}$ and $\underline{b}$. Also find the sine of the angle between them. $\underline{a}=\underline{i}+\underline{j}+\underline{k}$, $\underline{b}=\underline{i}-\underline{j}-\underline{k}$

\item Find a unit vector perpendicular to the plane containing $\underline{a}$ and $\underline{b}$. Also find the sine of the angle between them. $\underline{a}=5\underline{i}+\underline{j}-3\underline{k}$, $\underline{b}=-2\underline{i}+4\underline{j}+\underline{k}$

\item Find the area of the parallelogram, whose vertices are: $A(1,1,1)$; $B(4,2,3)$; $C(5,6,7)$; $D(2,5,5)$

\item Find the area of the parallelogram, whose vertices are: $A(4,5,6)$; $B(1,3,2)$; $C(-2,0,1)$; $D(1,2,5)$

\item Use the definition of the cross product, for any vectors $\underline{u},\underline{v},\underline{w}$ and scalar $k$, prove that: $\underline{u}\times\underline{v}= -\underline{v}\times\underline{u}$

\item Use the definition of the cross product, for any vectors $\underline{u},\underline{v},\underline{w}$ and scalar $k$, prove that: $\underline{u}\times(k\underline{v})=(k\underline{u})\times\underline{v}=k(\underline{u}\times\underline{v})$

\item Use the definition of the cross product, for any vectors $\underline{u},\underline{v},\underline{w}$ and scalar $k$, prove that: $\underline{u}\times(\underline{v}+\underline{w})=(\underline{u}\times\underline{v}) + (\underline{u}\times\underline{w})$

\item Prove that: $\sin(\alpha-\beta)=\sin\alpha\cos\beta - \cos\alpha\sin\beta$.

\item Use the definition of the cross product, prove that for any vectors $\underline{u}$ and $\underline{v}$: $(\underline{u}+\underline{v})\times(\underline{u}-\underline{v})=-2(\underline{u}\times\underline{v})$

\item Find the moment about the point $M(1,-3,3)$ of the force represented by $\overrightarrow{AB}$, where the coordinates of points $A(4,3,-1)$ and $B(-1,3,7)$ are given.

\item A force $\vec{F} = 6\underline{i}+4\underline{j}-4\underline{k}$ is applied at the point $A(1,-1,2)$. Find the moment of the force about the point $B(3,-2,3)$.

\item Given a force $\underline{F}=2\underline{i}+\underline{j}-3\underline{k}$ acting at a point $A(1,-2,1)$, find the moment of $\vec{F}$ about the point $B(2,0,2)$.

\item Find a vector perpendicular to each of the vectors. Also verify that $\underline{a}$ and $\underline{b}$ are (perpendicular) to $\underline{a}\times\underline{b}$. $\underline{a}=2\underline{i}-\underline{j}+\underline{k}$ and $\underline{b}=4\underline{i}+2\underline{j}-\underline{k}$

\item If $\underline{a}=4\underline{i}+3\underline{j}+\underline{k}$ and $\underline{b}=2\underline{i}-\underline{j}+2\underline{k}$, find a unit vector perpendicular to both $\underline{a}$ and $\underline{b}$. Also find the sine of the angle between the vectors $\underline{a}$ and $\underline{b}$.

\item Prove that $\sin(\alpha+\beta)=\sin\alpha\cos\beta + \cos\alpha\sin\beta$.

\item In any triangle $ABC$, prove that $\dfrac{a}{\sin A} = \dfrac{b}{\sin B} = \dfrac{c}{\sin C}$.

\item Find the area of the parallelogram whose vertices are: $P(0,0,0)$, $Q(-1,2,4)$, $R(2,-1,4)$ and $S(1,1,8)$.

\item Find the area of the triangle with vertices $A(1,-1,1)$, $B(2,1,-1)$ and $C(-1,1,2)$. Also find a unit vector perpendicular to the plane of triangle $ABC$.

\item Find the moment about the point $M(-2,4,6)$ of the force represented by $\overrightarrow{AB}$, where coordinates of points $A$ and $B$ are $(1,2,-3)$ and $(3,-4,2)$ respectively.

\item Find the volume of the parallelepiped whose given vectors are three edges: $\underline{u}=-3\underline{i}+2\underline{k}$; $\underline{v}=\underline{i}+2\underline{j}+\underline{k}$; $\underline{w}=\underline{j}+4\underline{k}$

\item Find the volume of the parallelepiped whose given vectors are three edges: $\underline{u}=\underline{i}-4\underline{j}-\underline{k}$; $\underline{v}=\underline{i}-\underline{j}-2\underline{k}$; $\underline{w}=2\underline{i}-3\underline{j}+\underline{k}$

\item Find the volume of the parallelepiped whose given vectors are three edges: $\underline{u}=\underline{i}-2\underline{j}+3\underline{k}$; $\underline{v}=2\underline{i}-\underline{j}-\underline{k}$; $\underline{w}=\underline{j}+\underline{k}$

\item Verify that $\underline{a}\cdot\underline{b}\times\underline{c}=\underline{b}\cdot\underline{c}\times\underline{a}=\underline{c}\cdot\underline{a}\times\underline{b}$ if $\underline{a}=3\underline{i}-\underline{j}+5\underline{k}$; $\underline{b}=4\underline{i}+3\underline{j}-2\underline{k}$ and $\underline{c}=2\underline{i}+5\underline{j}+\underline{k}$.

\item Find the constant $a$ such that the vectors are coplanar: $\underline{i}-2a\underline{j}-\underline{k}$, $\underline{i}-2\underline{j}+2\underline{k}$ and $a\underline{i}-2\underline{j}+\underline{k}$.

\item Prove that the points whose position vectors are $A(-6\underline{i}+3\underline{j}+2\underline{k})$, $B(3\underline{i}-2\underline{j}+4\underline{k})$, $C(5\underline{i}+7\underline{j}+3\underline{k})$, $D(-13\underline{i}+17\underline{j}-\underline{k})$ are coplanar.

\item Prove that the points whose position vectors are $A(3\underline{i}+2\underline{j}-\underline{k})$, $B(\underline{i}-2\underline{j}+\underline{k})$, $C(6\underline{i}+4\underline{j}-2\underline{k})$, $D(9\underline{i}+6\underline{j}-3\underline{k})$ are coplanar.

\item Prove that for any three non-zero vectors $\underline{u},\underline{v}$ and $\underline{w}$: $(\underline{u}+\underline{v})\cdot[(\underline{v}+\underline{w})\times(\underline{w}+\underline{u})]=2[\underline{u}\ \underline{v}\ \underline{w}]$

\item Consider a parallelepiped determined by the vectors $\underline{u}=2\underline{i}+4\underline{j}-3\underline{k}$, $\underline{v}=5\underline{i}-3\underline{j}+6\underline{k}$ and $\underline{w}=4\underline{i}-7\underline{j}-2\underline{k}$. If the base of the parallelepiped is defined by the vectors $\underline{u}$ and $\underline{v}$, then find the height of the parallelepiped.

\item A machine applies a force of 50 pounds along the positive $x$-axis on a wrench connected to a bolt. The pivot point of the wrench is at the origin $(0,0,0)$ and the force is applied at the point $(0\text{ ft}, 2\text{ ft}, 3\text{ ft})$. Determine the torque produced by this force about the pivot point.

\item A drone flies from point $(1,2,5)$ to point $(4,6,9)$ with each unit representing a metre. What is the magnitude of the displacement the drone experienced during this flight?

\item The vector $\underline{u}=50\underline{i}+75\underline{j}+65\underline{k}$ shows how many belts, pants and shirts were sold at a store. The vector $\underline{w}=1500\underline{i}+3500\underline{j}+3000\underline{k}$ shows the price (in rupees) of each item. Find $\underline{u}\cdot\underline{w}$ and explain what the result tells us in real life.

\item A force $\underline{F}=(20,-10,30)$N is applied at a point $P(2,-1,4)$ in 3D space. The pivot point is at $M(1,2,-3)$. Calculate the torque produced by this force about the pivot point $M$.

\item An electric shop sells three types of appliances: Fans, Heaters and Ovens. The monthly sales quantities are 500 units of Fans, 300 units of Heaters and 200 units of Ovens. The profit per unit for each appliance is Rs.\ 500 for Fans, Rs.\ 400 for Heaters and Rs.\ 2{,}000 for Ovens. Represent the monthly sales quantities and the profit per unit as vectors.

\item An electric shop sells three types of appliances: Fans, Heaters and Ovens. The monthly sales quantities are 500 units of Fans, 300 units of Heaters and 200 units of Ovens. The profit per unit for each appliance is Rs.\ 500 for Fans, Rs.\ 400 for Heaters and Rs.\ 2{,}000 for Ovens. Calculate the total monthly profit using vector operations.

\item Assuming $\underline{i}$, $\underline{j}$ and $\underline{k}$ are unit vectors in a Cartesian coordinate system, prove that $\underline{i}\cdot\underline{j}\times\underline{k}= \underline{j}\cdot\underline{k}\times\underline{i}= \underline{k}\cdot\underline{i}\times\underline{j}$.

\item A plumber exerts a force of 30 pounds along the negative $y$-axis on a lever connected to a machine. The pivot point of the lever is at the origin $(0,0,0)$, and the force is applied at the point $(1.2\text{ ft}, 0.5\text{ ft}, 0\text{ ft})$. Determine the torque produced by this force about the pivot point.

\item During a building construction, a crane exerts a force to pull a concrete block, represented by $\underline{F}= [4500,3300,2140]$ Newton. Each component corresponds to the force exerted along the $x$, $y$ and $z$ axes respectively. What is the magnitude of this force?

\item The components of $\underline{u}=300\underline{i}+250\underline{j}+180\underline{k}$ represent the respective number of jackets, shoes and handbags sold at a store. The components of $\underline{v}=3500\underline{i} + 4200\underline{j} + 6840\underline{k}$ represent the respective prices (in rupees) per unit for each product. Find $\underline{u}\cdot\underline{v}$ and explain what the result tells us in real life.

\end{longlist}

\end{document}